\documentclass[a4paper, 11pt, twoside, openright]{report}
%git-latexdiff 17807bf54f868 --main DP_parallelAddition.tex --cleanup none --latexmk


\usepackage{geometry}
\usepackage{cite}
\usepackage[english]{babel}
\usepackage[utf8]{inputenc}
% \usepackage[IL2]{fontenc}  
% \usepackage{amsmath,amsthm} 
% \usepackage{amsfonts}
\usepackage{amsmath, amsthm, amssymb, units, dsfont}

\usepackage[fixlanguage]{babelbib}
\selectbiblanguage{english}

\usepackage[pdfauthor={Jan Legersk\'y},
            pdfproducer={Jan Legersk\'y},
            pdfcreator={pdflatex},
            pdfencoding=unicode]{hyperref}
%\usepackage{bookmark}
\hypersetup{pdftitle={Construction of algorithms for parallel addition   in non-standard numeration systems}}
\hypersetup{pdfsubject={Master Thesis, Jan Legersk\'y}}
\hypersetup{pdfkeywords={parallel addition, non-standard numeration systems}}            

\usepackage[pdftex]{graphicx,color}


% fix left/right margins
\let\tmp\oddsidemargin
\let\oddsidemargin\evensidemargin
\let\evensidemargin\tmp
\reversemarginpar

\usepackage{fancyhdr}
	\pagestyle{fancy}
	\fancyhf{}
	\fancyhead[LO]{\slshape \rightmark}
	\fancyhead[RO,LE]{\thepage}
	\fancyhead[RE]{\slshape \leftmark}
	
\usepackage{emptypage}
%\setlength{\headheight}{16pt} 


\usepackage{amsmath, amsthm, amssymb, units, dsfont}
\usepackage{sidecap}
\usepackage{enumerate}
\usepackage{xcolor}

\usepackage{mathtools, mathdots}
\usepackage{pgffor}
\usepackage{pdflscape}
\usepackage{afterpage}
\usepackage{chngcntr}
\usepackage{multirow}
\usepackage{tabulary}

\usepackage{listings}
\usepackage{color}

\usepackage{algorithm}
\usepackage{algorithmic}

\usepackage{breqn}
\usepackage{hyperref}

\usepackage{pgfkeys}

\usepackage{changepage}

\definecolor{darkgreen}{rgb}{0,0.6,0}
\definecolor{darkred}{rgb}{0.6,0,0}
\definecolor{darkblue}{rgb}{0,0,0.6}
\definecolor{darkgrey}{rgb}{0.3,0.3,0.3}
\definecolor{grey}{rgb}{0.6,0.6,0.6} %comment
\definecolor{lightgrey}{rgb}{0.92,0.92,0.92}
\definecolor{terminal}{rgb}{0.9,0.9,0.6}
\definecolor{cmd}{rgb}{0.8,0.8,0.98}
\lstset{ 
        basicstyle=\ttfamily,
%         language=Matlab,                                % choose the language of the code
%       basicstyle=10pt,                                % the size of the fonts that are used for the code
%         numbers=left,                                   % where to put the line-numbers
        keywordstyle=\color{darkblue},
        commentstyle=\color{darkgreen},
        stringstyle=\color{darkred},
%         numberstyle=\footnotesize,                      % the size of the fonts that are used for the line-numbers
        stepnumber=1,                                           % the step between two line-numbers. If it's 1 each line will be numbered
        numbersep=5pt,                                  % how far the line-numbers are from the code
%       backgroundcolor=\color{white},          % choose the background color. You must add \usepackage{color}
        showspaces=false,                               % show spaces adding particular underscores
        showstringspaces=false,                         % underline spaces within strings
        showtabs=false,                                         % show tabs within strings adding particular underscores
%       frame=single,                                           % adds a frame around the code
%       tabsize=2,                                              % sets default tabsize to 2 spaces
%       captionpos=b,                                           % sets the caption-position to bottom
        breaklines=true,                                        % sets automatic line breaking
        breakatwhitespace=false,                        % sets if automatic breaks should only happen at whitespace
        escapeinside={\%*}{*)},                          % if you want to add a comment within your code
        emph={%  
    False, True%
    },emphstyle={\color{darkblue}}
}




%\newcommand{\komentar}[1]{\textcolor{red}{\MakeUppercase{#1}} \newline}
%\newenvironment{upravit}{\color{blue}}{}

\newcommand{\Zomega}{\mathbb{Z}[\omega]}
\newcommand{\Zbeta}{\mathbb{Z}[\beta]}

\newcommand{\ZZ}{\mathbb{Z}}
\newcommand{\QQ}{\mathbb{Q}}
\newcommand{\CC}{\mathbb{C}}
\newcommand{\NN}{\mathbb{N}}
\newcommand{\RR}{\mathbb{R}}


\newcommand{\A}{\mathcal{A}}
\newcommand{\B}{\mathcal{B}}
\newcommand{\Q}{\mathcal{Q}}

\newcommand{\Qw}[3][w]{\Q_{[#1_{-#2}, \dots, #1_{-#3}]}}
\newcommand{\Qwo}[2][w]{\Q_{[#1_{0}, \dots, #1_{-#2}]}}

\newcommand{\tuple}[3][w]{(#1_{-#2}, \dots, #1_{-#3})}
\newcommand{\tupleo}[2][w]{(#1_{0}, \dots, #1_{-#2})}

%\newcommand{\Qb}[1]{\mathcal{Q}_{[b^{#1}]}}
\newcommand{\Qb}[1]{\mathcal{Q}_{[\scriptstyle b]}^{\scriptstyle #1}}

\newcommand{\fin}[1]{\text{Fin}_{#1}(\beta)}

\newcommand{\multMat}[1]{\sum_{i=0}^{d-1} {#1}_i S^i}



\newcommand{\vect}[1]{\begin{pmatrix}
             {#1}_0 \\
             {#1}_1 \\
             \vdots \\
             {#1}_{d-1} 
             \end{pmatrix}}
             
\newcommand{\enum}[1]{({#1}_0,\ldots,{#1}_{d-1})}             

\newcommand{\vertiii}[1]{{\left\vert\kern-0.25ex\left\vert\kern-0.25ex\left\vert #1\right\vert\kern-0.25ex\right\vert\kern-0.25ex\right\vert}}
    
\newcommand{\norm}[2]{\left\lVert#1\right\rVert_{#2}}
\newcommand{\Mnorm}[2]{\vertiii{#1}_{#2}}
\newcommand{\normBeta}[1]{\norm{#1}{\beta}}
\newcommand{\MnormBeta}[1]{\Mnorm{#1}{\beta}}

\renewcommand\Re{\operatorname{Re}}
\renewcommand\Im{\operatorname{Im}}


\renewcommand{\algorithmicrequire}{\textbf{Input:}}
\renewcommand{\algorithmicensure}{\textbf{Ouput:}}
\algsetup{indent=2em}

 \usepackage{pifont}
 \renewcommand\checkmark{\ding{51}}
 \newcommand\xmark{\ding{55}}

 \newcommand{\var}[1]{\textit{#1}}
 \newcommand{\fun}[2]{\textbf{#1}\allowbreak{}(\var{#2})}

 \def\changemargin#1#2{\list{}{\rightmargin#2\leftmargin#1}\item[]}
 \let\endchangemargin=\endlist 


 \newenvironment{method}[2]{
 \noindent \textbf{#1(}\textit{#2}\textbf{)}
 \vspace{-5pt}
 \begin{changemargin}{3em}{0em}}
 {\end{changemargin}}

\def\Cpp{{C\nolinebreak[4]\hspace{-.05em}\raisebox{.4ex}{\tiny\bf ++}}}


 \pgfkeys{
  /phaseOnecaptions array/.is family, /phaseOnecaptions array,
  .unknown/.style = {\pgfkeyscurrentname/.initial = #1},
 }
 
 \newcommand\figurehascaptionOne[1]{\pgfkeys{/phaseOnecaptions array, #1}}
 \newcommand\getcaptionOne[1]{\pgfkeysvalueof{/phaseOnecaptions array/#1}}
 
 \pgfkeys{
  /phase2captions array/.is family, /phase2captions array,
  .unknown/.style = {\pgfkeyscurrentname/.initial = #1},
 }
 
 \newcommand\figurehascaptionTwo[1]{\pgfkeys{/phase2captions array, #1}}
 \newcommand\getcaptionTwo[1]{\pgfkeysvalueof{/phase2captions array/#1}}


% \hyphenation{coef-fi-cient}
% \hyphenation{Algorithm-For-Parallel-Addition}
% \hyphenation{Polynomial-Quotient-Ring}

\hyphenation{con-ver-gen-ce}
\hyphenation{non-con-ver-gen-ce}






\newtheorem{thm}{Theorem}[chapter]
\newtheorem{lem}[thm]{Lemma}

\theoremstyle{definition}
\newtheorem{defn}{Definition}[chapter]
\newtheorem{exmp}{Example}[chapter]

\renewcommand{\bibname}{References}

\setcounter{tocdepth}{1}




%%% k vyplnění ZAČÁTEK -------------------------------------------------
\newcommand{\cvut}{ČESKÉ VYSOKÉ UČENÍ TECHNICKÉ V~PRAZE}
\newcommand{\fjfi}{Fakulta jaderná a fyzikálně inženýrská}
\newcommand{\km}{Katedra matematiky}
\newcommand{\obor}{Matematická informatika}
%\newcommand{\zamereni}{Matematická informatika}

\newcommand{\nazevczBreaked}{Konstrukce algoritm\r u pro paraleln\'i s\v c\'it\'an\'i \\[8pt] v nestandardn\'ich \v c\'iseln\'ych soustav\'ach}
\newcommand{\nazevenBreaked}{Construction of algorithms for parallel addition  \\[8pt] in non-standard numeration systems}
\newcommand{\nazevcz}{Konstrukce algoritm\r u pro paraleln\'i s\v c\'it\'an\'i  v nestandardn\'ich \v c\'iseln\'ych soustav\'ach}
\newcommand{\nazeven}{Construction of algorithms for parallel addition   in non-standard numeration systems}
\newcommand{\autor}{Jan Legersk\'y}
\newcommand{\autorBc}{Bc. Jan Legersk\'y}
\newcommand{\vedouci}{Ing. \v St\v ep\'an Starosta, Ph.D.}
\newcommand{\pracovisteVed}{KAM FIT, \v CVUT v~Praze}
\newcommand{\konzultant}{---}

\newcommand{\klicova}{Paraleln\'i s\v c\'it\'an\'i, nestandardn\'i numera\v cn\'i syst\'emy, extending window method.}  % max. 5 klíčových slov
\newcommand{\keyword}{Parallel addition, non-standard numeration systems, extending window method.}
\newcommand{\abstrCZ}%{\begin{minipage}{0.7\textwidth}
{%Pr\'{a}ce je v\v{e}nov\'{a}na konstrukci algoritm\r{u} pro paraleln\'{\i} s\v{c}\'{\i}t\'{a}n\'{\i} v r\r{u}zn\'{y}ch numera\v{c}n\'{\i}ch soustav\'{a}ch. Zam\v{e}\v{r}uje se hlavn\v{e} na nestandardn\'{\i} numera\v{c}n\'{\i} syst\'{e}my s necelo\v{c}\'{\i}selnou, obecn\v{e} komplexn\'{\i}, b\'{a}z\'{\i} $\beta\in\Zomega$ a obecn\v{e} necelo\v{c}\'{\i}selnou abecedou $\A\subset\Zomega$ pro n\v{e}jak\'{e} algebraick\'{e} cel\'{e} \v{c}\'{\i}slo $\omega$. Navrhujeme takzvanou \emph{extending window method} s pou\v{z}it\'{\i}m z\'{a}kladn\'{\i}ho p\v{r}episovac\'{\i}ho pravidla $x-\beta$, kter\'{a} pro danou b\'{a}zi $\beta$ a abecedu $\A$ hled\'{a} algoritmus pro paraleln\'{\i} s\v{c}\'{\i}t\'{a}n\'{\i}. Metoda se skl\'{a}d\'{a} ze dvou f\'{a}z\'{\i}. Pro prvn\'{\i} z nich uv\'{a}d\'{\i}me posta\v{c}uj\'{\i}c\'{\i} podm\'{\i}nku konvergence, pro druhou m\'{a}me algoritmus, kter\'{y} kontroluje nutnou podm\'{\i}nku konvergence. Tuto metodu implementujeme v programovac\'{\i}m jazyce SageMath a uv\'{a}d\'{\i}me mno\v{z}stv\'{\i} otestovan\'{y}ch numera\v{c}n\'{\i}ch syst\'{e}m\r{u}.
} 
% \end{minipage}}
\newcommand{\abstrEN}{
%We focus on the construction of algorithms for parallel addition in different numeration systems, especially non-standard ones. The base $\beta$ is an element of $\Zomega$ and the alphabet, in general non-integer, is a subset of $\Zomega$ for some algebraic integer $\omega$. We design so-called extending window method with the basic rewriting rule $x-\beta$. The method searches for a parallel addition algorithm for a given base $\beta$ and alphabet $\A$. It consists of two phases. We have the sufficient condition of convergence of Phase 1. We introduce the algorithm which verifies the necessary condition of convergence of Phase 2. The method is implemented in SageMath and we provide several tested examples.
}
%%% k vyplnění KONEC ---------------------------------------------------


\begin{document}

\pagenumbering{roman}
\begin{titlepage}
%
%% 1. úvodní strana
% \thispagestyle{empty}
% \begin{center}
% 	{\Large \cvut \\[12pt] \fjfi\\}   % 260pt od nadpisu k "Bak. práce"
% 	\vspace{45pt}
% 	\includegraphics[height=100pt]{img/logoCVUT.pdf}\\  %bylo  75
% 	\vspace{160pt}
% 	{\Huge DIPLOMOV\'A PR\'ACE}
% \end{center}
% \vfill
% {
% 	\Large 2016 \hfill \autor
% }
% \newpage


%%% 2. úvodní strana
\thispagestyle{empty}
\begin{center}
	{\Large \cvut \\[10pt] \fjfi \\[10pt] \km\\}
	\vspace{45pt} %120
	\includegraphics[height=100pt]{img/logoCVUT.pdf}\\
	\vspace{90pt}
	{\Large DIPLOMOV\'A PR\'ACE}
	\vspace{90pt}
	
	{\Large\bf \nazevczBreaked}
	\vspace{30pt}
	
	{\Large\bf \nazevenBreaked}
\end{center}
\vfill
{
	\Large
	\begin{tabular}{ll}
	Vypracoval: & \autorBc\\[3pt]
	\v Skolitel: & \vedouci\\[3pt]
	Akademick\'y rok: & 2015/2016
	\end{tabular}
}
\cleardoublepage
\thispagestyle{empty}

%% sem přijde ofic. zadání
 \thispagestyle{empty}
 \Large
 Na toto místo přijde svázat \textbf{zadání diplomové práce}!
 
 \vspace{4mm}
 V~jednom z~výtisků musí být \textbf{originál} zadání, v~ostatních kopie.
 \normalsize
\cleardoublepage

%%% prohlášení
\thispagestyle{empty}
~
\vfill
\noindent\textbf{Čestné prohlášení}
\vspace{0.5cm}

Prohla\v suji na tomto m\'ist\v e, \v ze jsem p\v redlo\v zenou pr\'aci vypracoval samostatn\v e a \v ze jsem uvedl ve\v skerou pou\v zitou literaturu.
\vspace{1.5cm}

\noindent
\vspace{5mm}V Praze dne 4. května 2016\hfill
	\begin{tabular}{c}
	\\[40pt]
	\autor
	\end{tabular}
\cleardoublepage


%%% poděkování
\thispagestyle{empty}
~
\vfill
\noindent\textbf{Pod\v ekov\'an\'i}
\vspace{0.5cm}

Děkuji Ing. Štěpánu Starostovi, Ph.D., za vedení mé diplomové práce. Dále děkuji Ing. Mileně Svobodové, Ph.D., a prof. Ing. Editě Pelantové, CSc., za  podnětné diskuze o implementované metodě. Velký dík patří také mým rodičům za podporu během celého studia.

\begin{flushright}
\autor
\end{flushright}
\cleardoublepage


%%% abstrakt atp.
\thispagestyle{empty}

\noindent
\begin{tabular}{lp{0.78\textwidth}}
  {\em N\'azev pr\'ace:} & \bf \nazevcz \\[4mm]
  {\em Autor:} & \autor \\[4mm]
  {\em Obor:} & \obor \\[4mm]
%  {\em Zam\v e\v ren\'i:} & \zamereni \\[4mm]
  {\em Druh pr\'ace:} & Diplomová práce \\[4mm]
  {\em Vedoucí práce:} & \vedouci, \pracovisteVed \\[4mm]
  {\em Konzultant:} & \konzultant \\[4mm]
  {\em Abstrakt:} & \abstrCZ \\[4mm]
  {\em Kl\'i\v cov\'a slova:} & \klicova \\[20mm]

  {\em Title:} & \bf \nazeven \\[4mm]
  {\em Author:} & \autor \\[4mm]
  {\em Abstract:} & \abstrEN \\[4mm]
  {\em Key words:} & \keyword
\end{tabular}
\cleardoublepage


%%% obsah
\thispagestyle{empty}
\renewcommand\contentsname{\vspace{-3cm} Contents \vspace{-0.8cm}}
\tableofcontents
\thispagestyle{empty}

\cleardoublepage
\thispagestyle{empty}
\chapter*{List of symbols}
% \centering
\begin{tabular}{ll}
Symbol        & Description \\ \hline
$\NN$         & set of nonnegative integers $\{0,1,2,3,\dots\}$   \\
$\ZZ$         & set of integers $\{\dots,-2,-1,0,1,2,\dots\}$ \\
$\RR$           & set of real numbers \\
$\CC$           & set of complex numbers \\
$\QQ$           &set of rational numbers \\
$\QQ(\beta)$    &the smallest field containing  $\QQ$ and algebraic number  $\beta$ \\
$\#S$          & number of elements of the finite set $S$ \\
$C^*$            & complex conjugation and transposition of the complex matrix $C$ \\
\rule{0cm}{0cm}& \\
$m_\beta$       &monic minimal polynomial of the algebraic number $\beta$ \\
$\deg \beta$    &degree of the algebraic number $\beta$ (over $\QQ$)\\
\rule{0cm}{0cm}& \\
$(\beta,\A)$            & numeration system with the base $\beta$ and the alphabet $\A$\\
$(x)_{\beta,\A}$    &$(\beta,\A)$-representation of the number $x$\\
$\fin{\A}$          &set of all complex numbers with a finite $(\beta,\A)$-representation \\
$\A^\ZZ$        &set of all bi-infinite sequences of digits in $\A$\\
$\Zomega$       &the smallest ring containing $\ZZ$ and $\omega$ \\%&set of values of all polynomials with integer coefficients evaluated in $\omega$\\
$\pi$           &isomorphism from $\Zomega$ to $\ZZ^d$ ($d=\deg \omega$)\\
\rule{0cm}{0cm}& \\
$\B$            &alphabet of input digits\\
$q_j$           &weight coefficient for the $j$-th position \\
$\Q$            &weight coefficients set\\
$\Qwo{k}$ 		&set of possible weight coefficients for digits $\tupleo{k}\in\B^{k+1}$ \\
%\rule{0cm}{0cm}& \\
%$\lfloor x \rfloor$ & floor function of the number $x$ \\  
%$\Re x$           & real part of the complex number $x$ \\
%$\Im x$           & imaginary part of the complex number $x$
\end{tabular}
\thispagestyle{empty}
\cleardoublepage
\end{titlepage}
\pagenumbering{arabic}


%%%%%%%%%%%%%%%%%%%%%%%%%%%%%%%%%%%%%%%%%%%%%%%%%%%%%%%%%%%%%%%%%%%%%%%%
%     Z A Č Á T E K   P R Á C E                                        %
%%%%%%%%%%%%%%%%%%%%%%%%%%%%%%%%%%%%%%%%%%%%%%%%%%%%%%%%%%%%%%%%%%%%%%%%



\cleardoublepage\phantomsection
\addcontentsline{toc}{chapter}{Introduction}
\chapter*{Introduction}
\markboth{\MakeUppercase{Introduction}}{}

Parallel addition  is an important part of algorithms for multiplication and division. 
While carry propagation in standard algorithms requires to compute digits of the sum of $x_n x_{n-1}\cdots x_1 x_0 \bullet$ and $y_n y_{n-1}\cdots y_1 y_0 \bullet$  one by one, a parallel algorithm determines the $j$-th digit of the sum just from a fixed number of digits around $x_j$ and $y_j$. Thus it avoids the carry propagation  and all digits digits of the result can be outputted at the same time. Besides theoretical reasons, parallel addition is used for instance in floating point division algorithm in Pentium processors \cite{division}. The used base is 4 and the digit set is $\{-2,-1,0,1,2\}$.

A parallel addition algorithm  was introduced by A. Avizienis in 1961 \cite{avizienis}. The algorithm works with an integer base $\beta\geq3$ and an alphabet $\{-a, \dots, 0, \dots a\}$ where $a\in\NN$ is such that $\beta/2 <a \leq \beta-1$. Later, C. Y. Chow and J. E. Robertson presented a parallel addition algorithm for  base 2 and alphabet $\{-1,0,1\}$ \cite{chow}.

So-called non-standard numeration systems are extensively studied. Non-standard means that a base is not a positive integer. An example of such system is the Penney numeration system with the base $\imath-1$ and the alphabet is $\{0,1\}$. Complex numbers can be represented by this system without separating real and imaginary part. As division algorithms are also developed for non-standard numeration systems, we focus on construction of a parallel addition algorithm for these systems.

The numeration systems which allow parallel addition must be redundant, i.e, some numbers have more than one representation. 
C. Frougny, E. Pelantov\'a and M. Svobodov\'a provided parallel addition algorithms with an integer alphabet for all bases $\beta$ such that $|\beta|>1$ and no conjugate of $\beta$ equals 1 in modulus \cite{parAddNS}. Nevertheless, the alphabet is not minimal in general.
 The parallel addition algorithms for several bases (negative integers, complex numbers $\imath-1, 2\imath$ and $\sqrt{2}\imath$, quadratic Pisot unit and  non-integer rationals) with minimal integer alphabet are presented \cite{minAlph}.
 
Let $\omega$ be an algebraic integer. We are interested in construction of parallel addition algorithm for a given base $\beta\in\Zomega$ which is an algebraic integer such that $|\beta|>1$ and generally non-integer alphabet $\A\subset\Zomega$. Our aim is to implement a method which attempts to construct an algorithm automatically for the numeration system $(\beta,\A)$. The implementation is used to find as many numeration systems allowing parallel addition as possible. 

The structure of the thesis is the following: we review terminology and known results on numeration systems and parallel addition in Chapter~\ref{chap:preliminaries}. In Chapter~\ref{chap:ewm}, a general concept of addition is recalled and so-called extending window method for construction of parallel addition algorithm is outlined. This method, which consists of two phases, was proposed by M. Svobodov\'a \cite{milena}. 

Some useful tools related to the ring $\Zomega$ from which  are developed in Chapter~\ref{chap:Zomega}, including a norm derived from a companion matrix. Using these results, a necessary condition of convergence the extending window method is proved in Chapter~\ref{chap:convergence}. We show that it is also a sufficient condition of convergence of Phase~1. We establish check of non-convergence of Phase~2 in Theorem~\ref{thm:infinitePathInRauzyGraph}. Moreover, a lower bound on the size of an alphabet which allows parallel addition is stated. We also indicate a way how an alphabet with necessary properties can be generated.

Different variants of both phases of the extending window method are proposed in Chapter~\ref{chap:diffChoices}. In Chapter~\ref{chap:design}, we present algorithms based on the theorems from Chapter~\ref{chap:convergence} and we describe their implementation in SageMath.

Finally, several examples of results can be found in Chapter~\ref{chap:testing}. We compare various methods from Chapter~\ref{chap:diffChoices} and discuss successful numeration systems with a quadratic base and integer or non-integer alphabet. 

Appendix contains images illustrating the extending window method, interfaces of program and details of tested examples.

\chapter{Preliminaries}
\label{chap:preliminaries}
This chapter summarizes definitions related to numeration systems and parallel addition. We recall some known results on properties of a numeration system which allows parallel addition. At the end of the chapter, we define the ring $\Zomega$, where $\omega$ is an algebraic integer, and determine numeration systems within the scope of this thesis.
	In this chapter, we recall few definitions and results connected to numeration systems and parallelism. We define the set $\Zomega$ for an algebraic integer $\omega$ and we prove that $\Zomega$ is isomorphic to $\ZZ^d$. This property is used in Theorem \ref{thm:divisibility} which is an important tool for divisibility in $\Zomega$. Division in $\Zomega$ is necessary for the extending window method described in Chapter \ref{chap:methodDescription}.

\section{Numeration systems}
Firstly, we give a general definition of numeration system.
\begin{defn}
  Let $\beta \in \CC, |\beta|>1$ and $\A \subset \CC$ be a finite set containing 0. A pair $(\beta, \A)$ is called a \emph{positional numeration system} with \emph{base} $\beta$ and \emph{digit set} $\A$, usually called \emph{alphabet}.
\end{defn}
So-called standard numeration systems have an integer base $\beta$ and an alphabet $\A$ which is a set of contiguous integers. We restrict ourselves to the base $\beta$ which is an algebraic integer and possibly non-integer alphabet $\A$. 
\begin{defn}
Let $(\beta, \A)$ be a positional numeration system.  We say that a complex number $x$ has a \emph{$(\beta, \A)$-representation} if~ there exist digits $x_n,x_{n-1}, x_{n-2},\dots \in\A, n\geq 0$ such that $x=\sum_{j=-\infty}^n x_j \beta^j$.
\end{defn}
 We write briefly a \emph{representation} instead of a $(\beta, \A)$-representation if the base $\beta$ and the alphabet $\A$ follow from context. 

\begin{defn}
Let $(\beta, \A)$ be a positional numeration system. The set of all complex numbers with a finite $(\beta, \A)$-representation is defined by
$$
    \fin{\A}:=\left\{\sum_{j=-m}^n x_j \beta^j\colon n, m \in \NN, x_j \in \A \right\}\,.
$$
\end{defn}
   
For  $x\in\fin{\A}$, we write 
$$
(x)_{\beta,\A}= 0^\omega x_n x_{n-1}\cdots x_1 x_0 \bullet x_{-1} x_{-2} \cdots x_{-m} 0^\omega\,,
$$ 
where $0^\omega$ denotes right, respectively left-infinite sequence of zeros. Notice that indices are decreasing from left to right as it is usual to write the most significant digits first. In what follows, we omit the starting and ending $0^\omega$ when we work with numbers in $\fin{\A}$. We remark that existence of an algorithm (standard or parallel) producing a finite $(\beta,\A)$-representation of $x+y$ where $x,y\in\fin{\A}$ implies that the set $\fin{\A}$ is closed under addition, i.e.,
$$
\fin{\A} + \fin{\A} \subset \fin{\A}\,.
$$ 

Designing an algorithm for parallel addition requires some redundancy in numeration system. According to \cite{redundant}, a numeration system $(\beta,\A)$ is called \emph{redundant} if there exists $x \in \fin{\A}$ which has two different $(\beta,\A)$-representations. For instance, the number 1 has $(2,\{-1,0,1\})$-representations $1\bullet$ and $1(-1)\bullet$.
Redundant numeration system can enable us to avoid carry propagation in addition. On the other hand, there are some disadvantages. For example, comparison is problematic.  


\section{Parallel addition}
A local function, which is also often called a sliding block code, is used to mathematically formalize parallelism. 
\begin{defn}
Let $\A$ and $\B$ be alphabets. A function $\varphi:\B^\ZZ \rightarrow \A^\ZZ$ is said to be \emph{$p$-local} if there exist $r,t\in\NN$ satisfying $p=r+t+1$ and a function $\phi: \B^p \rightarrow \A$ such that, for any $w=(w_j)_{j\in\ZZ}\in\B^\ZZ$ and its image $z=\varphi(w)=(z_j)_{j\in\ZZ}\in\A^\ZZ$, we have $z_j=\phi(w_{j+t},\cdots,w_{j-r})$ for every $j\in\ZZ$. The parameter $t$, resp. $r$, is called \emph{anticipation}, resp. \emph{memory}.
\end{defn}
This means that each digit of the image $\varphi(w)$ is computed from $p$ digits of $w$ in a sliding window. Suppose that there is a processor on  each position with access to $t$ input digits on the left and $r$ input digits on the right. Then computation of $\varphi(w)$, where $w$ is a finite sequence, can be done in constant time independent on the length of $w$.   
  
\begin{defn}
\label{def:digitSetConversion}
Let $\beta$ be a base and $\A$ and $\B$ two alphabets containing 0. A function $\varphi:\B^\ZZ\rightarrow \A^\ZZ$ such that
  \begin{enumerate}
      \item for any $w=(w_j)_{j\in\ZZ}\in\B^\ZZ$ with finitely many non-zero digits, $z=\varphi(w)=(z_j)_{j\in\ZZ}\in\A^\ZZ$ has only finite number of non-zero digits, and
      \item $\sum_{j\in\ZZ} w_j \beta^j= \sum_{j\in\ZZ} z_j \beta^j$
  \end{enumerate}
  is called \emph{digit set conversion} in base $\beta$ from $\B$ to $\A$. Such a conversion $\varphi$ is said to be \emph{computable in parallel} if $\varphi$ is a $p$-local function for some $p\in\NN$. 
\end{defn}
In fact, addition on $\fin{\A}$ can be performed in parallel if there is a digit set conversion from $\A+\A$ to $\A$ computable in  parallel as we can easily output digitwise sum of two $(\beta,\A)$-representations in parallel.   


We recall few results about parallel addition in a numeration system with an integer alphabet. C. Frougny, E. Pelantov\'a and M. Svobodov\'a proved  the following sufficient condition of existence of an algorithm for parallel addition in \cite{parAddNS}.
  \begin{theo}
  \label{thm:suffConjugates}
  Let $\beta\in\CC$ be an algebraic number such that $|\beta|>1$ and all its conjugates in modulus differ from 1. There exists an alphabet $\A$ of contiguous integers containing 0 such that addition on $\fin{\A}$ can be performed in parallel.
  \end{theo}
  The proof of the theorem provides the algorithm for the alphabet of the form $\{-a,-a+1, \dots,0,\dots,a-1,a\}$. But in general, $a$ is not minimal.
    
The same authors showed in \cite{kBlock} that the condition on the conjugates of the base $\beta$ is also necessary:
  \begin{theo}
  Let the base $\beta\in\CC, |\beta|>1,$ be an algebraic number with a conjugate $\beta'$ such that $|\beta'|=1$. Let $\A\subset\ZZ$ be an alphabet of contiguous integers containing 0. Then addition on $\fin{\A}$ cannot be computable in parallel.
  \end{theo}
  
The question of minimality of the alphabet is studied in \cite{minAlph}. The following lower bound for the size of the alphabet is provided:
  \begin{theo}
  \label{thm:lowerBoundAlphabet}
  Let $\beta\in\CC, |\beta|>1,$  be an algebraic integer with the minimal polynomial $p$. Let $\A\subset\ZZ$ be an alphabet of contiguous integers containing 0 and 1. If addition on $\fin{\A}$ is computable in parallel, then $\#\A \geq |p(1)|$. Moreover, if $\beta$ is a positive real number, $\beta>1$, then $\#\A \geq  |p(1)|+2$.
  \end{theo}
  

In this thesis, we work in a more general concept as we consider also non-integer alphabets. First, we recall the following definition.
\begin{defn}
Let $\omega$ be a complex number. The set of values of all polynomials with integer coefficients evaluated in $\omega$ is denoted by
$$
    \ZZ[\omega] =\left\{\sum_{i=0}^n a_i \omega^i\colon n\in\NN, a_i\in\ZZ \right\} \subset \QQ(\omega)\,.
$$
\end{defn}
 Notice that $\ZZ[\omega]$ is a commutative ring (for our purposes, a ring is associative under multiplication and there is a multiplicative identity).     
    
From now on, let $\omega$ be an algebraic integer  which generates the set $\Zomega$ and let the base $\beta\in\Zomega$ be such that $|\beta|>1$. We remark that $\beta$ is also an algebraic integer as all elements of $\Zomega$ are algebraic integers. Finally, let the alphabet $\A$ be a finite subset of $\Zomega$ such that $0\in\A$.

Few parallel addition algorithms for such numeration system with a non-integer alphabet were found ad hoc. We introduce the method for construction of the parallel addition algorithm for a given numeration system $(\beta,\A)$ in Chapter \ref{chap:methodDescription}. 
  


\section{\texorpdfstring{Isomorphism of $\Zomega$ and $\ZZ^{d}$}{Isomorphism of Z[omega] and Zd}}
The goal of this section is to show a connection between the ring $\Zomega$ and the set $\ZZ^d$. Using Theorem \ref{thm:divisibility}, division in $\Zomega$ can be replaced by searching for an integer solution of a linear system. This is used for the implementation of the extending window method.

First we recall the notion of companion matrix which we use to define multiplication in $\ZZ^d$. By the minimal polynomial of an algebraic integer, we always mean the monic minimal polynomial.  
\begin{defn}
Let $\omega$ be an algebraic integer of degree $d\geq 1$ with the  minimal polynomial $p(x)=x^d +p_{d-1}x^{d-1}+ \cdots + p_1 x+p_0 \in \ZZ[x]$. The matrix 
$$
S := \begin{pmatrix}
            0 & 0 & \cdots & 0 & -p_0 \\
            1 & 0 & \cdots & 0 & -p_1 \\
            0 & 1 & \cdots & 0 & -p_2 \\
            \vdots &   & \ddots & & \vdots \\
            0 & 0 & \cdots & 1 & -p_{d-1} 
            \end{pmatrix} \in \ZZ^{d\times d}
$$
is called \emph{companion matrix} of the minimal polynomial of $\omega$.
\end{defn}
In what follows, the standard basis vectors of $\ZZ^d$  are denoted by 
$$
e_0=\begin{pmatrix}
              1 \\
              0 \\
              0 \\
              \vdots \\
              0
              \end{pmatrix}, \\
e_1=\begin{pmatrix}
              0 \\
              1 \\
              0 \\
              \vdots \\
              0
              \end{pmatrix}, \dots ,\\
e_{d-1}=\begin{pmatrix}
              0 \\        
              \vdots \\
              0 \\
              0\\
              1
              \end{pmatrix}\,.             
$$
% We remark that 1 in $e_i$ is in the $(i+1)$-st row because the index corresponds to the power of a companion matrix in the following definition. 

\begin{defn}
Let $\omega$ be an algebraic integer of degree $d\geq 1$, let $p$ be its minimal polynomial and let $S$ be its companion matrix. We define the mapping $\odot_\omega: \ZZ^d \times \ZZ^d \rightarrow \ZZ^d$ by 
$$
u \odot_\omega v := \left(\multMat{u}\right)\cdot \vect{v} \quad \text{ for all } u=\vect{u}, v=\vect{v} \in \ZZ^d\,.
$$ 
and we define powers of $u \in \ZZ^d$ by
\begin{align*}
    u^0&=e_0, \\
    u^{i}&= u^{i-1} \odot_\omega u \text{ for } i\in\NN\,.
\end{align*}
\end{defn}

We will see later that $\ZZ^d$ equipped with elementwise addition and multiplication $\odot_\omega$ builds a commutative ring. 
% It will follow from the isomorphism with $\Zomega$. 
Let us first recall an important property of a companion matrix  -- it is a root of its defining polynomial.
\begin{lem}
\label{lem:compMatrixIsRoot}
Let $\omega$ be an algebraic integer with a minimal polynomial $p$ and let $S$ be its companion matrix. Then
$$
p(S)=0\,.
$$
\end{lem}






Now we can prove that there is a correspondence between elements of $\Zomega$ and $\ZZ^d$.

\begin{theo}
Let  $\omega$ be an algebraic integer of degree $d$. Then 
$$
\Zomega =\left\{\sum_{i=0}^{d-1} a_i \omega^i \colon a_i\in\ZZ \right\},
$$ 
$(\ZZ^d,+,\odot_\omega)$ is a commutative ring and the mapping $\pi:\Zomega \rightarrow \ZZ^{d}$ defined by 
$$
\pi(u)=\vect{u} \quad \text{ for every } u=\sum_{i=0}^{d-1} u_i \omega^i \in \Zomega
$$
is a ring isomorphism.
\end{theo}


Due to this theorem we may work with integer vectors instead of elements of $\Zomega$ and multiplication in $\Zomega$ is replaced by multiplying by an appropriate matrix. 

The last theorem of this section is a practical tool for divisibility in $\Zomega$. To check whether an element of $\Zomega$ is divisible by another element, we look for an integer solution of a linear system. Moreover, this solution provides the result of  division in the positive case. 
\begin{theo}
\label{thm:divisibility}
Let $\omega$ be an algebraic integer of degree $d$ and let $S$ be the companion matrix of its minimal polynomial. Let $\beta=\sum_{i=0}^{d-1} b_i \omega^i$ be a nonzero element of $\Zomega$. Then for every $u\in\Zomega$
$$
u\in\beta\Zomega \iff S_\beta^{-1}\cdot \pi(u) \in \ZZ^d\,,
$$
where $S_\beta=\multMat{b}$.
\end{theo}
























  
   




\chapter{Design of extending window method}
\label{chap:ewm}
In this chapter, the general concept of addition and digit set conversion is recalled. We outline a so-called \emph{extending window method} which is due to M. Svobodov\'a \cite{milena}. The method consists of two phases. For a given numeration system $(\beta,\A)$, it attempts to construct a digit set conversion algorithm which is computable in parallel. We recall that $\omega$ is an algebraic integer, $\beta \in \Zomega$ is a base and $0\in\A\subset\Zomega$ is an alphabet. 


	
We recall the general concept of addition at the beginning of this chapter and then we describe a so-called \emph{extending window method} which is due to M. Svobodov\'a \cite{milena}. The method attempts to construct a conversion algorithm computable in parallel for a given numeration system $(\beta,\A)$. We recall that $\omega$ is an algebraic integer, $\beta \in \Zomega$ is a base and $0\in\A\subset\Zomega$ is an alphabet. 

\section{Addition}

The general idea of addition (standard or parallel) in any numeration system $(\beta,\A)$ is the following: we sum up two numbers digit-wise and then we convert the result with digits in $\A+\A$ into the alphabet $\A$. Obviously, digit-wise addition is computable in parallel, thus the problematic part is the digit set conversion of the obtained result. It can be easily done in a standard way but a parallel digit set conversion is non-trivial. A parallel conversion is based on the same  formulas as the standard one but the choice of so-called \emph{weight coefficients} differs.

Now, we go step by step more precisely. Let $x,y \in \fin{\A}$ with $(\beta,\A)$-representations $x_{n'}x_{{n'}-1}\cdots x_1 x_0\bullet x_{-1} x_{-2} \cdots x_{-m'}$ and $y_{n'}y_{{n'}-1}\cdots y_1 y_0\bullet y_{-1} y_{-2} \cdots y_{-m'}$ padded by zeros to have the same length $n'+m'+1$. We set 
  \begin{align*}
    w&=x+y =\sum_{i=-m'}^{n'} x_i\beta^i + \sum_{i=-m'}^{n'} y_i\beta^i = \sum_{i=-m'}^{n'} (x_i+y_i)\beta^i \\
    &=\sum_{i=-m'}^{n'} w_i\beta^i \,,
  \end{align*}
  where $w_i=x_i+y_i \in \A +\A$. Thus, $w_{n'} w_{{n'}-1}\cdots w_1 w_0 \bullet w_{-1} w_{-2} \cdots w_{-m'}$ is a  $(\beta, \A+\A)$-representation of $w\in \fin{\A+\A}$. 

We also use column notation for digit-wise addition in what follows, e.g.,    
%	\begin{center}
%	\begin{tabulary}{0.7\textwidth}{CCCCCCCCCCCL}
%	$x_{n'}$& $x_{{n'}-1}$ & $\cdots$ & $x_1$ &$x_0$ &$\bullet$  &$x_{-1}$ & $x_{-2}$ & $\cdots$ & $x_{-m'}$ & $=$& $(x)_{\beta,\A}$\\
%	$y_{n'}$& $y_{{n'}-1}$ & $\cdots$ & $y_1$ &$y_0$ &$\bullet$  &$y_{-1}$ & $y_{-2}$ & $\cdots$ & $y_{-m'}$ & $=$& $(y)_{\beta,\A}$\\ \hline
%	$w_{n'}$& $w_{{n'}-1}$ & $\cdots$ & $w_1$ &$w_0$ &$\bullet$  &$w_{-1}$ & $w_{-2}$ & $\cdots$ & $w_{-m'}$ & $=$& $(x+y)_{\beta,\A+\A}$\\\,.
%	\end{tabulary}	
%	\end{center}
 
  \begin{align*}
  x_{n'} \;x_{{n'}-1}\cdots x_1 \;x_0 &\bullet x_{-1} \;x_{-2}\, \cdots x_{-m'} \\[-3pt]
  y_{n'} \;y_{{n'}-1}\cdots y_1 \,\;y_0 &\bullet y_{-1} \;y_{-2} \;\cdots y_{-m'} \\[-10pt]
    \line(1,0){90} & \line(1,0){100} \\[-7pt]
  w_{n'} w_{{n'}-1}\cdots w_1 w_0 &\bullet w_{-1} w_{-2} \cdots w_{-m'}\,.
  \end{align*}
  
We search for a $(\beta,\A)$-representation of $w$, i.e., a  sequence 
  $$z_{n} z_{n-1}\cdots z_1 z_0 z_{-1} z_{-2} \cdots z_{-m}$$ such that $z_j \in \A$ and
  $$
    z_{n} z_{n-1}\cdots z_1 z_0 \bullet z_{-1} z_{-2} \cdots z_{-m}=(w)_{\beta,\A}\,.
  $$
  Note that the index of the first, resp. last, non-zero digit of the converted representation $z_{n} z_{n-1}\cdots z_1 z_0 \bullet z_{-1} z_{-2} \cdots z_{-m}=(w)_{\beta,\A}$ may differ from the original representation $w_{n'} w_{{n'}-1}\cdots w_1 w_0 \bullet w_{-1} w_{-2} \cdots w_{-m'}$. We assume that $ n\geq n'$ and $m\geq m'$, otherwise we pad the converted representation by zeros.
  
   Multiplication of a representation $w_{n'} w_{n-1}\cdots w_1 w_0 \bullet w_{-1} w_{-2} \cdots w_{-m'}$ by a power of $\beta$ is obvious:
  $$
  \beta^m \cdot w_{n'} w_{n'-1}\cdots w_1 w_0 \bullet w_{-1} w_{-2} \cdots w_{-m'} = w_{n} w_{n'-1}\cdots w_1 w_0 w_{-1} w_{-2} \cdots w_{-m'} \bullet
  $$  
  and after conversion
  $$
  z_{n} z_{n-1}\cdots z_1 z_0 z_{-1} z_{-2} \cdots z_{-m'}\bullet\cdots z_{-m} = \beta^m \cdot z_{n} z_{n-1}\cdots z_1 z_0 \bullet z_{-1} z_{-2} \cdots z_{-m'}\,. 
  $$  
Hence, without lost of generality, we consider only conversion of so-called $\beta$-integers -- numbers from $\fin{\A+\A}$ whose representations have all digits with negative indices equal to zero.
%  Particularly, let $(w)_{\beta, \A+\A}=w_{n'} w_{{n'}-1}\cdots w_1 w_0 \bullet$. We search for a number $n \in \NN$ and \mbox{$z_{n}, z_{n-1},\dots, z_1, z_0 \in\A$} such that $(w)_{\beta, \A}=z_{n} z_{n-1}\cdots z_1 z_0 \bullet$.   
  
  Digits $w_j$ are converted into the alphabet $\A$ by adding a suitable representation of zero digit-wise.
  For our purpose, we use the simplest possible representation which is deduced from the polynomial
  $$
    x-\beta \in \left(\Zomega\right)[x]\,.
  $$

We remark that any polynomial $R(x)=r_s x^s+r_{s-1}x^{s-1}+ \dots + r_1 x+r_0$ with coefficients $r_i \in \Zomega$ such that $R(\beta)=0$ gives us a possible representation of zero. The polynomial $R$ is called a \emph{rewriting rule}. One of the coefficients of $R$ which is greatest in modulus (so-called \emph{core coefficient}) is used for the conversion of a digit $w_j$. Nevertheless, the extending window method is strongly dependent on the rewriting rule, so we focus only on the simplest possible rewriting rule $R(x)=x-\beta$. Usage of an arbitrary rewriting rule $R$ is out of scope of this thesis.

 
As $0=\beta^{j} \cdot R(\beta)=1\cdot \beta^{j+1} -\beta \cdot \beta^{j}$, we have a representation of zero 
$$1 (-\!\beta) \underbrace{0 \cdots 0}_{j}\bullet = (0)_\beta\,. $$
for all $j \in \NN$. We multiply this representation by $q_j \in \Zomega$, which is called a \emph{weight coefficient}, to obtain another  representation of zero 
%$$q_j (-q_j\beta) \underbrace{0 \cdots 0}_{j}\bullet = (0)_\beta\,. $$ 
$q_j (-q_j\beta) 0 \cdots 0\bullet = (0)_\beta\,. $
This is digit-wise added to $w_{n} w_{n-1}\cdots w_1 w_0 \bullet$ to convert the digit $w_j$ into the alphabet $\A$. The conversion of $j$-th digit causes a \emph{carry} $q_{j}$ on the $(j+1)$-th position. 

In standard addition, the digit set conversion runs from the right ($j=0$) to the left until all non-zero digits and carries are converted into the alphabet $\A$:

	\begin{tabular}{rcccccccclcl}
	$w_{n'}$ & $w_{{n'}-1}$ & $\cdots$ & $w_{j+1}$ &\textcolor{red}{$w_{j}$} &$w_{j-1}$ & $\cdots$ & $w_1$ &$w_0$ &$\bullet$ & $=$&$(w)_{\beta,{\A+\A}}$\\
	  &   & &   &  &$q_{j-2}$ & $\iddots$ &  &  &\\
	   &   & &   & \textcolor{red}{$q_{j-1}$}  & $-\beta q_{j-1}$ &  &  &  &\\
	  &   &  $\iddots$&  $q_{j}$ &  \textcolor{red}{$-\beta q_{j}$} & & &  &  &\\ \cline{1-10}
	$z_{n} \cdots z_{n'}$ & $z_{{n'}-1}$ & $\cdots$ & $z_{j+1}$ &\textcolor{red}{$z_{j}$} &$z_{j-1}$ & $\cdots$ & $z_1$ &$z_0$ &$\bullet$  & $=$&$(w)_{\beta,{\A}}$\\
	\end{tabular}

%        \begin{align}
%        \label{eq:conversionScheme}
%            \hspace{100pt}  w_n w_{n-1}&&&\cdots& &w_{j+1}&\!\! &\textcolor{red}{w_j}  & \!\!  &w_{j-1} &&\cdots &&w_1 w_0\bullet \hspace{100pt} \notag\\[-5pt]
%                         &&&&       &       & &     &   &q_{j-2} &&\iddots  \notag \\[-3pt] 
%                         &&&&       &       & &\textcolor{red}{q_{j-1}}& -&\beta q_{j-1} \notag \\[-3pt]
%                         &&&&         &q_j&   \textcolor{red}{-}&\textcolor{red}{\beta q_j} &&\\[-8pt]
%                         &&&  \iddots      &   -&\beta q_{j+1}&   &\ && \notag \\[-17pt]
%          \intertext{\hspace{60pt}\line(1,0){300}}
%          \notag \\[-30pt]
%           z_{n'} \cdots z_{n} z_{n-1}&&&\cdots& &z_{j+1}& &\textcolor{red}{z_j}& &z_{j-1} &&\cdots &&z_1 \; z_0\bullet \notag                  
%        \end{align}
    Hence, the desired formula for conversion on the $j$-th position is 
    \begin{equation*}
        z_j=w_j + q_{j-1} - q_j \beta
    \end{equation*}
    for $j \in \NN$. We set $q_{-1}=0$ as there is no carry from the right on the 0-th position.
    
     The terms carry and weight coefficient are related to a position: while $q_{j-1}$ is a carry from the right  and $q_j$ is a chosen weight coefficient on the $j$-th position, $q_j$ is a carry from the right on the $(j+1)$-th position etc.

We remark that the conversion with the rewriting rule $x-\beta$ prolongs the part of non-zero digits only to the left as there is no carry from the left. Thus, all digits with negative indices of the converted sequence are zero.


     The fact that the conversion preserves the value of $w$ follows from adding a representation of zero:
\begin{align}
\label{eq:valuePreserving}
    \sum_{j\geq 0} z_j \beta^j &=w_0 - \beta q_0 + \sum_{j> 0} (w_j + q_{j-1} - q_j \beta) \beta^j \notag\\
    &=\sum_{j\geq 0} w_j \beta^j + \sum_{j>0} q_{j-1} \beta^j - \sum_{j\geq 0} q_j \cdot \beta^{j+1}  \\
    &=\sum_{j\geq 0} w_j \beta^j + \sum_{j>0} q_{j-1} \beta^j - \sum_{j> 0} q_{j-1} \cdot \beta^j \notag\\
    &=\sum_{j\geq 0} w_j \beta^j = w\,. \notag
\end{align}

    The weight coefficient $q_j$ must be chosen so that the converted digit is in the alphabet~$\A$, i.e., 
    \begin{equation}
    \label{eq:conversionFormula}
        z_j=w_j + q_{j-1} - q_j \beta \in \A\,.
    \end{equation} 
    The choice of weight coefficients is a crucial part of the construction of addition algorithms which are computable in parallel. The extending window method determining weight coefficients for a given input is described in Section~\ref{sec:methodDescription}.
    
    
     On the other hand, the following example shows that determining weight coefficients is trivial for  numeration systems such that an alphabet contains right one representative of each class modulo $\beta$.
     
     \begin{exmp}
        Assume now a numeration system $(\beta, \A)$, where
  $$
    \beta \in \NN\,,\beta  \geq 2\,, \A = \{0, 1, 2,\dots, \beta -1 \}\,.
  $$ 
       Notice that
    $$
        z_j \equiv w_j+q_{j-1} \mod \beta\,. 
    $$
  
  There is only one representative of each class modulo  $\beta$ in the standard numeration system $(\beta, \A)$. Therefore, the digit $z_j$ is uniquely determined for a given digit $w_j \in \A+\A$ and carry $q_{j-1}$ and thus so is the weight coefficient $q_j$. This means that $q_j=q_j(w_j,q_{j-1})$ for all $j\geq 0$. Generally,
  $$
  q_j=q_j(w_j,q_{j-1}(w_{j-1},q_{j-2}))=\dots =q_j(w_j ,\dots , w_1, w_0)
  $$
  and
  $$
  z_j=z_j(w_j ,\dots , w_1, w_0)\,,
  $$
  which implies that addition runs in linear time. For instance, the carry $q_{j-1}=1$ propagates through the whole result when we sum up $(\beta-1)(\beta-1)\dots(\beta-1)\bullet$ and $1\bullet$.
     
     \end{exmp}
  
  We require that the digit set conversion from $\A+\A$ into $\A$ is computable in parallel, i.e., there exist constants $r,t \in \NN_0$ such that for all $j\geq 0$ is $z_j=z_j(w_{j+t},\dots,w_{j-r})$. Anticipation $t$ equals zero since we use the rewriting rule $x-\beta$. To avoid the dependency on all less significant digits, we need variety in the choice of the weight coefficient $q_j$. This implies that the used numeration system must be redundant.
  

\section{Extending window method}
\label{sec:methodDescription}
In order to construct a digit set conversion in numeration system $(\beta,\A)$ from $\A+\A$ to $\A$ which is computable in parallel, we consider a more general case of digit set conversion from an \emph{input alphabet} $\B$ such that $\A \subsetneq \B$ instead of the alphabet $\A+\A$.
As mentioned above, the key problem is to find for every $j\geq 0$ a weight coefficient $q_j$ such that 
    $$
        z_j=\underbrace{w_j}_{\in \B} + q_{j-1} - q_j \beta \in \A 
    $$  
    for any input $w_{n'}w_{n'-1}\dots w_1 w_0 \bullet=(w)_{\beta,\B}, w\in \fin{\B}$. We remark that the weight coefficient $q_{j-1}$ is determined by the input $w_{j-1}\dots w_1 w_0 \bullet$. For a digit set conversion with the rewriting rule $x-\beta$ to be computable in parallel, the digit $z_j$ is required to satisfy $z_j=z_j(w_{j},\dots,w_{j-r})$ for a fixed memory $r$ in $\NN$.
    
    Note that the digit $z_j$ is given by the input digit $w_j$ and carry $q_{j-1}$ which is determined by input digits  $w_{j-1} w_{j-2}\dots$. Thus, if we find a weight coefficient $q_j$ for all possible combinations of input digits $w_j w_{j-1} w_{j-2}\dots$, then the position $j$ is not important. Therefore, we may strongly simplify our notation if we omit $j$ in subscripts. From now on, $w_0\in\B$ is a converted digit, $w_{-1} w_{-2}\dots\in\B$ are digits on right, $q_{-1}\in\Zomega$ is a carry from the right and we search for a weight coefficient $q_0\in\Zomega$ such that 
    $$
    z_0=w_0 + q_{-1} - q_0 \beta \in \A\,.
    $$
    
   
    We introduce two definitions before we describe the extending window method.
    \begin{defn}
    \label{def:weightCoefficientsSet}
        Let $\B$ be a set such that $\A \subsetneq \B$. Then any finite set $\Q\subset\Zomega$ containing~0 such that 
        $$
            \B + \Q \subset \A + \beta \Q
        $$  
        is called a \emph{weight coefficients set}.
    \end{defn}
    We see that if $\Q$ is a weight coefficients set, then
        $$
        (\forall w_0 \in \B)(\forall q_{-1}\in\Q)(\exists q_0 \in \Q )(\underbrace{w_0 + q_{-1} - q_0 \beta}_{z_0} \in \A )\,.
        $$
    In other words, there is a weight coefficient $q_0 \in \Q$ for a carry from the right $q_{-1}\in \Q$ and a digit $w_0$ in the input alphabet $\B$ such that $z_0$ is in the alphabet $\A$.  Notice that  the  carry from the right for the rightmost non-zero digit of the converted sequence which is $0$ is in $\Q$ by the definition.
    \begin{defn}
    Let $r$ be an integer and $q:\B^{r} \rightarrow \Q$ be a mapping such that 
    $$
    w_0+ q(w_{-1}, \dots, w_{-r}) - \beta q\tupleo{(r-1)} \in \A
    $$
    for all $w_0,w_{-1}, \dots, w_{-r} \in \B$, and $q(0,0,\dots,0)=0$. Then $q$ is called \emph{weight function} and $r$~is called \emph{length of window}.    
    \end{defn}

 Having a weight function $q$, we define a function $\phi:\B^{r+1}\rightarrow \A$ by
    \begin{equation}
    \label{eq:localConversion}
        \phi(w_{0}, \dots, w_{-r})=w_0+ \underbrace{q(w_{-1}, \dots, w_{-r})}_{=q_{-1}} - \beta \underbrace{q\tupleo{(r-1)}}_{=q_0}=:z_0\,,
    \end{equation} 
    which verifies that the digit set conversion is indeed a $(r+1)$-local function with anticipation $0$ and memory $r$. The requirement of zero output of the weight function $q$ for the input of $r$ zeros guarantees that $\phi(0,0,\dots,0)=0$. Thus, the first condition of Definition~\ref{def:digitSetConversion} is satisfied. The second one follows from the equation \eqref{eq:valuePreserving}. 
    
Let us summarize the construction of the digit set conversion by the rewriting rule \mbox{$x-\beta$}. We need to find weight coefficients for all possible combinations of digits of the input alphabet~$\B$. The rewriting rules multiplied by the weight coefficients are digit-wise added to an input sequence. In fact, it means that the equation  \eqref{eq:conversionFormula} is applied on each position. If the digit set conversion is computable in parallel, the weight coefficients are determined as the outputs of the weight function $q$ with some fixed length of window $r$.  

We search for a weight function $q$ for a given base $\beta$ and input alphabet $\B$ by the extending window method. It consists of two phases. First, we find some weight coefficients set $\Q$. We know that it is possible to convert an input sequence by choosing the weight coefficients from the set $\Q$. The set $\Q$ serves as the starting point for the second phase in which we increment the expected length of the window $r$ until the weight function $q$ is uniquely defined for each $\tupleo{(r-1)} \in \B^{r}$. Then, the local conversion is determined -- we use the weight function outputs as weight coefficients in the formula \eqref{eq:localConversion}.    

We describe the general concept of the extending window method in this chapter, while various possibilities of construction of sets during both phases are discussed in Chapter~\ref{chap:diffChoices}.
Note that  convergence of both phases is studied in Chapter~\ref{chap:convergence}.
      
\section{Phase 1 -- Weight coefficients set}
\label{subsec:phase1}
The goal of the first phase is to compute a weight coefficients set $\Q$, i.e., to find a set $\Q \ni 0$ such that 
$$
    \B + \Q \subset \A + \beta \Q\,.
$$  
We build a sequence $\Q_0, \Q_1, \Q_2,\dots$ iteratively so that we extend $\Q_k$ to $\Q_{k+1}$ in a way to cover all elements of the set $\B+\Q_k$ by elements of the extended set $\Q_{k+1}$, i.e.,
$$
\B+ \Q_k \subset \A + \beta \Q_{k+1}\,.
$$
This procedure is repeated until the extended weight coefficients set $\Q_{k+1}$ is the same as the previous set $\Q_{k}$. We remark that the expression ``a weight coefficient $q$ covers an element $x$'' means that there is a digit $a \in \A$ such that $x=a + \beta q$. 

In other words, we start with $\Q_0=\{0\}$ meaning that we search all weight coefficients $q_0$ necessary for digit set conversion for the case where there is no carry from the right, i.e., $q_{-1}=0$. We add them to the weight coefficients set $\Q_0$ to obtain the set $\Q_1$. Assume now that we have a set $\Q_k$ for some $k\geq 1$. The weight coefficients in $\Q_k$ now may appear as a carry $q_{-1}$. If there are no suitable weight coefficients $q_0$ in the weight coefficients set~$\Q_k$ to cover all sums of coefficients from $\Q_k$ and digits of the input alphabet $\B$, we extend $\Q_k$ to $\Q_{k+1}$ by  suitable coefficients. And so on until there is no need to add more elements, i.e., the extended set $\Q_{k+1}$ equals $\Q_k$. Then the weight coefficients set $\Q:=\Q_{k+1}$ satisfies Definition~\ref{def:weightCoefficientsSet}. 

\komentar{For better understanding, see Figures~\ref{img:phase1img1}--\ref{img:phase1img13} in Appendix~\ref{app:phase1} which illustrate the construction of the weight coefficients set $\Q$ for the Eisenstein base and a complex alphabet (see Example~\ref{ex:Eisenstein1-blockcomplex} for its description). }

The precise description of the algorithm in a pseudocode is in Algorithm~\ref{alg:weightCoefSet}. Observe that extending $\Q_k$ to $\Q_{k+1}$ is not unique. Various methods of choice are described in Section~\ref{sec:methodsOne} in Algorithm~\ref{alg:extendWeightCoefSet}.
    



Section~\ref{sec:convergencePhase2} discusses the convergence of Phase 1, i.e. whether it happens that  $\Q_{k+1}=\Q_k$ for some  $k$.
    
\begin{algorithm}
  \caption{Search for weight coefficients set (Phase 1)}
    \label{alg:weightCoefSet}
  \begin{algorithmic}[1]
    \STATE $k:=0$ 
    \STATE $Q_0:=\{0\}$
    \REPEAT
     \STATE Extend $\Q_k$ to $\Q_{k+1}$ (by Algorithm~\ref{alg:extendWeightCoefSet}) so that $$\B+ \Q_k \subset \A + \beta \Q_{k+1}$$
     \vspace{-20pt}
      \STATE  $k:=k+1$
      \UNTIL{$\Q_k = \Q_{k+1}$}      
      \STATE $\Q:=\Q_k$
    \RETURN $\Q$
  \end{algorithmic}
\end{algorithm}


% An added element from each list of \verb+candidates+ is chosen as the smallest one unless there is already a covering element contained in $\Q_{k}$ (Algorithm~\ref{alg:extendWeightCoefSet}).  


    

  
    
    





\section{Phase 2 -- Weight function}
\label{subsec:phase2}
    We want to find a length of the window $r$ and a weight function $q:\B^{r} \to \Q$. We start with the weight coefficients set $\Q$ obtained in Phase 1. The idea is to reduce necessary weight coefficients which are used to convert a given input digit up to a single value. This is done by enlarging the number of considered input digits, i.e. incrementing $r$.  If the window is extended to the right, we know more digits that cause a carry form the right. This may decrease the number of  possible carries from the right and hence, less weight coefficients to convert the input digit may be necessary.
     
    We introduce notation for sets of possible weight coefficients for given input digits.
        Let $\Q$ be a weight coefficients set and $w_0\in \B$. Denote by $\Q_{[w_0]}$ any set such that
        $$
            (\forall q_{-1} \in \Q)(\exists q_0 \in \Q_{[w_0]})(w_0 + q_{-1} - q_0 \beta \in \A)\,.
        $$
It means that we do not know any input digits on the right, therefore there might be any carry from the set $\Q$. However, we may determine a set $\Q_{[w_0]}$ of  weight coefficients which allow the conversion of $w_0$ to $\A$ since we know the input digit $w_0$.
        
        By induction with respect to $k \in \NN, k\geq 1$, for all $\tupleo{k}\in \B^{k+1}$ denote by $\Qwo{k}$ any subset of  $\Qwo{(k-1)}$ such that 
        $$
           (\forall q_{-1} \in \Qw{1}{k})(\exists q_0 \in \Qwo{k})(w_0 + q_{-1} - q_0 \beta \in \A)\,.
        $$
        
    
 
%    Recall the scheme \eqref{eq:conversionScheme} of the digit set conversion for better understanding of the notation and method:
%    \begin{align*}
%        \hspace{130pt}\cdots\; &w_{j+1}&\!\! &w_j  & \!\!  &w_{j-1}&\cdots w_{j-M+1} &w_{j-M}\cdots \hspace{130pt} \\[-3pt] 
%                         & &       & & & q_{j-2} \\[-3pt]
%                         & &       &q_{j-1}& -&\beta q_{j-1} \\[-3pt]
%                           &q_j&   -&\beta q_j &&\\[-3pt]      
%                           -&\beta q_{j+1}&   &  &&\\[-15pt]      
%    \intertext{\hspace{120pt}\line(1,0){250}} 
%          \vspace{-15pt}
%          \\[-30pt]
%     \cdots\; &z_{j+1}& &z_j& &z_{j-1}& \cdots z_{j-M+1}\; &z_{j-M}\cdots                     
%    \end{align*}     

Sets of possible weight coefficients and a weight function $q$ are constructed by Algorithm~\ref{alg:weightFunction}. The idea is to check all possible right carries $q_{-1}\in\Q$ and determine values $q_0\in\Q$ such that 
    $$
    z_0=w_0 + q_{-1} - q_0 \beta \in \A \,.
    $$  
    
    So we obtain a set $\Q_{[w_0]}\subset\Q$ of weight coefficients which are necessary to convert the digit~$w_0$ with any carry $q_{-1}\in\Q$. Assuming that we know the input digit $w_{-1}$, the set of possible carries from the right is also reduced to $\Q_{[w_{-1}]}$. Thus we may reduce the set $\Q_{[w_0]}$ to a set $\Qwo{1}\subset \Q_{[w_0]}$ which is necessary to cover all elements of $w_0 + \Q_{[w_{-1}]}$. 

In the $k$-th step, we search for a set $\Qwo{k}\subset\Qwo{(k-1)}$ such that 
               $$
              w_0 + \Qw{1}{k} \subset \A + \beta \Qwo{k}\,.
              $$
              The length of window is $k+1$, i.e., we know $k$ digits on the right. To  construct the set $\Qwo{k}$, we select from $\Qwo{(k-1)}$ such weight coefficients which are necessary to convert digit $w_0$ to the alphabet $\A$ with all possible carries from the set $\Qw{1}{k}$.
                 
    Proceeding in this manner may lead to a unique weight coefficient $q_0$ for enough long window.     
    If there is $r\in\NN, r\geq 1$ such that 
    $$
    \#\Qwo{(r-1)}=1
    $$
    for all $\tupleo{(r-1)} \in \B^r$, then the output $q\tupleo{(r-1)}$ is defined as the element of $\Qwo{(r-1)}$. 
    
    Similarly to Phase 1, the choice of $\Qwo{k}$ is not unique. We list different methods of choice in Section~\ref{sec:methodsTwo}, Algorithm~\ref{alg:minimalSet}.
    \komentar{ Figures~\ref{img:phase2img1}--\ref{img:phase2img7} in Appendix~\ref{app:phase2} illustrate the construction of the set $\Q_{[\omega,1,2]}$ for the Eisenstein numeration system.   }
    
        
    To verify that 
$$
	z_0=\phi(w_{0}, \dots, w_{-r})=w_0+ \underbrace{q\tuple{1}{r}}_{=q_{-1}} - \beta \underbrace{q\tupleo{(r-1)}}_{=q_0}
$$    
is in the alphabet $\A$, consider that $q_0=q\tupleo{(r-1)}$ is the only element of $\Qwo{(r-1)}$ which was constructed such that 
$$
w_0 + \Qw{1}{(r-1)} \subset \A +\beta \Qwo{(r-1)}\,.
$$
At the same time, $q_{-1}=q\tupleo{(r-1)}$ is the only element of $\Qw{1}{r}$ which is a subset of $\Qw{1}{(r-1)}$.

    

    
\begin{algorithm}
  \caption{Search for weight function $q$ (Phase 2)}
    \label{alg:weightFunction}
  \begin{algorithmic}[1]
    \REQUIRE{weight coefficients set $\Q$}
    \FORALL{$w_0 \in \B$} 
        \STATE Find set $\Q_{[w_0]} \subset \Q$ (by Algorithm~\ref{alg:minimalSet}) such that
          $$
          w_0 + \Q \subset \A + \beta \Q_{[w_0]}
          $$
    \ENDFOR
    \STATE $k:=0$
    \WHILE{$\max\{\#\Qwo{k}\colon \tupleo{k}\in \B^{k+1} \} > 1$}
        \STATE $k:= k +1$
        \FORALL{$\tupleo{k}\in \B^{k+1}$}
            \STATE Find set $\Qwo{k} \subset \Qwo{(k-1)}$ (by Algorithm~\ref{alg:minimalSet}) such that
              $$
              w_0 + \Qw{1}{k} \subset \A + \beta \Qwo{k}\,,
              $$
        \ENDFOR  
    \ENDWHILE  
    \STATE $r:= k+1$ 
    \FORALL{$\tupleo{(r-1)} \in \B^{r}$}  
        \STATE $q\tupleo{(r-1)}\in \B^{r}:=$ only element of $\Qwo{(r-1)}$
    \ENDFOR
    \RETURN $q$
  \end{algorithmic}
\end{algorithm}

Unfortunately, finiteness of Phase 2 is not guaranteed. But the non-convergence of Phase 2 with a specific property may be revealed by Algorithm~\ref{alg:oneletterSets} before the run of Phase 2 or by Algorithm~\ref{alg:checkCycles} during it. These algorithms are based on theorems in Chapter ~\ref{chap:convergence}. Modified Phase 2 which includes these algorithms can be found in Section~\ref{sec:modifiedPhase2}.



Notice that for a given length of window $r$, the number of calls of Algorithm~\ref{alg:minimalSet} within Algorithm~\ref{alg:weightFunction} is
$$
\sum_{k=0}^{r-1}  \#\B^{k+1} = \#\B \frac{\#\B^r-1}{\#\B-1}\,.
$$    
It implies that the time complexity grows exponentially. The required memory is also exponential because we have to store sets $\Qwo{k}$ at least for $k\in\{r-2, r-1\}$  for all $w_0,\dots, w_{-k} \in \B$.


\chapter{\texorpdfstring{Properties of $\Zomega$}{Properties of Z[omega]}}
\label{chap:Zomega}
We compile some properties  of the ring $\Zomega$ in this chapter. We first show that $\Zomega$ is isomorphic to $\ZZ^d$ and we use this fact to check divisibility in $\Zomega$.  We review some results from matrix theory to introduce a norm in $\Zomega$. This norm is used in the proof of convergence of Phase~1 in Chapter~\ref{chap:convergence}. In the last section of this chapter, we show how the number of congruence classes modulo $\beta$ in $\Zomega$ is determined. This result is used in the proof of a lower bound on the size of an alphabet which allows parallel addition in Section~\ref{sec:minimalAlphabet}.
	


\section{\texorpdfstring{Isomorphism of $\Zomega$ and $\ZZ^{d}$}{Isomorphism of Z[omega] and Zd}}
In this section, we recall that  the ring $\Zomega$ is isomorphic to the set $\ZZ^d$ equiped with multiplication, where $d$ is the degree of an algebraic integer $\omega$. This structure is useful as it allows to handle elements of $\Zomega$ as vectors. For example, division in $\Zomega$ can be replaced by searching for an integer solution of a linear system (Theorem \ref{thm:divisibility}) which is used in our implementation of the extending window method.


For defining multiplication in $\ZZ^d$, we recall the notion of companion matrix. 
\begin{defn}
Let  $p(x)=x^d +p_{d-1}x^{d-1}+ \cdots + p_1 x+p_0 \in \ZZ[x]$ be a monic polynomial with integer coefficients, $d\geq 1$. The matrix 
$$
S := \begin{pmatrix}
            0 & 0 & \cdots & 0 & -p_0 \\
            1 & 0 & \cdots & 0 & -p_1 \\
            0 & 1 & \cdots & 0 & -p_2 \\
            \vdots &   & \ddots & & \vdots \\
            0 & 0 & \cdots & 1 & -p_{d-1} 
            \end{pmatrix} \in \ZZ^{d\times d}
$$
is called \emph{companion matrix} of the polynomial $p$.
\end{defn}
Let $S$ be the companion matrix of a polynomial $p$. It is well known (see for instance \cite{horn}) that the characteristic polynomial of the companion matrix $S$ is $p$. The matrix $S$ is also root of the polynomial $p$.


%In what follows, the standard basis vectors of $\ZZ^d$  are denoted by 
%$$
%e_0=\begin{pmatrix}
%              1 \\
%              0 \\
%              0 \\
%              \vdots \\
%              0
%              \end{pmatrix}, \\
%e_1=\begin{pmatrix}
%              0 \\
%              1 \\
%              0 \\
%              \vdots \\
%              0
%              \end{pmatrix}, \dots ,\\
%e_{d-1}=\begin{pmatrix}
%              0 \\        
%              \vdots \\
%              0 \\
%              0\\
%              1
%              \end{pmatrix}\,.             
%$$
% We remark that 1 in $e_i$ is in the $(i+1)$-st row because the index corresponds to the power of a companion matrix in the following definition. 
We remark that the minimal polynomial of an algebraic integer $\omega$ is denoted by $m_\omega$ and is always meant to be monic.  
Multiplication in $\ZZ^d$ is defined in the following way.
\begin{defn}
Let $\omega$ be an algebraic integer of degree $d\geq 1$ and let $S$ be the companion matrix of $m_\omega$. We define a mapping $\odot_\omega: \ZZ^d \times \ZZ^d \rightarrow \ZZ^d$ by 
$$
u \odot_\omega v := \left(\multMat{u}\right)\cdot \vect{v} \quad \text{ for all } u=\vect{u}, v=\vect{v} \in \ZZ^d\,.
$$ 
%and we define powers of $u \in \ZZ^d$ by
%\begin{align*}
%    u^0&=e_0, \\
%    u^{i}&= u^{i-1} \odot_\omega u \text{ for } i\in\NN\,.
%\end{align*}
\end{defn}

The technical proof of Theorem~\ref{thm:isomorphismWithZd}, which shows that $\ZZ^d$ with multiplication $\odot_\omega$ is a ring isomorphic to $\Zomega$, can be found in \cite{vu}.

\begin{thm}
\label{thm:isomorphismWithZd}
Let  $\omega$ be an algebraic integer of degree $d$. Then 
$$
\Zomega =\left\{\sum_{i=0}^{d-1} a_i \omega^i \colon a_i\in\ZZ \right\},
$$ 
$(\ZZ^d,+,\odot_\omega)$ is a commutative ring and the mapping $\pi:\Zomega \rightarrow \ZZ^{d}$ defined by 
$$
\pi(u)=\vect{u} \quad \text{ for every } u=\sum_{i=0}^{d-1} u_i \omega^i \in \Zomega
$$
is a ring isomorphism.
\end{thm}

This theorem provides simple representation of elements of $\Zomega$ in computer -- they are represented by integer vectors and multiplication in $\Zomega$ is replaced by multiplying by an appropriate matrix. 

Divisibility in $\Zomega$ can be also transformed into $\ZZ^d$. To check whether an element of $\Zomega$ is divisible by another element, we look for an integer solution of a linear system given by Theorem~\ref{thm:divisibility}. Moreover, this solution provides the result of  division in the positive case. 
\begin{thm}
\label{thm:divisibility}
Let $\omega$ be an algebraic integer of degree $d$ and let $S$ be the companion matrix of its minimal polynomial. Let $\beta=\sum_{i=0}^{d-1} b_i \omega^i$ be a nonzero element of $\Zomega$. Then for every $u\in\Zomega$
$$
u\in\beta\Zomega \iff S_\beta^{-1}\cdot \pi(u) \in \ZZ^d\,,
$$
where $S_\beta=\multMat{b}$.
\end{thm}
The proof can be found in \cite{vu}.

	\section{$\beta$-norm}
The goal of this section is to construct a norm in $\Zomega$. We use the isomorphism with $\ZZ^d$ and some results from matrix theory. This norm is used for the proof of convergence of Phase~1 in Chapter~\ref{chap:convergence}.

First, we recall a simple way how a new norm can be constructed from another one.

\begin{lem}
\label{lem:vectNorm}
Let $\nu$ be a norm of the vector space $\CC^d$ and $P$ be a nonsingular matrix in $\CC^d$. Then the mapping $\mu:\CC^d\rightarrow \RR^+_0$ defined by $\mu(x)=\nu(Px)$ is also a norm of the vector space $\CC^d$.
\end{lem}
\begin{proof}
Let $x$ and $y$ be vectors in $\CC^d$ and $\alpha\in \CC$.  We use linearity of matrix multiplication, nonsingularity of matrix $P$ and the fact that $\nu$ is a norm to prove the following statements:
\begin{enumerate}
    \item $\mu(x)=\nu(Px)\geq 0\,,$
    \item $\mu(x)=0 \iff \nu(Px)=0 \iff Px=0 \iff x=0\,,$
    \item $\mu(\alpha x)=\nu(P(\alpha x))=\nu(\alpha Px)=|\alpha|\nu(Px)=|\alpha|\mu(x)\,,$
    \item $\mu(x+y)=\nu(P(x+y))=\nu(Px+Py)\leq \nu(Px)+\nu(Py)=\mu(x)+\mu(y)\,.$
\end{enumerate}
This  verifies that $\mu$ is a norm.
\end{proof}


Now we use Lemma \ref{lem:vectNorm} to define a new norm for a given diagonalizable matrix.
\begin{defn}
\label{def:newNorm}
Let $M\in\CC^{n\times n}$ be a diagonalizable matrix and $P\in\CC^{n\times n}$ be a nonsingular matrix which diagonalizes $M$, i.e., $M=P^{-1}DP$ for some diagonal matrix $D\in\CC^{n\times n}$. We define a vector norm $\norm{\cdot}{M}$ by  
\begin{equation}
\norm{x}{M}:=\norm{Px}{2}
\end{equation}
for all $x\in\CC^n$, where $\norm{\cdot}{2}$ is Euclidean norm. A matrix norm $\Mnorm{\cdot}{M}$ is induced by the norm $\norm{\cdot}{M}$, i.e, 
$$
\Mnorm{A}{M}=\sup_{\norm{x}{M}=1} \norm{Ax}{M}\,
$$
for all $A\in\CC^{n\times n}$.
\end{defn}

The following theorem is a known result from matrix theory -- for a given diagonalizable matrix, there is a matrix norm such that the norm of the matrix equals its spectral radius.
\komentar{je potreba nejakou citaci?}

\begin{thm}
\label{thm:norm}
Let $M\in\CC^{n\times n}$ be a diagonalizable matrix. Then %re exists a vector norm $\norm{\cdot}{M}$ such that 
$$
\rho(M)=\Mnorm{M}{M}\,,
$$
where $\rho(M)$ is the spectral radius of the matrix $M$. % and $\Mnorm{\cdot}{M}$ is the natural matrix norm induced by $\norm{\cdot}{M}$.
\end{thm}
\begin{proof}
First, we prove that $\Mnorm{M}{M}\geq\rho(M)$. For all eigenvalues $\lambda$ in the spectrum $\sigma(M)$ with a respective eigenvector $u$ such that $\norm{u}{M}=1$, we have
$$
\Mnorm{M}{M}=\sup_{\norm{x}{M}=1} \norm{Mx}{M}\geq \norm{Mu}{M}=\norm{\lambda u}{M}=|\lambda|\cdot\norm{u}{M}=|\lambda|\,.
$$
Secondly, we show that $\Mnorm{M}{M}\leq\rho(M)$. Following Definition~\ref{def:newNorm}, let $P\in\CC^{n\times n}$ be a  nonsingular matrix  and $D\in\CC^{n\times n}$ a diagonal matrix  with the eigenvalues of $M$ on the diagonal such that $PMP^{-1}=D$.

Let $y$ be a  vector such that $\norm{y}{M}=1$ and set $z=Py$. Notice that 
$$
\sqrt{z^*z}=\norm{z}{2}=\norm{Py}{2}=\norm{y}{M}=1\,.
$$
Consider
\begin{align*}
\norm{My}{M}&=\norm{PMy}{2}=\norm{DPy}{2}=\norm{Dz}{2}=\sqrt{z^*D^*Dz}\\
    &\leq \sqrt{\max_{\lambda\in\sigma(M)}|\lambda|^2 z^*z}=\max_{\lambda\in\sigma(M)}|\lambda|=\rho(M)\,\,.
\end{align*}
which implies the statement.
\end{proof}


Before we define a norm in $\Zomega$, we verify that a specific matrix given by an algebraic number $\beta\in\Zomega$ is diagonalizable. Lemma~\ref{lem:propertiesSbeta} summarizes also some other properties of this matrix and the norm which it induces according to Theorem~\ref{thm:norm}.
\begin{lem}
\label{lem:propertiesSbeta}
Let $\omega$ be an algebraic integer of degree $d$ and let $S$ be the companion matrix of its minimal polynomial $m_\omega$. Let $\beta=\sum_{i=0}^{d-1} b_i \omega^i$, where $b_i \in \ZZ$, be a nonzero element of $\Zomega$. Set $S_\beta=\multMat{b}$. Then
\begin{enumerate}[i)]
   \item The matrix $S_\beta$ is diagonalizable.
   \item The characteristic polynomial of $S_\beta$ is $m_\beta^k$ with $k=d / \deg \beta$.
   \item $|\det S_\beta|=|m_\beta(0)|^k$.
   \item $\norm{x}{S_\beta}=\norm{x}{S_\beta^{-1}}$ for all $x \in \CC^d$ and $\Mnorm{X}{S_\beta}=\Mnorm{X}{S_\beta^{-1}}$ for all $X \in \CC^{d\times d}$.
   \item $\Mnorm{S_\beta}{S_\beta}=\max \{|\beta'| \colon \beta' \text{ is conjugate of } \beta\}$ and $ \Mnorm{S_\beta^{-1}}{S_\beta}=\max \{\frac{1}{|\beta'|} \colon \beta' \text{ is conjugate of } \beta\}$.
\end{enumerate}  
\end{lem}
\begin{proof}
The characteristic polynomial of the companion matrix $S$ is the same as minimal polynomial of $\omega$ which has no multiple roots. Hence, $S$ is diagonalizable, i.e., $S=P^{-1}DP$ where $D$ is diagonal matrix with the conjugates of $\omega$ on the diagonal and $P$ is a nonsingular complex matrix. The matrix $S_\beta$ is also diagonalized by $P$:
$$
S_\beta=\sum_{i=0}^{d-1} b_i S^i= \sum_{i=0}^{d-1} b_i \left(P^{-1}DP\right)^i=P^{-1}\underbrace{\left(\sum_{i=0}^{d-1} b_i D^i\right)}_{D_\beta}P\,.
$$
There is a result in theory of algebraic numbers about conjugates under a field isomorphism -- let $\alpha$ be an algebraic number with a conjugate $\alpha'$ and $\sigma:\QQ(\alpha)\rightarrow\QQ(\alpha')$ be a field isomorphism. If $\gamma\in\QQ(\alpha)$, then $\sigma(\gamma)$ is a conjugate of $\gamma$. \komentar{je potreba nejakou citaci?}

Therefore, the diagonal elements of the diagonal matrix $D_\beta$ are conjugates of $\beta$. Since $S_\beta\in\ZZ^{d\times d}$, its characteristic polynomial $p_{S_\beta}$ has integer coefficients. There exists $k\in\NN, k\geq 1$ such that $p_{S_\beta}=m^k_\beta$ as all roots of $p_{S_\beta}$ must be conjugates of $\beta$. The value $k$ follows from the equality $d=\deg(m_\beta^k)=k \deg m_\beta$. 

The modulus of the determinant of $S_\beta$ equals the modulus of the absolute coefficient of the characteristic polynomial $p_{S_\beta}$ which is $|m_\beta(0)|^k$.

The matrix $S_\beta^{-1}$ is also diagonalized by $P$ since $S_\beta^{-1}=(P^{-1}D_\beta P)^{-1}=P^{-1}D_\beta^{-1}P$. Thus, the norms $\norm{\cdot}{S_\beta}$ and $\norm{\cdot}{S_\beta^{-1}}$ are the same and so are the induced matrix norms $\Mnorm{\cdot}{S_\beta}$ and $\Mnorm{\cdot}{S_\beta^{-1}}$.

The matrix $S_\beta$ is diagonalizable and its eigenvalues are the conjugates of $\beta$. Theorem~\ref{thm:norm} implies that 
$$
\Mnorm{S_\beta}{S_\beta}=\rho(S_\beta)= \max \{|\beta'| \colon \beta' \text{ is conjugate of } \beta\}\,. 
$$
For the second part of the last statement, we use the part \textit{iv)}, Theorem~\ref{thm:norm} and the fact that the eigenvalues of $S_\beta^{-1}$ are  reciprocal of the conjugates of $\beta$.
\end{proof}

Finally, we may define a norm in $\Zomega$.
\begin{defn}
Let $\pi$ be the isomorphism between $\Zomega$ and $(\ZZ^d,+,\odot_\omega)$. Using the notation of the previous lemma, we define \emph{$\beta$-norm}  $\normBeta{\cdot}:\Zomega \rightarrow \RR^+_0$ by 
$$
\normBeta{x}=\norm{\pi(x)}{S_\beta}
$$
for all $x\in\Zomega$.
\end{defn}
We can easily verify that $\beta$-norm is a norm in $\Zomega$:
\begin{enumerate}
    \item $\normBeta{x}=\norm{\pi(x)}{S_\beta}\geq 0\,,$
    \item $\normBeta{x}=0 \iff \norm{\pi(x)}{S_\beta}=0 \iff \pi(x)=0 \iff x=0\,,$
    \item $\normBeta{\alpha x}=\norm{\pi(\alpha x)}{S_\beta}=|\alpha|\norm{\pi(x)}{S_\beta}=|\alpha|\normBeta{x}\,,$
    \item $\normBeta{x+y}=\norm{\pi(x+y)}{S_\beta}=\norm{\pi(x)+\pi(y)}{S_\beta}\leq \norm{\pi(x)}{S_\beta}+\norm{\pi(y)}{S_\beta}=\normBeta{x}+\normBeta{y}\,,$
\end{enumerate}
for all $x,y \in \Zomega$ and $\alpha \in \Zomega$.

The important property of $\beta$-norm is that there is only finitely many elements of $\Zomega$ which are bounded in this norm by a given constant. The explanation is following -- images of elements of $\Zomega$ under the isomorphism $\pi$ are integer vector and there are only finitely many integer vectors in any finite dimensional vector space bounded by any norm. It follows from equivalence of all norms a finite dimensional vector space.


	\section{Number of congruence classes}

\begin{defn}
Let $M\in\ZZ^{d\times d}$ be a nonsingular integer matrix. Vectors $x,y \in \ZZ^d$ are \emph{congruent modulo $M$ in $\ZZ^d$}, if $x-y \in M\ZZ^d$.
\end{defn}

\begin{lem}
Let $M\in\ZZ^{d\times d}$ be a nonsingular integer matrix. The number of congruence classes modulo $M$ in $\ZZ^d$ is $|\det M|$.
\end{lem}
\begin{proof}
Set $y_i:=M e_i$ for $i\in\{0, \dots, d-1 \}$ and 
$$
B_{\enum{\alpha}}:=\left\{\sum_{i=0}^{d-1} (\alpha_i + t_i) y_i \colon t_i \in [0,1)\right\}
$$
for $\enum{\alpha} \in \ZZ^d$. Obviously,
$$
\RR^d=\bigcup_{\enum{\alpha} \in \ZZ^d} B_{\enum{\alpha}}\,.
$$
For fixed $\enum{\alpha} \in \ZZ^d$, the number of points of $\ZZ^d$ in $B_{\enum{\alpha}}$  is volume of $B_{\enum{\alpha}}$  which is $|\det M|$. Hence, it is enough to prove that there is exactly one representative of each congruence class in $B_{\enum{\alpha}}$. 

To show that there are representatives of all classes, assume $x\in\ZZ^d$. Since $\enum{y}$ is a basis of $\RR^d$, there are scalars $s_0, \dots, s_{d-1}\in \RR$ such that $x= \sum_{i=0}^{d-1} s_i y_i$. Set $\gamma_i:=\lfloor s_i \rfloor$ and $t_i:=s_i-\gamma_i$. Now
\begin{align*}
 x=\sum_{i=0}^{d-1} (\gamma_i+t_i) y_i=\sum_{i=0}^{d-1}t_i y_i +\sum_{i=0}^{d-1} (\gamma_i-\alpha_i)y_i + \sum_{i=0}^{d-1} \alpha_i y_i=\underbrace{\sum_{i=0}^{d-1} (\alpha_i+t_i)y_i}_{\in B_{\enum{\alpha}} } +M\underbrace{(\gamma-\alpha)}_{\in\ZZ^d}\,,
\end{align*}
where
$$
\alpha=\enum{\alpha}^T \quad \text{ and }\quad \gamma=\enum{\gamma}^T\,.
$$
Let $x=\sum_{i=0}^{d-1} s_i y_i \in \ZZ^d$ and $x'=\sum_{i=0}^{d-1}s'_i y_i \in \ZZ^d$ be distinct elements of $B_{\enum{\alpha}}$ which are congruent modulo $M$, i.e., there exists $z=\enum{z}^T\in\ZZ^d$ such that $x=x'+M z$. There is $i_0\in\{0, \dots , d-1\}$ such that $|z_{i_0}|\geq 1$ as $x\neq x'$. Thus, $|s_{i_0}-s'_{i_0}|=|z_{i_0}|\geq 1$ which contradicts that  $x, x'\in B_{\enum{\alpha}}$.
\end{proof}

\begin{thm}
Let $\omega$ be an algebraic integer of degree $d$ and  $\beta$ be an element of $\Zomega$ such that $\deg \omega = \deg\beta$. The number of congruence classes modulo $\beta$ in $\Zomega$ is $|m_\beta(0)|$.
\end{thm}
\begin{proof}
\komentar{dopsat dukaz}
\end{proof}	
	
\chapter{Convergence}
We discuss convergence of the extending window method designed in Chapter~\ref{chap:ewm}. We prove a necessary condition of convergence of the whole method and we show that it is also a sufficient condition of convergence of Phase~1, using the tools from the previous chapter. Next, we establish a condition which implies non-convergence of Phase~2. The condition is formulated by existence of an infinite walk in a specific directed graph. We also prove that there is the same lower bound on the size of an alphabet from $\Zbeta$ as for integer alphabets when the extending method converges. We propose a way how an alphabet can be generated for a given base.
\label{chap:convergence}

	\section{Convergence of Phase 1}
	\begin{lem}
\label{lem:suffCondPhase1}
    Let $\omega$ be an algebraic integer, $\deg \omega=d$, and $\beta$ be an expanding algebraic integer in $\Zomega$. Let $\A$ and $\B$ be finite subsets of $\Zomega$ such that $\A$ contains at least one representative of each congruence class modulo $\beta$ in $\Zomega$. Then there exists a finite set $\Q\subset\Zomega$  such that $ \B + \Q \subset \A + \beta \Q$.
\end{lem}

\begin{proof}
We use the isomorphism $\pi:\Zomega \rightarrow \ZZ^{d}$ and $\beta$-norm $\normBeta{\cdot}$ to bound the elements of $\Zomega$.
Let $\gamma$ be the smallest conjugate of $\beta$ in modulus. 
 Denote $C:=\max\{\normBeta{b-a}\colon a \in \A, b \in \B\}$. Consequently, set $R:=\frac{C}{|\gamma|-1}$ and $\Q:=\{q\in\Zomega\colon \normBeta{q}\leq R\}$. By Lemma~\ref{lem:propertiesSbeta}, we have 
 $$
 \Mnorm{S_\beta^{-1}}{S_\beta}=\max \{\frac{1}{|\beta'|} \colon \beta' \text{ is conjugate of } \beta\}=\frac{1}{|\gamma|}\,.
 $$ 
 Also, $|\gamma|>1$ as $\beta$ is an expanding integer.  Since $C>0$, the set $\Q$ is nonempty. Any element $x=b+q \in \Zomega$ with $b\in\B$ and $q\in\Q$ can be written as $x=a+\beta q'$ for some $a\in\A$  and $q'\in\Zomega$ due to existence of a representative of each congruence class in $\A$. Using the isomorphism $\pi$, we may write $\pi(q')=S^{-1}_\beta \cdot \pi(b-a+q)$. We prove that $q'$ is in $Q$:
\begin{align*}
    \normBeta{q'}&=\norm{\pi(q')}{S_\beta}=\norm{S^{-1}_\beta \cdot \pi(b-a+q)}{S_\beta}\leq \Mnorm{S^{-1}_\beta}{S_\beta}  \normBeta{b-a+q} \\
    &\leq  \frac{1}{|\gamma|} (\normBeta{b-a} +\normBeta{q})=\frac{1}{|\gamma|} (C+R)=\frac{C}{|\gamma|} (1+\frac{1}{|\gamma|-1}) =R\,.
\end{align*}
 
 Hence $q'\in\Q$ and thus  $x=b+q \in \A + \beta \Q$. 
 
 Since there are only finitely many elements of $\ZZ^{d}$ bounded by the constant $R$, the set $Q$ is finite.
\end{proof}

\begin{theo}
\label{thm:suffCondPhase1}
Let $\omega$ be an algebraic integer and $\beta\in\Zomega$. Let the alphabet $\A\subset\Zomega$ be such that $\A$ contains at least one representative of each congruence class modulo $\beta$ in $\Zomega$. Let $\B\subset\Zomega$ be the input alphabet. 

If $\beta$ is expanding, Phase 1 of the extending window method converges.
\end{theo}
\begin{proof}
We have the constant $R$ and finite set $\Q$ from Lemma \ref{lem:suffCondPhase1} for the alphabet $\A$ and input alphabet $\B$. We prove by induction that  all intermediate weight coefficient sets $\Q_k$ in Algorithm \ref{alg:weightCoefSet} are subsets of the finite set $\Q$. 

We start with $\Q_0=\{0\}$ which is bounded by any positive constant. Suppose that the intermediate weight coefficients set $\Q_k$ has elements bounded by the constant $R$. We see from the previous  proof that the candidates obtained by Algorithm \ref{alg:searchCand} for the set $\Q_k$ are also bounded by $R$. Thus, the next intermediate weight coefficients set $\Q_{k+1}$ has elements bounded by the constant $R$, i.e., $\Q_{k+1}\subset\Q$. 

Since $\#\Q$ is finite and $\Q_0\subset\Q_1\subset\Q_2\subset\cdots\Q$,  Phase 1 successfully ends. 
\end{proof}




	\section{Convergence of Phase 2}
	\label{sec:convergencePhase2}

We have no simple sufficient or necessary condition of convergence of Phase~2 in the sense of properties of a base $\beta$ or an alphabet $\A$. Nevertheless, the convergence can be controled during a run of algorithm. An easy check of non-convergence can be done by searching $\Q_{[b, \dots, b]}$ for each $b\in\B$ separately. This was already described in \cite{vu}, but we recall it with a simpler proof. For its purposes, we introduce a notion of stable Phase~2, which is used also in the main result of this section -- the control of convergence during Phase 2 is transformed into searching for a cycle in an oriented graph.

Firstly, we menshion some equivalent conditions of non-convergence of Phase~2. It enables us to handle easier with non-convergence in later  proofs.
\begin{lem}
\label{lem:equivalentStatementsForNonConvergenePhaseTwo}
The following statements are equivalent:
\begin{enumerate}[i)]
	\item Phase~2 does not converge,
	\item $\forall \,k\in \NN \,\exists\, \tupleo{k}\in\B^{k+1} \colon \#\Qwo{k}\geq 2$,
	\item $\exists \,(w_{-k})_{k\geq 0}, w_{-k}\in\B \,\exists\, k_0\forall k\geq k_0 \colon \#\Qwo{k}=\#\Qwo{(k-1)}\geq 2$.
\end{enumerate}
\end{lem}
\begin{proof}
\textit{i)}$\iff$\textit{ii):} The while loop in Algorithm~\ref{alg:weightFunction} ends if and only if there exist $k\in\NN$ such that $\#\Qwo{k}=1$ for all $\tupleo{k}\in\B^{k+1}$.

\textit{ii)}$\iff$\textit{iii):} There is an infinite sequence $(w_{-k})_{k\geq 0}$ such that $\#\Qwo{k}\geq 2$ for all $k\in\NN$ since $\Qwo{k}\supset\Qwo{(k+1)}$. Hence, the sequence of integers $(\#\Qwo{k})_{k\geq 0}$ is eventually constant. The opposite implication is trivial.
\end{proof}

We need to ensure that choice of a possible weight coefficient set $\Qwo{k}\subset \Qwo{(k-1)}$ is determined by an input digit $w_0$ and a set $\Qw{1}{0}$, while the influence of the set $\Qwo{(k-1)}$ is limited. It is formalized in the following definition.

\begin{defn}
Let $\B$ be an alphabet of input digits. We say that Phase~2 is \emph{stable}, if the following condition on the sets of possible weight coefficients which are produced by the algorithm holds: 
$$
\Qw{1}{k}=\Qw{1}{(k-1)} \implies \Qwo{k}=\Qwo{(k-1)}
$$
for all $k\in\NN, k\geq 2$ and for all $\tupleo{(k+1)} \in \B^{k+1}$.

\komentar{ii)$\#\Qw{1}{k}=1 \implies \#\Qwo{k}=1$, neni potreba pokud nebude opacna implikace}
\end{defn}
The definition may seem too restrictive, but note that $\Qwo{k}$ if fully determined by $\Qw{1}{k}$, $\Qwo{(k-1)}$ and $w_0$ for a fixed deterministic way of choice of possible weight coefficients sets. Therefore, the assumption $\Qw{1}{k}=\Qw{1}{(k-1)}$ implies that the only difference in the choice of $\Qwo{k}$ an $\Qwo{(k-1)}$ is that $\Qwo{k}$ is a subset of  $\Qwo{(k-1)}$, while $\Qwo{(k-1)}$ is chosen as a subset of $\Qwo{(k-2)}$. At the same time, $\Qwo{(k-1)}$ is a subset of $\Qwo{(k-2)}$. Thus, the requirement that Phase~2 is stable means that the same possible weight coefficients set is found even if it is searched as a subset of smaller set. This is natural way how an algorithm of choice should be constructed -- the set $\Qwo{k}$ is searched such that
$$
\B+\Qw{1}{k} \subset \A+\beta \Qwo{k}\,,
$$
i.e., there is no reason to choose the set $\Qwo{k}\subsetneq\Qwo{(k-1)}$ as we know that
$$
\B+\underbrace{\Qw{1}{(k-1)}}_{=\Qw{1}{k}} \subset \A+\beta \Qwo{(k-1)}\,.
$$
In other word, if $\Qwo{k}\subsetneq\Qwo{(k-1)}$, the set $\Qwo{(k-1)}$ might have been chosen smaller.

Now we use that finiteness of Phase~2 implies that there exists a length of window $m$ such that the set $\Qb{m}$ contains only one element for all $b\in\B$, where $\Qb{m}$ is a shorter notation for
$$
\Q_{[\underbrace{\scriptstyle b,\dots,b}_m]}\,.
$$
The following theorem was proved in \cite{vu} with the assumption that Phase~2 is deterministic. Briefly, it says that $\#\Qb{m}$ must decrease every time we increase $m$, otherwise Phase~2 does not converge. When we consider only inputs of the form $bb\dots b$ for some $b\in\B$, determinism implies that Phase~2 is stable. The given proof with Phase~2 being stable is slightly shorter.

\begin{thm}
\label{thm:bbbCondition}
Let $m_0 \in \NN$ and $b\in\B$ be such that sets $\Qb{m_0}$ and $\Qb{m_0-1}$ produced by stable Phase~2 have the same size. Then
$$
    \#\Qb{m} = \#\Qb{m_0} \qquad \forall m\geq m_0-1\,.
$$ 
Particularly, if $\#\Qb{m_0}\geq 2$, Phase~2 does not converge.
\end{thm}
\begin{proof}
As $\Qb{m_0} \subset \Qb{m_0-1}$, the assumption of the same size implies
$$
    \Qb{m_0} = \Qb{m_0-1}\,.
$$
By the assumption that Phase~2 is stable, we have
\begin{align*}
 \Qb{m_0} = \Qb{m_0-1} &\implies  \Qb{m_0+1} = \Qb{m_0} \\
 						&\implies  \Qb{m_0+2} = \Qb{m_0+1} \\
 						&\vdots
\end{align*}
This implies the statement.

If $\#\Qb{m_0}\geq 2$, then the statement \textit{iii)} in Lemma~\ref{lem:equivalentStatementsForNonConvergenePhaseTwo} holds for the sequence $(b)_{k\geq 0}$.
\end{proof}



The condition of convergence during Phase~2 is formulated as a searching for an infinite path in a so-called Rauzy graph. This term comes from combinatorics on words. The vertices of our graph are combinations of input digits for which the size of their possible weight coefficients sets did not decreased with an increment of lenght of the window. Whereas in combinatorics on words, vertices are given as factors of some language. But the edges are placed identically -- if some combination of digits without the first one equals another one without the last digit.
\begin{defn}
Let $\B$ be an alphabet of input digits. Let Phase~2 is stable. Let $k\in\NN, k\geq 2$. We set
$$
V:=\left\{\tuple{1}{k}\in\B^k \colon \#\Qw{1}{k}=\#\Qw{1}{(k-1)}\right\}
$$
and
$$
E:=\left\{\tuple{1}{k}\rightarrow\tuple[w']{1}{k}\in V\times V \colon \tuple{2}{k}=\tuple[w']{1}{(k-1)}\right\}\,.
$$
The oriented graph $G_k=(V,E)$ is called \emph{Rauzy graph of Phase~2 (for the window of length~$k$)}.
\end{defn}

The structure of Rauzy graph $G_k$ signifies whether the non-decreasing combinations are such that they cause nonconvergence of Phase~2. Existence of an infinite walk in $G_k$ implies that Phase~2 does not converge:

\begin{thm}
\label{thm:infinitePathInRauzyGraph}
Let Phase~2 is stable.  If there exists $k_0\in\NN, k_0\geq 2$, and $\tupleo{k_0}\in\B^{k_0+1}$ such that
\begin{enumerate}[i)]
	\item $\#\Qwo{(k_0-1)}>1$ and
	\item there exist an infinite walk $(\tuple[w^{(i)}]{1}{k_0})_{i\geq 1}$ in $G_{k_0}$ which starts in the vertex  $$\tuple[w^{(1)}]{1}{k_0}=\tuple{1}{k_0}\,,$$ 
\end{enumerate}
then Phase~2 does not converge.
\end{thm}
\begin{proof}
Set
$$(w_k)_{k\geq 0}:=w_0, w_1^{(1)}, \dots , w^{(1)}_{k_0-1}, w^{(1)}_{k_0}, w_{k_0}^{(2)},w_{k_0}^{(3)},w_{k_0}^{(4)},\dots$$
We prove that $\#\Qwo{k}=\#\Qwo{(k_0-1)}>1$ for all $k\geq k_0-1$, i.e., the condition \textit{iii)} in Lemma~\ref{lem:equivalentStatementsForNonConvergenePhaseTwo} is satisfied.

%The general ideas of the proof are following: if $\tuple[w']{1}{k_0}$ is a vertex of $G_{k_0}$, then $$\Qw[w']{1}{(k_0-1)}=\Qw[w']{1}{k_0}\,.$$ And the assumption that Phase~2 is stable implies that $$\Qw[w']{1}{(k_0-1)}=\Qw[w']{1}{k_0}\implies \Qwo[w']{(k_0-1)}=\Qwo[w']{k_0}$$ for all $w'_0\in\B$.

%Now we show that $\Qwo{k}=\Qwo{(k+1)}$ for all $k\geq k_0-1$. 
Let $l\in\NN$. Since $\tuple{(1+l)}{(k_0+l)}$ is a vertex of $G_{k_0}$, the set $\Qw{l}{(k_0+l)}$ equals $\Qw{l}{(k_0+l-1)}$. As Phase~2 is stable, we have
\begin{align*}
\Qw{l}{(k_0+l)}&=\Qw{l}{(k_0+l-1)}\\
\implies \Qw{(l-1)}{(k_0+l)}&=\Qw{(l-1)}{(k_0+l-1)}\\
&\vdots \\
\implies \Qw{1}{(k_0+l)}&=\Qw{1}{(k_0+l-1)}\\
\implies \Qwo{(k_0+l)}&=\Qwo{(k_0+l-1)}\,.
\end{align*}
Hence, $\#\Qwo{k}=\#\Qwo{(k_0-1)}>1$ for all $k\geq k_0-1$.
\end{proof}
\komentar{bylo by fajn dokazat i opacny smer, jedine, pres co se neumim dostat je kdyby existovala jen aperiodicka posloupnoust, kvuli ktere to nekonverguje}



We remark that existence of an infinite walk in a finite graph is equivalent to existence of a cycle in the graph. Thus, if there is an infinite walk, we may find another one whose sequence of vertices is eventually periodic. We use this fact in Section~\ref{sec:modifiedPhase2} which describes modified algorithm for Phase~2 which implements the result of Theorem~\ref{thm:infinitePathInRauzyGraph}.










	

	\section{Minimal alphabet $\A$}
\label{sec:minimalAlphabet}	

Frougny, Pelantov\'a and Svodov\'a \cite{minAlph} proved a lower bound on the size of an alphabet $0\in\A\subset\ZZ$ of consecutive integers which enables parallel addition. In this section, we prove the same bound for an arbitrary alphabet $\A\in\Zbeta$.  We recall their auxiliary results in Theorem~\ref{thm:reprBetaMinusOne} an Lemma~\ref{lem:alphabetRestrictions}, but only for a parallel digit set conversion without anticipation as our rewriting rule $x-\beta$ does not require memory. 

We remark that we indicate by the assumption $\A\in\Zbeta$  that we work in $\Zbeta$ instead of $\Zomega$. Notice that congruence classes modulo $\beta$ in $\Zomega$ and $\Zbeta$ are generally different. It implies that even an integer alphabet behaves differently if $\beta\in\Zomega$ and $\Zomega\neq\Zbeta$. On the other hand, if $\beta=\pm \omega +c$, where $c\in\ZZ$, then $\Zomega=\Zbeta$

The following theorem says that all classes modulo $\beta$ in $\Zomega$ which are contained in $\A+\A$ must have their representatives in $\A$.
\begin{thm}
\label{thm:reprBetaMinusOne}
Let $\omega$ be an algebraic integer. Let the base $\beta\in\Zomega$ be such that $|\beta|>1$ and the alphabet $\A\subset\Zomega$ be such that $0\in\A$. If there exists a $p$-local digit set conversion defined by the function $\phi\colon (\A+\A)^p\rightarrow \A$ and $p=r+1$, then the number $\phi(b,\dots,b)-b$ belongs to the set $(\beta-1)\Zomega$ for any $b\in\A+\A$. 
\end{thm}
\begin{proof}
Let $b\in\A+\A$ and $a=\phi(b, \dots,b)$. For $n\in\NN, n\geq 1$, we denote $S_n$ the number represented by
$$
^{\omega}\!0 \underbrace{b\dots b}_{n}\bullet \underbrace{b\dots b}_{r}0^\omega\,.
$$
The representation of $S_n$ after the digit set conversion is of the form
$$
^{\omega}\!0 \underbrace{w_{r}\dots w_{1}}_{\beta^n W}\underbrace{a\dots a}_{n}\bullet \underbrace{\widetilde{w_1}\dots \widetilde{w_r}}_{\beta^{-r}\widetilde{W}}0^\omega\,,
$$
where 
$$W=\sum_{j=1}^r w_j \beta^{j-1} \qquad \text{and} \qquad \widetilde{W}=\sum_{j=1}^r\widetilde{w_j} \beta^{r-j}\,.$$
Since both representations have same value, we have
\begin{align}
\label{eq:reprBetaMinusOne}
b \sum_{j=-r}^{n-1} \beta^j &= W \beta^n + a \sum_{j=0}^{n-1} \beta^j + \beta^{-r}\widetilde{W} \notag \\
b \sum_{j=-r}^{-1} \beta^j +b\frac{\beta^n-1}{\beta-1} &= W \beta^n + a \frac{\beta^n-1}{\beta-1} + \beta^{-r}\widetilde{W}\,,
\end{align}
for all $n\geq 1$. We subtract this equation for $n$ and $n-1$ to obtain
$$
b\frac{\beta^n-\beta^{n-1}}{\beta-1}=W(\beta^n-\beta^{n-1}) + a\frac{\beta^n-\beta^{n-1}}{\beta-1}\,.
$$
We simplify it to
\begin{equation}
\label{eq:reprBetaMinusOneFinal}
b=W(\beta-1) + a\,.
\end{equation}
Hence, $a=\phi(b, \dots,b)\equiv b$ modulo $\beta-1$.
\end{proof}

If a base $\beta$ has a real conjugate greater than one, then there are some extra requirements on the alphabet $\A$. For simplicity, we assume that the base $\beta$ itself is real and greater than one. We show later that this assumption is without loss of generality.
\begin{lem}
\label{lem:alphabetRestrictions}
Let $\omega$ be a real algebraic integer and the base $\beta\in\Zomega$ be such that $\beta>1$. Let the alphabet $\A\subset\Zomega$ be such that $0\in\A$ and denote $\lambda=\min \A$ and $\Lambda=\max \A$. If there exists a $p$-local digit set conversion defined by the function $\phi\colon (\A+\A)^p\rightarrow \A$ and $p=r+1$, then:
\begin{enumerate}[i)]
	\item $\phi(b,\dots,b)\neq \lambda$ for all $b\in\A+\A$ such that $b>\lambda$.
	\item $\phi(b,\dots,b)\neq \Lambda$ for all $b\in\A+\A$ such that $b<\Lambda$.
	\item If $\Lambda\neq 0$, then $\phi(\Lambda,\dots,\Lambda)\neq \Lambda$.
	\item If $\lambda\neq 0$, then $\phi(\lambda,\dots,\lambda)\neq \lambda$.
\end{enumerate}
\end{lem}
\begin{proof}
To prove \textit{i)}, assume in contradiction that $\phi(b,\dots,b)= \lambda$. We proceed in the same manner as in Theorem \ref{thm:reprBetaMinusOne}, the equation \eqref{eq:reprBetaMinusOne} implies
$$
b \sum_{j=-r}^{-1} \beta^j +b\frac{\beta^n-1}{\beta-1} =  \beta^n W + \lambda \frac{\beta^n-1}{\beta-1} + \beta^{-r}\widetilde{W}\,.
$$
We apply also the equation c to obtain
\begin{align*}
b \sum_{j=-r}^{-1} \beta^j +b\frac{\beta^n-1}{\beta-1} &=  \beta^n \frac{b-\lambda}{\beta-1} + \lambda \frac{\beta^n-1}{\beta-1} + \beta^{-r}\widetilde{W}\,.
\end{align*}
Now we may simplify and estimate
\begin{align*}
b \sum_{j=-r}^{-1} \beta^j +\frac{-b}{\beta-1} &=  \frac{-\lambda}{\beta-1} + \beta^{-r}\sum_{j=1}^r\widetilde{w_j} \beta^{r-j} \\
b \underbrace{\left(\sum_{j=1}^{r} \frac{1}{\beta^j} -\frac{1}{\beta-1}\right)}_{-\sum_{j=r+1}^{\infty} \frac{1}{\beta^j}} &=  -\lambda\frac{1}{\beta-1} + \sum_{j=1}^r\widetilde{w_j} \beta^{-j}\geq  \lambda\underbrace{\left(-\frac{1}{\beta-1}+\sum_{j=1}^{r} \frac{1}{\beta^j} \right)}_{-\sum_{j=r+1}^{\infty}\frac{1}{\beta^j}}\,.
\end{align*}
Hence $b\leq\lambda$ which is a contradiction. The proof of \textit{ii)} can be done in the same way.

For \textit{iii)}, assume that $\phi(\Lambda,\dots,\Lambda)= \Lambda$. Now consider a number $T_q$ represented by
$$
^{\omega}\!0 \bullet \underbrace{\Lambda\dots\Lambda}_{r} \underbrace{(2\Lambda)\dots(2\Lambda)}_{q} 0^\omega\,.
$$
After the digit set conversion, a representation is
$$
^{\omega}\!0  \underbrace{w_r\dots w_1}_{W} \bullet z_1\dots z_{r+q} 0^\omega\,.
$$
The value $T_q$  preserves, thus,
$$
\Lambda\sum_{j=1}^r \beta^{-j} +2\Lambda \sum_{j=r+1}^{r+q} \beta^{-j}=W+\sum_{j=1}^{r+q}z_j\beta^{-j}\,.
$$
But $W=0$ from the equation \eqref{eq:reprBetaMinusOneFinal}. We estimate
\begin{align*}
\Lambda\sum_{j=1}^{r+q} \beta^{-j} +\Lambda \sum_{j=r+1}^{r+q} \beta^{-j}&=\sum_{j=1}^{r+q}z_j\beta^{-j}\leq \Lambda\sum_{j=1}^{r+q}\beta^{-j}\\
\Lambda \sum_{j=r+1}^{r+q} \beta^{-j}&\leq 0\,.
\end{align*}
This contradicts that $\Lambda$ is positive as it is a nonzero, maximal element of the alphabet $\A$ which contains 0. The proof of \textit{iv)} is analogous.
\end{proof}

In order to prove the lower bound, we need to show that the alphabet $\A$ must contain all representatives modulo $\beta$ and $\beta-1$. 
\begin{thm}
Let $\beta$ be an algebraic integer such that $|\beta|>1$. Let $0\in \A\subset \Zbeta$ be an alphabet such that $1\in \A[\beta]$. If addition in the numeration system $(\beta, \A)$ which uses the rewriting rule $x-\beta$ is computable in parallel, then the alphabet $\A$ contains at least one representative of each congruence class modulo $\beta$ and $\beta-1$ in $\Zbeta$. 
\label{thm:representativesInAlphabet}
\end{thm}
\begin{proof}
The existence of an algorithm for addition with the rewriting rule $x-\beta$ implies that the set $\A[\beta]$ is closed under addition. By the assumption $1\in \A[\beta]$, the set $\NN$ is subset of  $\A[\beta]$. Since $0\in\A$, we have $\beta \cdot \A[\beta] \subset \A[\beta]$. Hence, $\NN[\beta] \subset \A[\beta]$.

For any element  $x=\sum_{i=0}^N x_i \beta^i\in \Zbeta$ there is an element $x'=\sum_{i=0}^N x'_i \beta^i\in \NN[\beta]$ such that $x\equiv_\beta x'$  since $m_\beta (0)\equiv_\beta 0$ and $\beta^i\equiv_\beta 0$. As $x'\in \NN[\beta] \subset \A[\beta]$, we have
$$
x\equiv_\beta x'=\sum_{i=0}^{n}a_i\beta^i \equiv_\beta a_0\,,
$$
where $a_i\in \A$. Hence, for any element $x\in\Zomega$, there is a letter $a_0\in\A$ such that $x\equiv_\beta a_0$.

In order to prove that there is at least one representative of each congruence class modulo $\beta-1$ in the alphabet $\A$, we consider again an element $x=\sum_{i=0}^N x_i \beta^i\in \Zbeta$. Similarly, there is an element $x'=\sum_{i=0}^N x'_i \beta^i\in \NN[\beta]$ such that $x\equiv_{\beta-1} x'$  since $m_{\beta-1} (0)\equiv_{\beta-1} 0$ and $(\beta-1)^i\equiv_{\beta-1} 0$.

Since $x'\in \NN[\beta]\subset \A[\beta]$,
$$
x'=\sum_{i=0}^{n}a_i\beta^i\,,
$$
where $a_i\in \A$. We prove by induction with respect to $n$ that $x'\equiv_{\beta-1} a$ for some $a\in\A$.
If $n=0$, $x'=a_0$. Now we use the fact that if there is a parallel addition algorithm, for each letter $b \in\A+\A$, there is $a\in\A$ such that $b \equiv_{\beta-1} a$ (Theorem~\ref{thm:reprBetaMinusOne}). For $n+1$, we have
\begin{align*}
x'&=\sum_{i=0}^{n+1}a_i\beta^i =a_0 + \sum_{i=1}^{n+1}a_{i}\beta^i\\
    &=a_0 + \beta \sum_{i=0}^{n}a_{i+1}\beta^i - \sum_{i=0}^{n}a_{i+1}\beta^i+ \sum_{i=0}^{n}a_{i+1}\beta^i \\
    &\equiv_{\beta-1} a_0 + (\beta-1)\sum_{i=0}^{n}a_{i+1}\beta^i + a \equiv_{\beta-1}a_0 +a \equiv_{\beta-1}a' \in\A\,,
\end{align*}
where we use the induction assumption
$$
\sum_{i=0}^{n}a_{i+1}\beta^i\equiv_{\beta-1} a\,.
$$
\end{proof}
Unfortunately, the claim cannot be generalized to modulo in $\Zomega$ -- there are numeration systems with integer alphabets which allow parallel addition, but these alphabets does not contain all representatives modulo $\beta-1$ in $\Zomega$, see Table~\ref{tab:resultsQuadrInt} and Examples \ref{ex:integerAB},  \ref{ex:integerAJ}, \ref{ex:integerAO} and \ref{ex:integerAR}. Nevertheless, each class modulo $\beta -1$ in $\Zomega$ which is contained in $\A+\A$ must still has its representative in $\A$ according to Theorem~\ref{thm:reprBetaMinusOne}.

The following lemma summarizes that if we have a parallel addition algorithm for a base $\beta$, then we easily obtain an algorithm also for conjugates of $\beta$.
\begin{lem}
\label{lem:parAddAlgForConjugate}
Let $\omega$ be an algebraic integer with a conjugate $\omega'$. Let $\beta\in\Zomega, |\beta|>1$ and let $\sigma:\QQ(\omega)\rightarrow \QQ(\omega')$ be an isomorphism such that $|\sigma(\beta)|>1$. Let $\varphi$ be a digit set conversion  in the base $\beta$ from $\A+\A$ to $\A$. There exists  is a digit set conversion $\varphi'$ in the base $\beta'$ from $\A'+\A'$ to $\A'$ where $\beta'=\sigma(\beta)$ and $\A'=\{\sigma(a) \colon a\in\A\}$.
\end{lem}
\begin{proof}
Let $\phi:\A^p\rightarrow\A$ be a mapping which defines $\varphi$ with $p=r+t+1$. We define a mapping $\phi':\A^p\rightarrow \A$ by 
$$
\phi'(w'_{j+t}, \dots, w'_{j-r})=\sigma\left(\phi\left(\sigma^{-1}(w'_{j+t}), \dots, \sigma^{-1}(w'_{j-r})\right)\right)\,.
$$
Next, we define a digit set conversion  $\varphi':(\A'+\A')\rightarrow\A'$ by $\varphi'(w')=(z'_j)_{j\in\ZZ}$ where $w'=(w'_j)_{j\in\ZZ}$ and $z'_j=\phi'(w'_{j+t}, \dots, w'_{j-r})$. Obviously, if $w'$ has only finitely many nonzero entries, then there is only finitely many nonzeros in $(z'_j)_{j\in\ZZ}$   since
$$
\phi'(0, \dots, 0)=\sigma\left(\phi\left(\sigma^{-1}(0), \dots, \sigma^{-1}(0)\right)\right)=\sigma\left(\phi\left(0, \dots, 0\right)\right)=\sigma\left(0\right)=0\,.
$$
The value of the number represented by $w'$ is also preserved:
\begin{align*}
\sum_{j\in\ZZ}w'_j {\beta'}^j&=\sum_{j\in\ZZ}\sigma(w_j) \sigma(\beta)^j=\sigma\left(\sum_{j\in\ZZ}w_j\beta^j\right) \\
&=\sigma\left(\sum_{j\in\ZZ}z_j\beta^j \right)=\sigma\left(\sum_{j\in\ZZ}\phi\left(w_{j+t}, \dots,w_{j-r}\right)\beta^j\right) \\
&=\sum_{j\in\ZZ}\sigma(\phi\left(w_{j+t}, \dots,w_{j-r}\right)){\beta'}^j=\sum_{j\in\ZZ}z'_j {\beta'}^j
\end{align*}
where $w_j=\sigma^{-1}(w'_j)$ for $j\in\ZZ$ and $\varphi((w_j)_{j\in\ZZ})=(z_j)_{j\in\ZZ}$.
\end{proof}

Finally, we put together that the alphabet $\A$ contains all representative modulo $\beta$ and $\beta-1$, number of congruence classes and restrictions on the alphabet for a base with a real conjugate greater than one.
\begin{thm}
Let $\beta$ be an algebraic integer such that $|\beta|>1$. Let $0\in \A\subset \Zbeta$ be an alphabet such that $1\in \A[\beta]$. If addition in the numeration system $(\beta, \A)$ which uses the rewriting rule $x-\beta$ is computable in parallel, then
$$
\#\A \geq \max \{|m_\beta(0)|, |m_\beta(1)|\}\,.
$$
Moreover, if $\beta$ is such that it has a real conjugate greater than 1, then 
$$
\#\A \geq \max \{|m_\beta(0)|, |m_\beta(1)|+2\}\,.
$$
\end{thm}
\begin{proof}
By Theorem~\ref{thm:representativesInAlphabet}, there are all representatives modulo $\beta$ and modulo $\beta-1$ in the alphabet $\A$. The numbers of congruence classes are $|m_\beta(0)|$ and $|m_{\beta-1}(0)|$ by Theorem~\ref{thm:numbCongruenceClasses}. Obviously, $m_{\beta-1}(x) = m_\beta (x+1)$. Thus $m_{\beta-1}(0) = m_\beta (1)$.

Let $\phi$ be a mapping which defines the parallel addition. According to Lemma~\ref{lem:parAddAlgForConjugate}, we may assume that $\beta$ is real and greater than 1 in the proof of the second part. The assumption $1\in \A[\beta]$ implies that $\Lambda>0$. Thus, there are at least three elements in the alphabet $\A$, because $\A\ni\phi(\Lambda,\dots,\Lambda)\neq \lambda$ and $\A\ni\phi(\Lambda,\dots,\Lambda)\neq \Lambda$ by Lemma \ref{lem:alphabetRestrictions}. It also implies that there are at least two representatives modulo $\beta-1$ in the alphabet in the class which contains $\Lambda$, since $\phi(\Lambda,\dots,\Lambda)\equiv_{\beta-1} \Lambda$.  

If $\lambda\equiv_{\beta-1}\Lambda$, there must be one more element of the alphabet $\A$ in this class, since $\lambda \neq \Lambda$. Therefore, $\#\A\geq |m_\beta(1)|+2$. 

The case that $\lambda\not\equiv_{\beta-1}\Lambda$ is divided into two subcases. Suppose now that $\lambda\neq 0$. Then $\phi(\lambda,\dots,\lambda)\neq \lambda$ and hence there is one more element in the alphabet in the class containing $\lambda$. Thus, there are at least two congruence classes which contain at least two elements of the alphabet $\A$. Therefore, $\#\A\geq |m_\beta(1)|+2$.

%If $\lambda=0$, then all elements of $\A+\A$ are nonnegative and $\phi(b,\dots,b)\neq 0$ for all $b\in(\A+\A)\setminus 0$. Suppose for contradiction, that there is only $|m_\beta(1)|+1$ elements in $\A$ -- one in each congruence class modeulo $\beta-1$ and one more in the class which contains $\Lambda$. The set $\mathcal{D}=\{\phi(d,\dots,d)\colon d\in\Lambda+\A\}\subset\A$ has $|m_\beta(1)|+1$ elements, but none of them is congruent to 0 as there are nonzero and the class containing zero has only one element by the assumption. Therefore, the elements of the set $\mathcal{D}$ belong to only $|m_\beta(1)|-1$ congruence classes. Hence, there exists $e,f,g,h\in\A, e\neq f,g\neq h, h\neq f$ such that $\phi(e+\Lambda,\dots, e+\Lambda)\equiv_{\beta-1}\phi(f+\Lambda,\dots, f+\Lambda)$ and $\phi(g+\Lambda,\dots, g+\Lambda)=\phi(h+\Lambda,\dots, h+\Lambda)$. Since 
%$$
%e+\Lambda\equiv_{\beta-1}\phi(e+\Lambda,\dots, e+\Lambda)\equiv_{\beta-1}\phi(f+\Lambda,\dots, f+\Lambda)\equiv_{\beta-1}f+\Lambda
%$$
%and
%$$
%g+\Lambda\equiv_{\beta-1}\phi(g+\Lambda,\dots, g+\Lambda)=\phi(h+\Lambda,\dots, h+\Lambda)\equiv_{\beta-1}h+\Lambda\,,
%$$ 
%also $e\equiv_{\beta-1}f$ and $g\equiv_{\beta-1}h$ which is a contradiction. In the same manner, the assumption that there is no nonzero element of the alphabet $\A$ which is congruent to 0 leads to contradiction for arbitrarily large alphabet.

If $\lambda=0$, then all elements of $\A+\A$ are nonnegative and $\phi(b,\dots,b)\neq 0$ for all $b\in(\A+\A)\setminus 0$. Suppose for contradiction, that there is no nonzero element of the alphabet $\A$ congruent to 0. We know that there is at least one representative of each congruence class class modulo $\beta-1$ in $\A$ and at least two representatives in the congruence class which contains $\Lambda$. Let $k\in\NN$ denote the number of elements which are in $\A$ extra to one element in each congruence class, i.e., $\#A=|m_\beta(1)|+k$.  For $d\in\Lambda+\A$, the value $\phi(d,\dots,d)\in\A$ is not congruent to 0 as it is nonzero and the class containing zero has only one element by the assumption. Therefore, the values  $\phi(d,\dots,d)\in\A$ for $|m_\beta(1)|+k$ distinct letters $d\in\Lambda+\A$ belong to only $|m_\beta(1)|-1$ congruence classes. Hence, there exists $e_1,\dots e_k, e_{k+1}\in\A$, pairwise distinct, and $f_1,\dots f_k, f_{k+1}\in\A$ such that $e_i\neq f_i$ and $\phi(e_i+\Lambda,\dots, e_i+\Lambda)\equiv_{\beta-1}\phi(f_i+\Lambda,\dots, f_i+\Lambda)$ for all $i, 1\leq i\leq k+1$. Since 
$$
e_i+\Lambda\equiv_{\beta-1}\phi(e_i+\Lambda,\dots, e_i+\Lambda)\equiv_{\beta-1}\phi(f_i+\Lambda,\dots, f_i+\Lambda)\equiv_{\beta-1}f_i+\Lambda\,.
$$ 
also $e_i\equiv_{\beta-1}f_i$ for all $i, 1\leq i\leq k+1$. This is a contradiction since it implies that $\#\A=|m_\beta(1)|+k+1$. Hence, classes containing $\lambda$ and $\Lambda$ have both at least one more  element of the alphabet $\A$, i.e., $\#\A\geq |m_\beta(1)|+2$. 
\end{proof}

Note that the larger necessary alphabet for a base with a real conjugate greater than one is caused by the fact that minimal and maximal element of $\A$ both require another element in $\A$ which is in the same congruence class. This should be considered when an alphabet for such a base is generated.

Since the proof is based on Theorem~\ref{thm:representativesInAlphabet}, it cannot be easily extended to alphabets which are subsets of  $\Zomega$.





\chapter{Different methods in~Phases~1~and~2}
	\label{chap:diffChoices}
	We mentioned in Chapter~\ref{chap:ewm} that there is a lot of variability in both phases of the extending window method. Within this chapter, we propose various methods how an intermediate weight coefficients set $\Q_k$ can be extended to $\Q_{k+1}$ in Phase~1 and how the set $\Qwo{k}$ can be constructed in Phase~2. The presented methods are based on many preliminary tests.

	\section{Different methods in Phase 1}
\label{sec:methodsOne}


We recall that the ambiguous part of Phase~1 is extending an intermediate weight coefficient set $\Q_k$ to $\Q_{k+1}$ so that 
$$
\B+\Q_k \subset \A+\beta \Q_{k+1}\,.
$$
It means that a weight coefficient $q$ such that $x=a+\beta q$ for some $a\in\A$ must be found for all $x\in \B+\Q_k$. Since the alphabet $\A$ is redundant, there might be more such weight coefficients. Let $C_x\subset \Zomega$ be a set of all these candidates for some $x\in \B+\Q_k$, i.e., $C_x=\{q\in\Zomega\colon x=a+\beta q, a\in\A\}$. Set $C:=\{C_x\colon x\in \B+\Q_k\}$. For a given $x\in \B+\Q_k$, the set $C_x$ is constructed so that all digits $a\in\A$ are tested whether $x-a$ is divisible by $\beta$ according to Theorem~\ref{thm:divisibility}. See Algorithm~\ref{alg:extendWeightCoefSet} which constructs the set $C$.

Now we extend the set $\Q_k$ to $\Q_{k+1}$ so that $C_x \cap \Q_{k+1} \neq \emptyset$ for all $C_x \in C$. We describe five methods how it can be done (1a, 1b, 1c, 1d and 1e).

As $Q_k\subset \Q_{k+1}$, set $\Q_{k+1}:=\Q_k$. For methods 1a, 1b and 1d, add  the element of all $C_x\in C$ such that $\#C_x=1$ to $\Q_{k+1}$. Next, select elements from all $C_x\in C$ such that $C_x \cap \Q_{k+1} = \emptyset$ according to a chosen method and add them to $\Q_{k+1}$. Selection for different methods is following:
\begin{enumerate}[ ]
	\item 1a -- all elements of $C_x$,
	\item 1b and 1c -- all smallest elements in absolute value of $C_x$,
	\item 1d and 1e -- all smallest elements in $\beta$-norm of $C_x$,
\end{enumerate}
Note that more elements may be added by method 1e than 1d, resp. 1c than 1b, since adding necessary elements before ($\#C_x=1$)  may cause that $C_{x'}\cap \Q_{k+1} \neq \emptyset$ for some $C_{x'}\in C$. 

The procedure is summarized in Algorithm~\ref{alg:extendWeightCoefSet}. 
Another methods may decrease the number of added elements for instance by picking only one of all smallest elements.
	


We can slightly improve performance by substituting the set $\B + \Q_{k}$ by $(\B + \Q_{k})\setminus (\B + \Q_{k-1})$ on the line~\ref{line:replaceBySmaller} in Algorithm~\ref{alg:searchCand}, because
$$
\B + \Q_{k-1} \subset \A + \beta \Q_{k}\,.
$$
Since $\Q_{k+1} \supset \Q_{k}$, $C_x \cap \Q_{k+1} \neq \emptyset$ for  $x\in \B + \Q_{k-1}$.

\begin{algorithm}
  \caption{Extending intermediate weight coefficients set}
    \label{alg:extendWeightCoefSet}
  \begin{algorithmic}[1]
    \REQUIRE{previous weight coefficients set $\Q_{k}$, method number $M\in\{$1a, 1b, 1c, 1d, 1e$\}$}
    \STATE $\widetilde{\Q}:=\Q_{k}$
    \IF {$M\in\{$1a,1b,1d$\}$}
    	\FORALL {$C_x \in C$}
    		\IF {$\# C_x =1$}
		    	\STATE Add the element of $C_x$ to $\widetilde{\Q}$
			\ENDIF
		\ENDFOR
	\ENDIF
    \STATE $\Q_{k+1}:=\widetilde{\Q}$
    \STATE By Algorithm~\ref{alg:searchCand}, find set of candidates $C$
    \FORALL{$C_x \in C$}
        \IF{$C_x \cap \widetilde{\Q}=\emptyset$}
	        \IF {$M=$ 1a}
	        	\STATE Add all elements of $C_x$ to $\Q_{k+1}$
	        \ELSIF {$M\in\{$1b, 1c$\}$}
	        	\STATE Add all smallest elements in absolute value of $C_x$ to $\Q_{k+1}$ 
	        \ELSIF {$M\in\{$1d, 1e$\}$}
	        	\STATE Add all smallest elements in $\beta$-norm of $C_x$ to $\Q_{k+1}$
	        \ENDIF
        \ENDIF
    \ENDFOR
    \RETURN $\Q_{k+1}$
  \end{algorithmic}
\end{algorithm}



\begin{algorithm}
  \caption{Search for set of candidates $C$}
    \label{alg:searchCand}
  \begin{algorithmic}[1]
    \REQUIRE{the previous weight coefficients set $\Q_{k}$, alternatively also the set $\Q_{k-1}$}
    \STATE $C:=\{\emptyset\}$
    \FORALL[Alternatively, $x \in (\B + \Q_{k})\setminus (\B + \Q_{k-1})$]{$x \in \B + \Q_{k}$} \label{line:replaceBySmaller}
      \STATE $C_x:=\emptyset$
      \FORALL{$a \in \A$}
          \IF{$(x-a)$ is divisible by $\beta$ in $\Zomega$ (using Theorem~\ref{thm:divisibility})}
              \STATE Add $\frac{x-a}{\beta}$ to $C_x$
            \ENDIF
      \ENDFOR 
      \STATE Add the set $C_x$ to $C$
  \ENDFOR
  \RETURN $C$
  \end{algorithmic}
\end{algorithm}  


	\section{Different methods in Phase 2}
\label{sec:methodsTwo}
We propose some methods which can be used in Phase~2 to construct a set $\Qwo{k}$.

The set of possible weight coefficients $\Qwo{k}$ must be a subset of the previous set of possible weight coefficients $\Qwo{(k-1)}$. It is constructed such that
$$
 w_0 + \Qw{1}{k} \subset \A + \beta \Qwo{k}\,,
$$
i.e, it is determined also by the input digit $w_0$ and the set of carries from the right $\Qw{1}{k}$. Whereas in Phase~1 the constructed set $\Q_{k+1}$ is extended from $\Q_k$, now we  reduce the set $\Qwo{(k-1)}$ to obtain $\Qwo{k}$. Set 
$$
D_x:= \left\{q_0\in\Qwo{(k-1)}\colon\, \exists \, a\in\A: x=a+\beta q_0 \right\}\,,
$$
where $x\in w_0 + \Qw{1}{k}$, and 
$$D:=\left\{D_x\colon x\in w_0 + \Qw{1}{k}\right\}\,.$$
 The goal is to find $\Qwo{k}$ such that $D_x \cap \Qwo{k} \neq \emptyset$ for all $x\in w_0 + \Qw{1}{k}$. In other words, if at least one element of each covering sets $D_x\in D$ is contained in $\Qwo{k}$, then 
$$
 w_0 + \Qw{1}{k} \subset \A + \beta \Qwo{k}\,.
$$
We follow Algorithm~\ref{alg:minimalSet} now. Choice of possible weight coefficients set starts with $\Qwo{k}':=\emptyset$. We add the elements of all $D_x\in D$ such that $\#D_x=1$ to $\Qwo{k}'$. We remove from covering set $D$ all sets $D_x\in D$ such that $\Qwo{k}' \cup D_x \neq \emptyset$. 

Next, while $D$ is nonempty, we pick an element $q$ from $\bigcup D$ according to a chosen method, we add $q$ to $\Qwo{k}'$ and we remove all sets $D_x$ which contain $q$ from $D$. Finally, the possible weight coefficient set $\Qwo{k}=\Qwo{k}'$ is found.

Before we explain different methods  to pick an element from $\bigcup D$, we define a linear order on $\Zomega$.

\begin{defn}
\label{def:lexicographical}
Let $n$ be a positive integer. Let $(x_1, \dots, x_n)$ and $(y_1, \dots, y_n)$ are elements of $\ZZ^n$. We say that $(x_1, \dots, x_n)\preceq (y_1, \dots, y_n)$ if 
$$
x_1<y_1 \text{ or } \left(x_1=y_1 \wedge (x_2, \dots, x_n)\preceq (y_2, \dots, y_n)\right)\,.
$$
If $x$ and $y$ are elements of $\Zomega$, then we say that $x$ is \emph{lexicographically smaller} than $y$ if
$$
\pi(x)^T \preceq \pi(y)^T\,,
$$
where $^T$ indicates transpose of column vector.
\end{defn}
Roughly speaking, elements of $\Zomega$ are first ordered according their absolute coeffient, then quadratic, etc.

Algorithm~\ref{alg:pickElement} outlines methods 9, 15, 22 and 23. All methods select a set $T$ of elements of $\bigcup D$. If there is more than one element in $T$, then the lexicographically smallest element is picked.

Method 9 computes the point of gravity $g$ of elements of $\bigcup D$ considered as complex numbers. The set $T$ are elements of  $\bigcup D$ which are closest to $g$ in absolute value. 

For remaining methods, set
$$D':=\left\{D_x \in D \colon \#D_x=\min_{D_x\in D} \#D_x \right\}\,.$$
Method 15 computes the point of gravity $g$ of already picked elements $\Qwo{k}'$ considered as complex numbers. The set $T$ are elements of  $\bigcup D'$ which are closest to $g$ in absolute value. 

Methods 22, resp. 23, builds the set $T$ to be the smallest elements of  $\bigcup D'$ in absolute value, resp.  $\beta$-norm. 

The motivation of these methods is to obtain the set $\Qwo{k}$ small and ``compact'' as it should be easier to cover them in the next iteration of Phase~2. Moreover, methods 15, 22 and 23 prefer elements from the smallest covering sets $D_x$ in the hope that the picked elements are also in the bigger covering sets.

Other implemented methods can be found in the source code. 

\begin{algorithm}
  \caption{Search for set $\Qwo{k}$ }
    \label{alg:minimalSet}
  \begin{algorithmic}[1]
    \REQUIRE{Input digit $w_0$, set of possible carries $\Qw{1}{k}$, previous set of possible weight coefficients $\Qwo{(k-1)}$}
    \STATE $D:=\{\emptyset\}$
    \FORALL{$x \in w_0 + \Qw{1}{k}$}
    	\STATE $D_x:= \{q_0\in\Qwo{(k-1)}\colon\, \exists \, a\in\A: x=a+\beta q_0 \}$
        \STATE Add $D_x$ to $D$
    \ENDFOR
    \STATE $\Qwo{k}':=\emptyset$
    \FORALL{$D_x\in D$}
    	\IF {$\#D_x=1$}
    		\STATE Add the element $q\in D_x$ to $\Qwo{k}'$
	        \STATE Remove sets $D_{x'}$ such that $q \in D_{x'}$  from the set $D$ 
	    \ENDIF
	\ENDFOR
    \WHILE{$D\neq\emptyset$}
        \STATE By Algorithm~\ref{alg:pickElement},  pick an element $q$ from $\bigcup D$
        \STATE Add the element $q$ to $\Qwo{k}'$
        \STATE Remove sets $D_x$ such that $q \in D_x$  from the set $D$ 
    \ENDWHILE
	\STATE $\Qwo{k}:=\Qwo{k}'$
    \RETURN $\Qwo{k}$
  \end{algorithmic}
\end{algorithm}


\begin{algorithm}
   \caption{Choose one element from the set of covering $D$ }
   \label{alg:pickElement}   
  \begin{algorithmic}[1]
    \REQUIRE{set of coverings $D$, method number $M\in\{9,15,22,23\}$, already added elements $\Qwo{k}'$}
    \IF {$M=9$}
    	\STATE $g:=$ point of gravity of elements of $\bigcup D$ as complex numbers
    	\STATE $T:=$ elements of $\bigcup D$ which are closest to $g$ in absolute value
    \ELSE
    	\STATE $m:=\min\{\#D_x \colon D_x\in D \}$
    	\STATE $D':=\{D_x \in D \colon \#D_x=m\}$
    	\IF {$M=15$}
    		\STATE $g:=$ point of gravity of elements of $\Qwo{k}'$ as complex numbers
    		\STATE $T:=$ elements of $\bigcup D'$ which are closest to $g$ in absolute value
		\ELSIF {$M=22$}
			\STATE $T:=$ elements of $\bigcup D'$ which are smallest in absolute value
    	\ELSIF {$M=23$}
			\STATE $T:=$ elements of $\bigcup D'$ which are smallest in $\beta$-norm
    	\ENDIF
    \ENDIF    
    \RETURN the lexicographically smallest element of $T$ according to Definition~\ref{def:lexicographical}

  \end{algorithmic}
\end{algorithm}




	
	


\chapter{Design and implementation}

In the first section of this chapter, we propose algorithms which may reveal non-convergence of Phase~2. They are based on the theorems from Chapter~\ref{chap:convergence}. We present a simple algorithm which makes Phase~2 stable. We develop Algorithm~\ref{alg:weightFunction_modified} that executes  Phase~2 with all modifications. All designed algorithms are implemented in SageMath. The program is described in Section~\ref{sec:implementation}.
\label{chap:design}
	\section{Modified Phase 2}
\label{sec:modifiedPhase2}

In this section, we introduce algorithms based on the theorems proved in Section~\ref{sec:convergencePhase2}. The first one checks convergence of Phase 2 for $bb\dots b$ inputs. Next, we show that it is possible to make Phase 2 stable by wrapping the choice of a set of possible weight coefficients into a simple while loop. Finally, we present an algorithm for Phase~2 which includes all modifications -- the mentioned check for $bb\dots b$ inputs and control of non-convergence by searching for an infinite walk in Rauzy graph $G_k$.

\subsection*{Checking $bb\dots b$ inputs}

Algorithm~\ref{alg:oneletterSets} was proposed in \cite{vu}. It checks whether Phase~2 stops when it processes input digits $bb\dots b$.
Sets $\Qb{m}$  can be easily constructed separately for each $b\in\B$ for given $m$. We build the set $\Qb{m}$ for input digits $bb\dots b$ in the same way as in Phase 2. This means that we first search for $\Qb{1}$ such that 
$$
b + \Q \subset \A + \beta \Qb{1}\,.
$$
Until the set $\Qb{m}$ contains only one element, we increment the length of  window $m$ and, using Algorithm \ref{alg:possibleWeightCoefSetStable}, we find a subset $\Qb{m+1}$ of the set $\Qb{m}$ such that
$$
b + \Qb{m} \subset \A + \beta \Qb{m+1}\,.
$$
In each iteration, we check whether the set $\Qb{m+1}$ is strictly smaller than the set $\Qb{m}$. If not, we know by Theorem \ref{thm:bbbCondition} that Phase 2 does not converge because $\#\Qb{m}$ is eventually a constant greater than 2.

Hence, non-finiteness of Phase~2 can be revealed by running Algorithm \ref{alg:oneletterSets} for each input digit $b\in\B$.
\begin{algorithm}
  \caption{Check the input $bb\dots b$}
    \label{alg:oneletterSets}
  \begin{algorithmic}[1]
    \REQUIRE{Weight coefficient set $\Q$, digit $b\in\B$}
    \ENSURE{TRUE if there is a unique weight coefficient for input $bb\dots b$, FALSE otherwise}
    \STATE Find minimal set $\Qb{1} \subset \Q$ such that
      $$
      b + \Q \subset \A + \beta \Qb{1}\,.
      $$
      \vspace{-20pt}
    \STATE $m:=1$
    \WHILE{$\#\Qb{m} > 1$}
        \STATE $m:= m +1$
        \STATE By Algorithm \ref{alg:possibleWeightCoefSetStable}, find minimal set $\Qb{m} \subset \Qb{m-1}$ such that
          $$
          b + \Qb{m-1} \subset \A + \beta \Qb{m}\,.
          $$  
          \vspace{-20pt}
        \IF{$\#\Qb{m}=\#\Qb{m-1}$}
            \RETURN FALSE
        \ENDIF
    \ENDWHILE  
    \RETURN TRUE
  \end{algorithmic}
\end{algorithm}


\subsection*{Stable Phase 2}
Recall that in order to use Theorem~\ref{thm:infinitePathInRauzyGraph} to check non-convergence of Phase~2, we require Phase~2 to be stable. An algorithm of choice of possible weight coefficients set can be easily modified to ensure stability of Phase~2, i.e., 
$$\Qw{1}{k}=\Qw{1}{(k-1)} \implies \Qwo{k}=\Qwo{(k-1)}$$
for all $\tuple{1}{k}\in\B^k$. If the algorithm is deterministic, then the set $\Qwo{k}$ is determined by $\Qw{1}{k}$, $\Qwo{(k-1)}$ and $w_0$. Similarly, the set $\Qwo{(k-1)}$ is determined by  the set $\Qw{1}{(k-1)}$, $\Qwo{(k-2)}$ and $w_0$. 

Suppose that $\Qw{1}{k}=\Qw{1}{(k-1)}$. Now, the only difference between $\Qwo{k}$ and $\Qwo{(k-1)}$ is that $\Qwo{k}$ is searched as a subset of $\Qwo{(k-1)}$, whereas $\Qwo{(k-1)}$ as a subset of $\Qwo{(k-2)}$. In order to find  $\Qwo{(k-1)}$, we use Algorithm~\ref{alg:minimalSet} repeatedly instead of once. In each iteration, the input digit $w_0$ and the set $\Qw{1}{(k-1)}$ remains the same but we search for $\Qwo{(k-1)}$ as a subset of $\Q'_w$, where $\Q'_w$ is the output of the previous iteration or $\Qwo{(k-2)}$ in the first iteration. The algorithm stops when $\Qwo{(k-1)}=\Q'_w$. This loop is described in Algorithm~\ref{alg:possibleWeightCoefSetStable}.

Now, when the set $\Qwo{k}$ is searched as a subset of $\Qwo{(k-1)}$ in the same manner, it runs with the same inputs as the last iteration of the search for  $\Qwo{(k-1)}$. Hence, $\Qwo{k}=\Qwo{(k-1)}$.


\begin{algorithm}
  \caption{Stable search for possible weight coefficient set $\Qwo{k}$}
    \label{alg:possibleWeightCoefSetStable}
  \begin{algorithmic}[1]
    \REQUIRE{Input digit $w_0$, set of possible carries $\Qw{1}{k}$, previous set of possible weight coefficients $\Qwo{(k-1)}$}
	\STATE $\Q'_w:=\Qwo{(k-1)}$
	\WHILE{TRUE}
		\STATE By Algorithm \ref{alg:minimalSet}, find the set of possible weight coefficients  $\Qwo{k}$ as a subset of $\Q'_w$ instead of the previous set of weight coefficients $\Qwo{(k-1)}$.
		\IF{$\Qwo{k}=\Q'_w$}
			\RETURN $\Qwo{k}$
		\ENDIF
		\STATE $\Q'_w:=\Qwo{k}$
	\ENDWHILE
  \end{algorithmic}
\end{algorithm}


\subsection*{Infinite walk in Rauzy graph $G_k$}
As the Rauzy graph $G_k$ is finite, there exist an infinite walk in $G_k$ if and only if there exists an oriented cycle. Algorithm~\ref{alg:checkCycles}  checks whether some walk starting in $\tuple{1}{k}$ enters such a cycle eventually. First, it  checks if the vertex $\tuple{1}{k}$ is in the graph $G_k$. If yes, all vertices which are accessible by appropriately directed edge from $\tuple{1}{k}$ are entered by Algorithm~\ref{alg:enterVertices}. It is called recursively to enter all accessible vertices from the given one. The already traversed path is passed in each call and if an vertex is visited second time, then the cycle is found and algorithm ends. If all branches of the recursion ends up without visiting some vertex twice, then there is no cycle and thus no infinite path.

Note that saving of the traversed path requires only the label of the first vertex and last digits of the labels of next visited vertices due to the construction of Rauzy graph.
\begin{algorithm}
  \caption{Check if there is in an infinite walk in $G_k$ starting in $\tuple{1}{k}$}
    \label{alg:checkCycles}
  \begin{algorithmic}[1]
    \REQUIRE{Rauzy graph $G_k$, combination of input digits $\tuple{1}{k}$}
    \ENSURE{Return TRUE if TRUE is returned in any step of the recursion, that is when a walk $w_1, w_2, \dots$ enters a directed cycle in $G_k$. Otherwise return FALSE.}
	\IF {$\tuple{1}{k}\in G_k$}
		\STATE By Algorithm~\ref{alg:enterVertices}, enter next vertices from $\tuple{1}{k}$ with the traversed path  $\tuple{1}{k}$.
	\ELSE
		\RETURN FALSE
	\ENDIF
	\RETURN FALSE
  \end{algorithmic}
\end{algorithm}

\begin{algorithm}
  \caption{Enter vertices from $\tuple[w']{1}{k}$}
    \label{alg:enterVertices}
  \begin{algorithmic}[1]
  	\REQUIRE{Rauzy graph $G_k$, vertex $\tuple[w']{1}{k}$, traversed path $(w_1, \dots, w_l)$.}
	\FORALL {$w'_{k+1}\in \B$ such that $\tuple[w']{1}{k}\rightarrow \tuple[w']{2}{(k+1)}\in G_k$}
		\IF {$\tuple[w']{2}{(k+1)}$ is in the traversed path $(w_1, \dots, w_l)$}
			\RETURN TRUE
		\ELSE
			\STATE By Algorithm~\ref{alg:enterVertices}, enter next vertices from $\tuple[w']{2}{(k+1)}$ with the traversed path   $(w_1, \dots, w_l, w'_{(k+1)})$.
		\ENDIF
	\ENDFOR
  \end{algorithmic}
\end{algorithm}


\subsection*{Phase 2 with control of non-convergence}
Algorithm~\ref{alg:weightFunction_modified} modifies the basic proposal of Phase~2 to reveal its possible non-convergence. First, the necessary condition given by Theorem~\ref{thm:bbbCondition} is checked, i.e., the convergence of Phase~2 for inputs $bb\dots b$ is verified by Algorithm~\ref{alg:oneletterSets} for all $b\in\B$. 

\begin{algorithm}
  \caption{Modified search for a weight function (Phase 2)}
    \label{alg:weightFunction_modified}
  \begin{algorithmic}[1]
    \REQUIRE{weight coefficients set $\Q$}
    \FORALL{$b\in\B$}
    	\IF  {\NOT Check the input $bb\dots b$ by Algorithm \ref{alg:oneletterSets}}
    		\RETURN Phase 2 does not converge for input $bb\dots b$.
    	\ENDIF
    \ENDFOR
    \FORALL{$w_0 \in \B$} 
        \STATE By Algorithm~\ref{alg:possibleWeightCoefSetStable}, find set $\Q_{[w_0]} \subset \Q$ such that
          $$
          w_0 + \Q \subset \A + \beta \Q_{[w_0]}
          $$\vspace{-20pt}
    \ENDFOR
    \STATE $k:=0$
    %\STATE $G_0:=$ empty graph
    \WHILE{$\max\{\#\Qwo{k}:\tupleo{k} \in \B^{k+1} \} > 1$}
        \STATE $k:= k +1$
        \FORALL{$\{\tupleo{k} \in \B^{k+1}\colon \#\Qwo{(k-1)}>1\}$}
            \STATE By Algorithm~\ref{alg:possibleWeightCoefSetStable}, find set $\Qwo{k} \subset \Qwo{(k-1)}$ such that
              $$
              w_0 + \Qw{1}{k} \subset \A + \beta \Qwo{k}\,.
              $$\vspace{-20pt}
              \IF {$\Qwo{k} = \Qwo{(k-1)}$ \AND $k\geq 2$}
              	\STATE Add the vertex $\tupleo{k}$ to the Rauzy graph $G_{k+1}$.
              	\IF{Check infinite walks in $G_k$ starting in $\tuple{1}{k}$ by Alg.~\ref{alg:checkCycles}}
              		\RETURN Phase 2 does not converge.
              	\ENDIF
              \ENDIF
     %    	\STATE $G_{k}:=G_{k+1}$
        \ENDFOR  
    \ENDWHILE  
    \STATE $r:= k+1$
    \FORALL {$l\in\NN, l\leq r$}
	    \FORALL{$\{\tupleo{l}\in\B^{l+1} \colon \Qwo{l} \text{ was found, } \#\Qwo{l}=1\}$}
	    	\FORALL{$\tuple{(l+1)}{r}\in\B^{r-l}$}
	    		\STATE $q\tupleo{r}:=$ only element of $\Qwo{l}$
	    	\ENDFOR
	    \ENDFOR
	\ENDFOR
    \RETURN $q$
  \end{algorithmic}
\end{algorithm}

Then we proceed in the same way as in Algorithm~\ref{alg:weightFunction} with the following modifications: it is sufficient to process only $\tupleo{k} \in \B^{k+1}$ such that $\#\Qwo{(k-1)}>1$ since if $\#\Qwo{(k-1)}=1$, then $\Qwo{k'}=\Qwo{(k-1)}$ for all $k'>k$. 

Moreover, we check possible non-convergence according to Theorem~\ref{thm:infinitePathInRauzyGraph}. If $\Qwo{k}$ equals $\Qwo{(k-1)}$, then the vertex $\Qwo{k}$ is added into the Rauzy graph $G_{k+1}$ and the Rauzy graph $G_k$ is examined by Algorithm~\ref{alg:checkCycles} whether it contains an infinite walk starting in $\tuple{1}{k}$. Note that $\Qwo{k} = \Qwo{(k-1)}$ is a necessary condition for $\tuple{1}{k}$ be a vertex of $G_k$. It also implies that $\#\Qwo{k}>1$.

We remark that  non-convergence caused by an input $bb\dots b$ for some $b\in\B$ would be revealed also as a cycle in a Rauzy graph. Nevertheless, the number of calls of Algorithm~\ref{alg:minimalSet} during Algorithm~\ref{alg:oneletterSets} is at most $\#Q$ for each $b\in\B$, while it grows exponentially with the length of window in search for a weight function. Thus, we save a lot of computational time in cases which fail already on $bb\dots b$ inputs. At the same time, the cost of this check is low in other cases.








	\section{Implementation}

Our implementation of the design is based on the program attached to \cite{vu}. The chosen programming language is SageMath. It is Python-based language with numerous implemented mathematical structures. That is the main reason of our choice -- the extending window method requires to handle elements of $\Zomega$ and arithmetic operations in this set. Moreover, SageMath provides various data structures and plotting tools.  The code is simpler and more similar to pseudocode than if we implemented in pure \Cpp. Due to it, we may concern ourselves with the algorithmic part of the problem instead of difficult programming. Unfortunately, SageMath is much slower than \Cpp.


The implementation is object-oriented. It consists of five classes. Class \emph{AlgorithmForParallelAddition} contains structures for computations in $\Zomega$. The build-in classes \emph{PolynomialQuotientRing} and \emph{NumberField} are used to represent elements of $\Zomega$ as an algebraic and complex numbers. The class also links together all functions and instances of other classes which are necessary to construct an algorithm for digit set conversion from $\B$ to $\A$ by the extending window method. The input parameters are an algebraic integer $\omega$ given by its minimal polynomial $m_\omega$ and an approximate complex value, a base $\beta\in\Zomega$, an alphabet $\A\subset\Zomega$ and input alphabet $\B\subset\Zomega$. 

Phase 1 of the extending window method is implemented in class \emph{WeightCoefficientsSetSearch} and Phase 2 in \emph{WeightFunctionSearch}. Class \emph{WeightFunction} serves to save a function $q$ computed in Phase 2. The last class \emph{ExceptionParAdd} is inherited from build-in class \emph{Exception} to distinguish between errors which are raised by the algorithm  of the extending window method and other ones.

We use notation from previous chapters in descriptions of the classes. We list only the most important methods of each class, see commented source code for all of them.  

For basic use, load \verb+AlgorithmForParallelAddition.sage+, create an instance of \emph{AlgorithmForParallelAddition} and call \textbf{findWeightFunction()}.

We also provide an interface -- a shell script with given parameters is executed. If the network access is enabled and the modul \verb+gspread+ is installed (see \cite{gspread}), then results of computation are automatically saved to Google spreadsheet \href{https://docs.google.com/spreadsheets/d/1TnhrHdefHfHa0WSeVs4q6XVj3epjPlPlnoekE0E1xeM/edit?usp=sharing}{ParallelAddition\_results}. The spreadsheet can be also used for loading input parameters. Both interfaces are described at the end of this section. The whole implementation is on the attached DVD or it can be downloaded from  \url{https://github.com/Legersky/ParallelAddition}.



\subsection*{Class AlgorithmForParallelAddition}
This class constructs necessary structures for computation in $\Zomega$. It is \emph{PolynomialQuotientRing} obtained as a \emph{PolynomialRing} over integers factored by polynomial $p$. This is used for representation of elements of $\Zomega$ and arithmetic. We remark that it is independent on the choice of root of the  minimal polynomial $p$. But as we need also comparisons of numbers in~$\Zomega$ in modulus, we specify $\omega$ by its approximate complex value and we form a factor ring of rational polynomials by using class \emph{NumberField}. This enables us to get absolute values of elements of $\Zomega$ which can be then compared.

Method \textbf{findWeightFunction()} links together both phases of the extending window method to find the weight function $q$. That is used in the methods for addition and digit set conversion to process them as local functions. There are also verification methods.

Moreover, many methods for printing, plotting and saving outputs are implemented.

The constructor of class \emph{AlgorithmForParallelAddition} is 

\begin{method}{\_\_init\_\_}{minPol\_str, embd, alphabet, base, name='NumerationSystem', inputAlphabet=' ',\\
 printLog=True, printLogLatex=False, verbose=0}
Take \var{minPol\_str} which is a string of symbolic expression in the variable $x$ of an irreducible polynomial~$p$. The closest root of  \var{minPol\_str} to \var{embd} is used as the ring generator $\omega$ (see more in documentation of \emph{NumberField} in SageMath \cite{sage}). The structures for $\Zomega$ are constructed as described above. Setters \fun{setAlphabet}{alphabet}, \fun{setInputAlphabet}{inputAlphabet} and \fun{setBase}{base} are called.  Messages saved to logfile during existence of an instance are printed (using \LaTeX) on standard output depending on \var{printLog} and \var{printLogLatex}. The level of messages for a development is set by \var{verbose}. 
\end{method}

Methods for the construction of an addition algorithm which is computable in parallel by the designed extending window method are the following:

\begin{method}{\_findWeightCoefSet}{ max\_iterations, method\_number}
Create an instance of \emph{WeightCoefficientsSetSearch(method\_number)} and call its method \fun{findWeightCoefficientsSet}{max\_iterations} to obtain a weight coefficients set $\Q$.
\end{method}

\begin{method}{\_findWeightFunction}{ max\_input\_length,method\_number}
Create an instance of \emph{WeightFunctionSearch(method\_number)} and call its methods \fun{check\_one\_letter\_inputs}{max\_input\_length} and \\ \fun{findWeightFunction}{max\_input\_length} to obtain a weight function $q$.
\end{method}


\begin{method}{findWeightFunction}{ max\_iterations, max\_input\_length, method\_weightCoefSet=None,\\ method\_weightFunSearch=None}
It is verified that there are all representatives mod $\beta$ in the alphabet and that all elements of the input alphabet have their representatives mod $\beta-1$ in the alphabet, see Section \ref{sec:alphabet}.
Return the weight function $q$ obtained by calling \fun{\_findWeightCoefSet\\}{max\_iterations, method\_weightCoefSet} and \fun{\_findWeightFunction}{max\_input\_length, method\_weightFunSearch}.
\end{method}

The important function for the searching for possible weight coefficients is

\begin{method}{divideByBase}{divided\_number}
Using Theorem \ref{thm:divisibility}, check if \var{divided\_number} is divisible by the base $\beta$. If it is so, return the result of division, else return \var{None}.
\end{method}


Methods for the addition and the digit set conversion computable in parallel are following:

\begin{method}{addParallel}{a,b}
Sum up numbers represented by the lists of digits \var{a} and \var{b} digitwise and convert the result by \fun{parallelConversion}{}. 
\end{method}


\begin{method}{parallelConversion}{\_w}
Return $(\beta,\A)$-representation of the number represented by the list \var{\_w} of digits in the input alphabet $\B$. According to the equation \eqref{eq:conversionFormula}, it is computed locally by using the weight function $q$.
\end{method}


\begin{method}{localConversion}{w}
Return converted digit according to the equation \eqref{eq:conversionFormula} for the list of input digits \var{w}.
\end{method}


The correctness of the implementation of the extending window for a given numeration system can be verified by
 
\begin{method}{sanityCheck\_conversion}{ num\_digits}
Check whether the values of all possible numbers of the length \var{num\_digits} with digits in the input alphabet $\B$ are the same as their $(\beta, \A)$-representation obtained by \fun{parallelConversion}{}.   
\end{method}




\subsection*{Class WeightCoefficientsSetSearch}
Class \emph{WeightCoefficientsSetSearch} implements Phase 1 of the extending window method described in Section \ref{subsec:phase1} with different methods how an intermediate weight coefficients set $\Q_k$ is extended to $\Q_{k+1}$. Algorithm~\ref{alg:extendWeightCoefSet} explains some of these methods. The most important method of the class is  \fun{findWeightCoefficientsSet}{} which returns a weight coefficients set $\Q$.

The constructor of the class is 
\begin{upravit}
\begin{method}{\_\_init\_\_}{ algForParallelAdd, method}
Initialize a ring generator $\omega$, base $\beta$, an alphabet $\A$ and input alphabet $\B$ by values obtained from \var{algForParallelAdd}. The parameter \var{method} characterizes how an intermediate weight coefficients set $\Q_k$ is extended to $\Q_{k+1}$. If \var{None}, then method 13 from Algorithm~\ref{alg:extendWeightCoefSet} is used as default.
\end{method}

Methods implementing Phase 1 are the following:

\begin{method}{\_findCandidates}{C}
Following Algorithm \ref{alg:searchCand}, return the list of lists \var{candidates} such that each element in \var{C} is covered by any value of the appropriate list in \var{candidates}.  
\end{method}


\begin{method}{\_chooseQk\_FromCandidates}{candidates}
Take the previous intermediate weight coefficients set $\Q_{k}$ as the class attribute and choose from \var{candidates} the intermediate weight coefficients set $\Q_{k+1}$ by Algorithm \ref{alg:extendWeightCoefSet}.
\end{method}


\begin{method}{\_getQk}{C}
Links together methods \fun{\_findCandidates}{C} and \fun{\_chooseQk\_FromCandidates}{} to return itermediate weight coefficients set $\Q_k$.
\end{method}


\begin{method}{findWeightCoefficientsSet}{ maxIterations}
According to Algorithm \ref{alg:weightCoefSet}, return the weight coefficients set $\Q$ which is build iteratively by using \fun{\_getQk}{}. The computation is aborted if the number of iterations exceeds \var{maxIterations}. 
\end{method}




\subsection*{Class WeightFunctionSearch}

This class implements modified Phase 2 of the extending window method from Section \ref{sec:modifiedPhase2} with different methods of choice of possible weight coefficients sets and control of non-convergence. Methods 2a, 2b, 2c, 2d and 2e from Algorithm~\ref{alg:pickElement} correspond to 9, 15, 22, 23 and 14 respectively. A weight function $q$ is returned by method \fun{findWeightFunction}{}. The constructor of the class is

\begin{method}{\_\_init\_\_}{algForParallelAdd, weightCoefSet, method, maxInputs}
The ring generator $\omega$, base $\beta$, alphabet $\A$ and input alphabet $\B$ are initialized by the values obtained from \var{algForParallelAdd}. The weight coefficients set $\Q$ is set to \var{weightCoefSet}. The parameter \var{method} (the number of an experimental method or \verb+'2a'+, \verb+'2b'+, etc.) determines the way of the choice of a possible weight coefficients set $\Qwo{k}$ from $\Qwo{(k-1)}$, see Algorithm~\ref{alg:pickElement}. If \var{method} is \var{None}, then is the default method is 2b. Possible weight coefficients sets  are saved in a dictionary \var{\_Qw\_w}, which is set to be empty.
\end{method}

The following methods are used for search for a weight function $q$:

\begin{method}{\_find\_weightCoef\_for\_comb\_B}{combinations}
Take all combinations of input digits $\tupleo{(k-1)} \in \B^{k}$ in \var{combinations} such that $\#\Qwo{(k-1)}>1$, extend them by all letters $w_{-k}\in\B$ and find a possible weight coefficients set $\Qwo{k}$ by the method \fun{\_findQw}{$\tupleo{(k-1)}$}. If there is only one element in $\Qwo{k}$, it is saved as a solved input of the weight function $q$. Otherwise, the set $\Qwo{k}$ is saved in \var{\_Qw\_w} as an unsolved combination which requires extending of the window. The unsolved combinations are returned.  
\end{method}

\begin{method}{\_findQw}{w\_tuple}
Return a set $\Qwo{k}=\Q_{[\text{w\_tuple}]}$ by wrapping \fun{\_findQw\_once}{} into a while loop according to Algorithm~\ref{alg:possibleWeightCoefSetStable}. If $\Qwo{k}=\Qwo{(k-1)}$, then add a vertex $\tupleo{k}$ to a Rauzy graph $G_{k+1}$ and call \fun{\_checkCycles}{w\_tuple} in $G_{k}$.
\end{method}

\begin{method}{\_findQw\_once}{w\_tuple,Qw\_prev} 
Return a set $\Qwo{k}'=\Q'_{[\text{w\_tuple}]}$ obtained by Algorithm~\ref{alg:minimalSet} as a subset of \var{Qw\_prev}$=\Qwo{(k-1)}$. The set of possible carries from the right $\Qw{1}{k}$ is taken from the class attribute \var{\_Qw\_w}. The methods of Algorithm~\ref{alg:pickElement} are implemented here along with the experimental ones.
\end{method}

\begin{method}{\_checkCycles}{ w\_tuple}
\end{method}

\begin{upravit}
\begin{method}{findWeightFunction}{max\_input\_length}
Return the weight function $q$ unless the length of window exceeds \var{max\_input\_length}. Then an exception is raised. It implements Algorithm \ref{alg:weightFunction} by repetetive calling of the method \fun{\_find\_weightCoef\_for\_comb\_B}{} which extends length of window. This is done until all possible combinations of input digits are solved for some length of window $m$, i.e. $\max\{\#\Q_{[w_j,\dots, w_{j-m+1}]}:(w_j,\dots, w_{j-m+1}) \in \B^m \} = 1$.
\end{method}


\begin{method}{check\_one\_letter\_inputs}{max\_input\_length}
The method checks by Algorithm \ref{alg:oneletterSets} if there is a unique weight coefficient for inputs $({b,b,\dots,b})\in\B^m$ for some length of window $m$. Using Theorem \ref{thm:suffCondPhase1}, an exception is raised in the case of an infinite loop. Otherwise the list of inputs $(b,b,\dots,b)$ which have the largest length of the window is returned.
\end{method}

%\begin{method}{\_checkCycles}{ w\_tuple}
%
%
%
%
%\begin{method}{isSublist}{    def isSublist(\_list, \_sublist}
%
%\end{method}
%
%
%\begin{method}{find\_next\_letter}{    def find\_next\_letter(\_w,witness\_seq}
%
%\end{method}
%
%
%
%
%\begin{method}{\_findQw\_once}{w\_tuple,Qw\_prev}
%
%\end{method}
%
%
%\begin{method}{findWeightFunction}{ max\_input\_length}
%
%\end{method}
%
%
%\begin{method}{check\_one\_letter\_inputs}{ max\_input\_length}
%
%\end{method}






\subsection*{Class WeightFunction}
This class serves for saving the weight function $q$. The constructor is

\begin{method}{\_\_init\_\_}{B}
Set the input alphabet to \var{B} and the maximum length of window to 1. Initialize the attribute \var{\_mapping} to an empty dictionary for saving the weight function $q$. 
\end{method}

The methods for saving and calling are following:

\begin{method}{addWeightCoefToInput}{\_input, coef}
Save the weight coefficient \var{coef} for \var{\_input} to the class attribute \var{\_mapping}. The digits of \var{\_input} must be in the input alphabet.
\end{method}

\begin{method}{getWeightCoef}{w}
The digits of the list \var{w} are taken from the left until the weight coefficient in the dictionary \var{\_mapping} is found. 
\end{method}

The result of the method \fun{getWeightCoef}{} is used to make this class callable, i.e., if \var{\_q} is an instance of \emph{WeightFunction}, then \var{\_q}.\fun{getWeightCoef}{w} is the same as \var{\_q}(\var{w}).







\subsection*{User interfaces}
We provide two options of loading inputs for running the implemented extending window method. SageMath must be installed as they are executed as shell scripts.

The implementation of the extending window method is launched in a shell by typing \verb+sage ewm_inputs.sage+. The parameters of a numeration system and setting of outputs and computation is given in the head of the file \verb+ewm_inputs.sage+. See Appendix \ref{app:interfaces} for an example of such a file.


The name of the numeration system, minimal polynomial of generator $\omega$, an approximate value of $\omega$, the base $\beta$, alphabet $\A$ and input alphabet $\B$ are set in the part INPUTS. Different methods of choice for Phase 1 and 2 can be set. If there are more methods in the lists, then methods for Phase~1 are compared first. Next, each distinct result is processed with each method for Phase~2.  

The sanity check is launched according to the boolean value in the part SANITY CHECK. 

The boolean values in the part SAVING determines which formats of the outputs are saved. All outputs are saved in the folder \verb+./outputs/<name>/+. General information about the computation can be saved in the \LaTeX{} format, a computed weight function and local digit set conversion in the CSV file format.  A log of the whole computation can be saved as a text file.

Figures of the alphabet, input alphabet or weight coefficients set are saved in the PNG format in the folder \verb+./outputs/<name>/img/+ according to the boolean values in the part IMAGES. Images of individual steps of both phases of the extending window method can be also saved. For Phase 2, searching for a weight coefficient  is plotted for given input digits.  

The program prints out all inputs and then it computes a weight function $q$ by calling \fun{findWeightFunction}{}. Intermediate weight coefficients sets in each iteration of Phase 1 and the obtained weight coefficients set $\Q$ are printed out. Convergence of Phase~2 for combinations given by repetition $b\in\B$ is verified by \fun{check\_one\_letter\_inputs}{}. The processed length of window is showed during computing of Phase 2. At the end, the final length of window, elapsed time and info about saved outputs are printed. Results are also saved in the Google spreadsheet \href{https://docs.google.com/spreadsheets/d/1TnhrHdefHfHa0WSeVs4q6XVj3epjPlPlnoekE0E1xeM/edit?usp=sharing}{ParallelAddition\_results} in the worksheet \verb+results+ and \verb+comparePhase1+ if there are more methods for Phase~1.

The second option is to execute the script \verb+ewm_gspreadsheet.sage+ (see Appendix \ref{app:interfaces}). Parameters are loaded from the worksheet \verb+inputs+ in the Google spreadsheet \href{https://docs.google.com/spreadsheets/d/1TnhrHdefHfHa0WSeVs4q6XVj3epjPlPlnoekE0E1xeM/edit?usp=sharing}{ParallelAddition\_results}. The column A marks whether a row should be tested. The columns B--G, i.e., \emph{Name}, \emph{Alphabet}, \emph{Input alphabet}, \emph{Approximate value of ring generator omega}, \emph{Minimal polynomial omega} and \emph{Base} must be filled. If the column \emph{Input alphabet} is empty, then the input alphabet $\A+\A$ is used. Methods for Phase~1, resp. 2, are given in the cell C1, resp. C2.

Program runs in the same way as before, but results are saved only to the Google spreadsheet. Notice that column A and cells with the methods may be changed  after executing the script, but do not modify other cells or order of rows. The reason is that the program reads the methods at the beginning and it remembers position of rows to be tested, but the parameters are loaded on the fly.



\chapter{Testing and results}
\label{chap:testing}
	\komentar{tabulka vsech uspesnych plus nejake neuspesne a okomentovat}

%\afterpage{
%\clearpage   % To flush out all floats, might not be what you want
%\begin{landscape}
%	\begin{table}[h]
%		\begin{center}
%		\begin{tabular}{l|c cc| c c c} 
 $\omega$                & Base                                                    & $\#\A$ & $m_\beta$ & Phase 1 & bbb & Phase 2 \\ \hline
$1/2\cdot I\sqrt{15} - 1/2 $ & $  \omega $ & $  6    $ & $  x^2 + x + 4     $  & \checkmark      & \checkmark  & \xmark  \\ \hline
$1/2\cdot I\sqrt{11} + 1/2 $ & $  -2\omega - 2  $ & $  27   $ & $  x^2 + 6x + 20    $  & \checkmark      & \checkmark  & \checkmark      \\
$1/2\cdot I\sqrt{11} + 1/2 $ & $  -\omega+ 1  $ & $  3    $ & $  x^2 - x + 3  $  & \checkmark      & \checkmark  & \xmark  \\ \hline
$1/2\cdot I\sqrt{7} - 1/2  $ & $  \omega  $ & $  4    $ & $  x^2 + x + 2    $  & \checkmark      & \checkmark  & \checkmark      \\
$1/2\cdot I\sqrt{7} + 1/2  $ & $  \omega    $ & $  2    $ & $  x^2 - x + 2      $  & \checkmark      & \checkmark  & \xmark  \\
$1/2\cdot I\sqrt{7} + 3/2  $ & $  \omega                                                   $ & $  4    $ & $  x^2 - 3x + 4                                              $& \checkmark  & \xmark  & -     \\ \hline
%$1/2\cdot I\sqrt{3} - 1/2  $ & $  2\omega                                                 $ & $  8    $ & $  x^2 + 2x + 4                                          $  & \checkmark      & \checkmark  & \checkmark      \\
$I\sqrt{3}            $ & $  -\omega                                                  $ & $  4    $ & $  x^2 + 3                                                $  & \checkmark      & \checkmark  & \checkmark      \\
$1/2\cdot I\sqrt{3} + 1/2  $ & $  -3\omega + 2                                            $ & $  7    $ & $  x^2 - x + 7                                            $  & \checkmark      & \checkmark  & \xmark  \\ \hline
$I\sqrt{2}            $ & $  -\omega                                                  $ & $  3    $ & $  x^2 + 2                                                $  & \checkmark      & \checkmark  & \checkmark      \\
%$I\sqrt{2} + 1        $ & $  -\omega  - 3  $ & $  27   $ & $  x^2 + 8x + 18                                         $  & \checkmark      & \checkmark  & \checkmark      \\
$I\sqrt{2} - 1        $ & $  \omega                                                   $ & $  6    $ & $  x^2 + 2x + 3                                          $  & \checkmark      & \checkmark  & \xmark  \\ \hline
$I - 1                $ & $  \omega                                                   $ & $  5    $ & $  x^2 + 2x + 2                                          $  & \checkmark      & \checkmark  & \checkmark      \\
$I                    $ & $  -3\omega                                                $ & $  10   $ & $  x^2 + 9                                                $  & \checkmark      & \checkmark  & \xmark  \\ \hline
$-1/2\sqrt{5} + 3/2   $ & $  -2\omega - 2                                            $ & $  31   $ & $  x^2 + 10x + 20                                        $  & \checkmark      & \checkmark  & \checkmark      \\
$-1/2\sqrt{5} + 3/2   $ & $  3\omega   - 3 $ & $  13   $ & $  x^2 - 3x - 9                                          $  & \checkmark      & \checkmark  & \xmark  \\ \hline
$1/2\sqrt{13} - 3/2   $ & $  \omega   - 3   $ & $  27   $ & $  x^2 + 9x + 17                                         $  & \checkmark      & \checkmark  & \xmark  \\
$1/2\sqrt{13} + 1/2   $ & $  -2\omega + 2                                            $ & $  16   $ & $  x^2 - 2x - 12                                         $  & \checkmark      & \xmark  & -  \\ \hline
$1/2\sqrt{17} - 3/2   $ & $  \omega   - 3   $ & $  26   $ & $  x^2 + 9x + 16                                         $  & \checkmark      & \checkmark  & \checkmark           
\end{tabular}
%		\end{center}
%	\caption{xxx}
%	\label{tab:results}
%	\end{table}
%\end{landscape}
%}


\begin{table}[h]
	\begin{center}
	\begin{tabular}{l|c cc| c c c} 
 $\omega$                & Base                                                    & $\#\A$ & $m_\beta$ & Phase 1 & bbb & Phase 2 \\ \hline
$1/2\cdot I\sqrt{15} - 1/2 $ & $  \omega $ & $  6    $ & $  x^2 + x + 4     $  & \checkmark      & \checkmark  & \xmark  \\ \hline
$1/2\cdot I\sqrt{11} + 1/2 $ & $  -2\omega - 2  $ & $  27   $ & $  x^2 + 6x + 20    $  & \checkmark      & \checkmark  & \checkmark      \\
$1/2\cdot I\sqrt{11} + 1/2 $ & $  -\omega+ 1  $ & $  3    $ & $  x^2 - x + 3  $  & \checkmark      & \checkmark  & \xmark  \\ \hline
$1/2\cdot I\sqrt{7} - 1/2  $ & $  \omega  $ & $  4    $ & $  x^2 + x + 2    $  & \checkmark      & \checkmark  & \checkmark      \\
$1/2\cdot I\sqrt{7} + 1/2  $ & $  \omega    $ & $  2    $ & $  x^2 - x + 2      $  & \checkmark      & \checkmark  & \xmark  \\
$1/2\cdot I\sqrt{7} + 3/2  $ & $  \omega                                                   $ & $  4    $ & $  x^2 - 3x + 4                                              $& \checkmark  & \xmark  & -     \\ \hline
%$1/2\cdot I\sqrt{3} - 1/2  $ & $  2\omega                                                 $ & $  8    $ & $  x^2 + 2x + 4                                          $  & \checkmark      & \checkmark  & \checkmark      \\
$I\sqrt{3}            $ & $  -\omega                                                  $ & $  4    $ & $  x^2 + 3                                                $  & \checkmark      & \checkmark  & \checkmark      \\
$1/2\cdot I\sqrt{3} + 1/2  $ & $  -3\omega + 2                                            $ & $  7    $ & $  x^2 - x + 7                                            $  & \checkmark      & \checkmark  & \xmark  \\ \hline
$I\sqrt{2}            $ & $  -\omega                                                  $ & $  3    $ & $  x^2 + 2                                                $  & \checkmark      & \checkmark  & \checkmark      \\
%$I\sqrt{2} + 1        $ & $  -\omega  - 3  $ & $  27   $ & $  x^2 + 8x + 18                                         $  & \checkmark      & \checkmark  & \checkmark      \\
$I\sqrt{2} - 1        $ & $  \omega                                                   $ & $  6    $ & $  x^2 + 2x + 3                                          $  & \checkmark      & \checkmark  & \xmark  \\ \hline
$I - 1                $ & $  \omega                                                   $ & $  5    $ & $  x^2 + 2x + 2                                          $  & \checkmark      & \checkmark  & \checkmark      \\
$I                    $ & $  -3\omega                                                $ & $  10   $ & $  x^2 + 9                                                $  & \checkmark      & \checkmark  & \xmark  \\ \hline
$-1/2\sqrt{5} + 3/2   $ & $  -2\omega - 2                                            $ & $  31   $ & $  x^2 + 10x + 20                                        $  & \checkmark      & \checkmark  & \checkmark      \\
$-1/2\sqrt{5} + 3/2   $ & $  3\omega   - 3 $ & $  13   $ & $  x^2 - 3x - 9                                          $  & \checkmark      & \checkmark  & \xmark  \\ \hline
$1/2\sqrt{13} - 3/2   $ & $  \omega   - 3   $ & $  27   $ & $  x^2 + 9x + 17                                         $  & \checkmark      & \checkmark  & \xmark  \\
$1/2\sqrt{13} + 1/2   $ & $  -2\omega + 2                                            $ & $  16   $ & $  x^2 - 2x - 12                                         $  & \checkmark      & \xmark  & -  \\ \hline
$1/2\sqrt{17} - 3/2   $ & $  \omega   - 3   $ & $  26   $ & $  x^2 + 9x + 16                                         $  & \checkmark      & \checkmark  & \checkmark           
\end{tabular}
	\end{center}
\caption{xxx}
\label{tab:results}
\end{table}



\cleardoublepage\phantomsection
\addcontentsline{toc}{chapter}{Summary}
\chapter*{Summary}
\markboth{\MakeUppercase{Summary}}{}
The main goal of this thesis was to improve the extending window method with rewriting rule $x-\beta$ which attempts to construct a parallel addition algorithm for a given base $\beta$ and alphabet $\A$ and discuss its convergence.

We  recalled necessary definitions and notation in the scope of numeration systems and parallel addition. The concept of parallel addition and extending  window method was explained. We proposed various methods for both phase of the method. Some of them involve $\beta$-norm, which was constructed by using the companion matrix of the minimal polynomial $m_\beta$.

We solved convergence of Phase~1 -- a sufficient condition is that the base $\beta$ is expanding, i.e., all its conjugates are greater than one in modulus. We showed that it is also a necessary condition of convergence of the whole extending method. 

The check of convergence of Phase~2 for $bb\dots b$ inputs was reviewed. Next, we introduced the notion of stable Phase~2 and a Rauzy graph. We proved Theorem~\ref{thm:infinitePathInRauzyGraph} which may reveal non-convergence of Phase~2 by searching for a cycle in the Rauzy graph.

We indicated that there is the same lower bound on the size of an alphabet from $\Zbeta$ as of an integer one. The necessary conditions on the alphabet which allows parallel addition were summarized and the way of generating such an alphabet was sketched.

Algorithms based on the obtained theoretical results were designed. Algorithm~\ref{alg:weightFunction_modified} summarizes Phase~2 with the control of non-convergence. The extending window method was implemented with all proposed algorithms in SageMath.

The shell interfaces are provided for running the implementation. All results are automatically saved to Google spreadsheet which enables easy sorting.

We have tested large number of numeration systems with different methods for both phases. The controls of non-convergence reveal failure of Phase~2 very efficiently. Thus, many inputs may be processed. We compared various methods for Phase~1 and 2 for selected numeration systems including those for which a parallel addition algorithm was found manually.

 To conclude, there is about 140 numeration systems for which the implemented extending window method successfully found a parallel addition algorithm. The numeration system have integer or non-integer alphabet, often of the minimal size.

The question of sufficient condition of convergence of Phase~2 remains open.




\newpage
\bibliography{literatura}
\addcontentsline{toc}{chapter}{References}
\bibliographystyle{amsplain}

%\cleardoublepage\phantomsection
\appendix
\chapter*{Appendices}
\markboth{\MakeUppercase{Appendices}}{}
%\pagenumbering{Roman}
\addcontentsline{toc}{chapter}{Appendices}
\addtocontents{toc}{\protect\setcounter{tocdepth}{0}}
\renewcommand{\thesection}{\Alph{section}}
\counterwithin{exmp}{section}
\section{Illustration of Phase 1}
\label{app:phase1}   
Figures \ref{img:phase1img3} -- \ref{img:phase1img12} illustrates first and last iterations of the construction of the weight coefficients set $\Q$ for the Eisenstein base $\beta = -\frac{3}{2} + \frac{\imath \sqrt{3}}{2}$ with the complex alphabet $\mathcal{A} =\{0, 1, -1, \omega, -\omega, -\omega - 1, \omega + 1\}$ and input alphabet $\B=\A+\A$. The second last iteration is skipped.

\figurehascaptionOne{1 = The starting set $\Q_0{=}\{0\}$.}
\figurehascaptionOne{2 = The set $\B+\Q_0$ need to be covered.}
\figurehascaptionOne{3 = The set $\Q_0$ does not cover the set $\B+\Q_0${,} i.e.{,} the set $\A+\beta \cdot \Q_0$ is not superset of $\B+\Q_0$.}
\figurehascaptionOne{4 = The set $\Q_0$ is extended to $\Q_1$ to cover all elements of $\B+\Q_0$.}
\figurehascaptionOne{5 = The set $\B+\Q_1$ need to be covered.}
\figurehascaptionOne{6 = The set $\Q_1$ does not cover the set $\B+\Q_1${,} i.e.{,} the set $\A+\beta \cdot \Q_1$ is not superset of $\B+\Q_1$.}
\figurehascaptionOne{7 = The set $\Q_1$ is extended to $\Q_2$ to cover all elements of $\B+\Q_1$.}
\figurehascaptionOne{8 = The set $\B+\Q_2$ need to be covered.}
\figurehascaptionOne{9 = The set $\Q_2$ does not cover the set $\B+\Q_2${,} i.e.{,} the set $\A+\beta \cdot \Q_2$ is not superset of $\B+\Q_2$.}
\figurehascaptionOne{10 = The set $\Q_2$ is extended to $\Q_3$ to cover all elements of $\B+\Q_2$.}
\figurehascaptionOne{11 = The set $\B+\Q_3$ need to be covered.}
\figurehascaptionOne{12 = In the last iteration{,} the set $\Q_3$ covers the set $\B+\Q_3${,} i.e.{,} the set $\A+\beta \cdot \Q_3$ is superset of $\B+\Q_2$. The weight coefficients set $\Q$ equals $\Q_3$.}
\figurehascaptionOne{13 = The final weight coefficients $\Q{=}\Q_3$.}


\foreach \n in {3,4,6,7,12} {%
\begin{SCfigure}[][htbp]
    \centering
    \caption{\getcaptionOne{\n}}
    \label{img:phase1img\n}
    \includegraphics[height=0.27\textheight]{img/eisenstein/phase1_image_\n.png}
\end{SCfigure}
    }

\newpage


\section{Illustration of Phase 2}
The construction of set $\Q_{[\omega,1,2]}$ for the Eisenstein base $\beta = -\frac{3}{2} + \frac{\imath \sqrt{3}}{2}$ with the complex alphabet $\mathcal{A} =\{0, 1, -1, \omega, -\omega, -\omega - 1, \omega + 1\}$  and input alphabet $\B=\A+\A$ is illustrated on Figures \ref{img:phase2img3} -- \ref{img:phase2img7}.
\label{app:phase2}    

\figurehascaptionTwo{1 = Phase 2 starts with the weight coefficients set $\Q$ from Phase 1.}
\figurehascaptionTwo{2 = The set $\omega+\Q$ need to be covered.}
\figurehascaptionTwo{3 = The elements of $\omega+\Q$ are covered by the set $\Q_{[\omega]}\subset\Q$.}
\figurehascaptionTwo{4 = The set $\omega+\Q_{[1]}$ need to be covered.}
\figurehascaptionTwo{5 = The elements of $\omega+\Q_{[1]}$ are covered by the set $\Q_{[\omega,1]}\subset\Q_{[\omega]}$.}
\figurehascaptionTwo{6 = The set $\omega+\Q_{[1,2]}$ need to be covered.}
\figurehascaptionTwo{7 = The elements of $\omega+\Q_{[1]}$ are covered by the set $\Q_{[\omega,1,2]}\subset\Q_{[\omega,1]}$ which has only one element{.} This element is the output of the weight function $q{(\omega,1,2)}$.}



\foreach \n in {3,5,7} {%
\begin{SCfigure}[][htbp]
    \centering
    \caption{\getcaptionTwo{\n}}
    \label{img:phase2img\n}
    \includegraphics[height=0.23\textheight]{img/eisenstein/phase2_image_\n.png}
\end{SCfigure}
    }

\newpage

\section{Interfaces}
\label{app:interfaces}
File \verb+ewm_inputs.sage+:
\lstinputlisting[language=Python]{ewm_inputs.sage}

File \verb+ewm_gspreadsheet.sage+:
\lstinputlisting[language=Python]{ewm_gspreadsheet.sage}




\section{Tested examples}
\label{app:examples}

\subsection*{Unsuccessful examples comparing different methods}
The reasons of failure of Phase~2 for numeration systems in Section~\ref{sec:compareMethods} can be found here. See Tables \ref{tab:resultsPhaseOne}, \ref{tab:resultsPhaseTwo} and \ref{tab:alphabets} for parameters of the numeration systems.
\begin{exmp}
\label{ex:compareAB}

\textbf{Eisenstein\_1--block\_integer}

Phase 1 (methods 1b, 1c, 1d, 1e):
\begin{enumerate}[ ]
\item  2a -- Check of $b,b,\dots,b$ inputs fails for $b\in \{2, 3, 5, 6, -5, -4, -3\}$.
\item  2b -- Check of $b,b,\dots,b$ inputs fails for $b\in \{2, 3, 6, -6, -4, -3\}$.
\item  2c -- Check of $b,b,\dots,b$ inputs fails for $b\in \{0, 1, 3, 4, 6, -6, -4, -3, -1\}$.
\item  2d -- Check of $b,b,\dots,b$ inputs fails for $b\in \{3, 4, 6, -6, -4, -3\}$.
\item  2e -- Check of $b,b,\dots,b$ inputs fails for $b\in \{2, 3, 6, -2, -6, -4, -3\}$.
\end{enumerate}


Phase 1 (methods 1a):
\begin{enumerate}[ ]
\item  2a -- Check of $b,b,\dots,b$ inputs fails for $b\in \{0, 2, 4, 5, -2, -5, -4\}$.
\item  2b -- Check of $b,b,\dots,b$ inputs fails for $b\in \{0, 2, 4, 5, -2, -5, -4\}$.
\item  2c -- Check of $b,b,\dots,b$ inputs fails for $b\in \{0, 2, 3, 5, 6, -2, -6, -5, -3\}$.
\item  2d -- Check of $b,b,\dots,b$ inputs fails for $b\in \{0, 3, 4, 5, 6, -6, -5, -4, -3\}$.
\item  2e -- Check of $b,b,\dots,b$ inputs fails for $b\in \{0, 2, 4, 5, -2, -5, -4\}$.
\end{enumerate}


\end{exmp}




\begin{exmp}
\label{ex:compareAC}

\textbf{Eisenstein\_2--block\_complex}

Phase 1 (methods 1a, 1b, 1c, 1d, 1e):
\begin{enumerate}[ ]
\item  2a -- Check of $b,b,\dots,b$ inputs fails for $b\in \{2\omega - 1, \omega + 1, -2\omega, -4, -\omega - 2\}$.
\item  2b -- Check of $b,b,\dots,b$ inputs fails for $b\in \{2\omega - 1, \omega + 1, -2\omega, -4, -\omega - 2\}$.
\item  2c -- Check of $b,b,\dots,b$ inputs fails for $b\in \{2\omega - 1, \omega + 1, -2\omega, -4, -\omega - 2\}$.
\item  2d -- Check of $b,b,\dots,b$ inputs fails for $b\in \{2\omega - 1, \omega + 1, -2\omega, -4, -\omega - 2\}$.
\item  2e -- Check of $b,b,\dots,b$ inputs fails for $b\in \{2\omega - 1, \omega + 1, -2\omega, -4, -\omega - 2\}$.
\end{enumerate}


\end{exmp}




\begin{exmp}
\label{ex:compareAD}

\textbf{Eisenstein\_2--block\_integer}

Phase 1 (methods 1a, 1b, 1c, 1d, 1e):
\begin{enumerate}[ ]
\item  2a -- Check of $b,b,\dots,b$ inputs fails for $b\in \{0, 1, 2, 2\omega - 4, \omega - 2, 4\omega, 3\omega - 5, \omega - 1, -\omega + 3, -2\omega + 5, 2\omega - 3\}$.
\item  2b -- Check of $b,b,\dots,b$ inputs fails for $b\in \{0, 1, 2, 2\omega - 4, 2\omega - 2, \omega - 2, \omega, 4\omega, 3\omega - 5, \omega - 1, -\omega + 3, -2\omega + 5, 2\omega - 3\}$.
\item  2c -- Check of $b,b,\dots,b$ inputs fails for $b\in \{0, 1, 2, 2\omega - 4, \omega - 2, \omega, 4\omega, 3\omega - 5, \omega - 1, -\omega + 3, -2\omega + 5, 2\omega - 3\}$.
\item  2d -- Check of $b,b,\dots,b$ inputs fails for $b\in \{1, 2, 2\omega - 4, \omega - 2, 4\omega, 3\omega - 5, \omega - 1, -\omega + 3, -2\omega + 5, 2\omega - 3\}$.
\item  2e -- Check of $b,b,\dots,b$ inputs fails for $b\in \{0, 1, 2, 2\omega - 4, 2\omega - 2, \omega - 2, \omega, 4\omega, 3\omega - 5, \omega - 1, -\omega + 3, -2\omega + 5, 2\omega - 3\}$.
\end{enumerate}


\end{exmp}




\begin{exmp}
\label{ex:compareAF}

\textbf{Penney\_1--block\_integer}

Phase 1 (methods 1b, 1c, 1d, 1e):
\begin{enumerate}[ ]
\item  2a -- Check of $b,b,\dots,b$ inputs fails for $b\in \{0, 3, 4, -4, -3\}$.
\item  2b -- Check of $b,b,\dots,b$ inputs fails for $b\in \{0, 3, -3\}$.
\item  2c -- Check of $b,b,\dots,b$ inputs fails for $b\in \{0, 3, 4, -4, -3, -2\}$.
\item  2d -- Check of $b,b,\dots,b$ inputs fails for $b\in \{0, 3, 4, -4, -3, -2\}$.
\item  2e -- Check of $b,b,\dots,b$ inputs fails for $b\in \{0, 3, -3\}$.
\end{enumerate}


Phase 1 (methods 1a):
\begin{enumerate}[ ]
\item  2a -- Check of $b,b,\dots,b$ inputs fails for $b\in \{0, 1, 2, -1, -2\}$.
\item  2b -- Check of $b,b,\dots,b$ inputs fails for $b\in \{0, 1, 2, -1, -2\}$.
\item  2c -- Check of $b,b,\dots,b$ inputs fails for $b\in \{0, 1, 2, 3, 4, -4, -3, -2\}$.
\item  2d -- Check of $b,b,\dots,b$ inputs fails for $b\in \{0, 1, 2, 3, 4, -4, -3, -2\}$.
\item  2e -- Check of $b,b,\dots,b$ inputs fails for $b\in \{2, -2\}$.
\end{enumerate}


\end{exmp}




\begin{exmp}
\label{ex:compareAK}

\textbf{Quadratic+1+0--5\_integer}

Phase 1 (methods 1a, 1b, 1c, 1d, 1e):
\begin{enumerate}[ ]
\item  2a -- Check of $b,b,\dots,b$ inputs fails for $b\in \{-2\}$.
\end{enumerate}


\end{exmp}




\begin{exmp}
\label{ex:compareAL}

\textbf{Quadratic+1+2+3\_complex}

Phase 1 (methods 1a, 1b, 1c, 1d, 1e):
\begin{enumerate}[ ]
\item  2a -- Phase 2   fails because  the sequence $(1, -2\omega - 2, -\omega - 2, 2\omega + 2, -2\omega - 2, -\omega - 2, 2\omega + 2, -2\omega - 2, \dots ,-2\omega - 2, -\omega - 2, 2\omega + 2, -2\omega - 2, \dots)$ leads to an infinite loop.
\item  2c -- Phase 2   fails because  the sequence $(1, -2\omega - 2, -\omega - 2, 2\omega + 2, -2\omega - 2, -\omega - 2, 2\omega + 2, -2\omega - 2, \dots ,-2\omega - 2, -\omega - 2, 2\omega + 2, -2\omega - 2, \dots)$ leads to an infinite loop.
\item  2d -- Phase 2   fails because  the sequence $(0, \omega + 2, -1, -1, 2\omega + 1, -1, -1, 2\omega + 1, -1, \dots ,-1, -1, 2\omega + 1, -1, \dots)$ leads to an infinite loop.
\end{enumerate}


\end{exmp}




\begin{exmp}
\label{ex:compareAM}

\textbf{Quadratic+1+3+4\_complex}

Phase 1 (methods 1b):
\begin{enumerate}[ ]
\item  2c -- Phase 2   fails because  the sequence $(2, 2\omega + 1, -\omega - 2, -2, 2\omega, 2\omega + 1, -\omega - 2, -2, 2\omega, \dots ,2\omega + 1, -\omega - 2, -2, 2\omega, \dots)$ leads to an infinite loop.
\item  2d -- Check of $b,b,\dots,b$ inputs fails for $b\in \{\omega + 3\}$.
\end{enumerate}


Phase 1 (methods 1c):
\begin{enumerate}[ ]
\item  2a -- Phase 2   fails because  the sequence $(2, 2\omega + 3, -2\omega - 2, 2, 2\omega + 3, -2\omega - 2, 2, \dots ,2, 2\omega + 3, -2\omega - 2, 2, \dots)$ leads to an infinite loop.
\item  2c -- Phase 2   fails because  the sequence $(2, 2\omega + 2, -\omega + 1, 2, 2\omega + 2, -\omega + 1, 2, \dots ,2, 2\omega + 2, -\omega + 1, 2, \dots)$ leads to an infinite loop.
\item  2d -- Check of $b,b,\dots,b$ inputs fails for $b\in \{\omega + 3\}$.
\end{enumerate}


Phase 1 (methods 1d, 1e):
\begin{enumerate}[ ]
\item  2a -- Phase 2   fails because  the sequence $(2, 2\omega + 3, -2\omega - 2, 2, 2\omega + 3, -2\omega - 2, 2, \dots ,2, 2\omega + 3, -2\omega - 2, 2, \dots)$ leads to an infinite loop.
\item  2c -- Phase 2   fails because  the sequence $(2, 2\omega + 2, -\omega + 1, 2, 2\omega + 2, -\omega + 1, 2, \dots ,2, 2\omega + 2, -\omega + 1, 2, \dots)$ leads to an infinite loop.
\item  2d -- Check of $b,b,\dots,b$ inputs fails for $b\in \{\omega + 3\}$.
\end{enumerate}


Phase 1 (methods 1a):
\begin{enumerate}[ ]
\item  2c -- Phase 2   fails because  the sequence $(2, 2\omega + 2, -\omega + 1, 2, 2\omega + 2, -\omega + 1, 2, \dots ,2, 2\omega + 2, -\omega + 1, 2, \dots)$ leads to an infinite loop.
\item  2d -- Check of $b,b,\dots,b$ inputs fails for $b\in \{\omega + 3\}$.
\end{enumerate}


\end{exmp}




\begin{exmp}
\label{ex:compareAN}

\textbf{Quadratic+1+3+5\_complex1}

Phase 1 (methods 1a):
\begin{enumerate}[ ]
\item  2a -- Check of $b,b,\dots,b$ inputs fails for $b\in \{2\omega + 2, \omega + 4, -\omega - 4, -2\omega - 2\}$.
\item  2b -- Check of $b,b,\dots,b$ inputs fails for $b\in \{2\omega + 2, -2\omega - 2\}$.
\item  2c -- Check of $b,b,\dots,b$ inputs fails for $b\in \{2\omega + 2, \omega + 4, -\omega - 4, -2\omega - 2\}$.
\item  2d -- Check of $b,b,\dots,b$ inputs fails for $b\in \{2\omega + 2, 2\omega + 4, \omega + 4, -\omega - 4, -2\omega - 4, -2\omega - 2\}$.
\item  2e -- Check of $b,b,\dots,b$ inputs fails for $b\in \{2\omega + 2, -2\omega - 2\}$.
\end{enumerate}


Phase 1 (methods 1b, 1c):
\begin{enumerate}[ ]
\item  2a -- Check of $b,b,\dots,b$ inputs fails for $b\in \{-2\omega - 2\}$.
\item  2b -- Phase 2   fails because  the sequence $(2, 2\omega + 2, 2\omega + 2, 2, 2\omega + 2, 2\omega + 2, \dots ,2, 2\omega + 2, 2\omega + 2, \dots)$ leads to an infinite loop.
\item  2c -- Check of $b,b,\dots,b$ inputs fails for $b\in \{2\omega + 2, -2\omega - 2\}$.
\item  2d -- Check of $b,b,\dots,b$ inputs fails for $b\in \{2\omega + 2, -2\omega - 2\}$.
\item  2e -- Phase 2   fails because  the sequence $(2, 2\omega + 2, 2\omega + 2, 2, 2\omega + 2, 2\omega + 2, \dots ,2, 2\omega + 2, 2\omega + 2, \dots)$ leads to an infinite loop.
\end{enumerate}


Phase 1 (methods 1d, 1e):
\begin{enumerate}[ ]
\item  2a -- Check of $b,b,\dots,b$ inputs fails for $b\in \{2\omega + 2, -2\omega - 2\}$.
\item  2b -- Check of $b,b,\dots,b$ inputs fails for $b\in \{2\omega + 2, -2\omega - 2\}$.
\item  2c -- Check of $b,b,\dots,b$ inputs fails for $b\in \{2\omega + 2, -2\omega - 2\}$.
\item  2d -- Check of $b,b,\dots,b$ inputs fails for $b\in \{2\omega + 2, 2\omega + 4, -2\omega - 4, -2\omega - 2\}$.
\item  2e -- Check of $b,b,\dots,b$ inputs fails for $b\in \{2\omega + 2, -2\omega - 2\}$.
\end{enumerate}


\end{exmp}




\begin{exmp}
\label{ex:compareAO}

\textbf{Quadratic+1+3+5\_complex2 }

Phase 1 (methods 1b, 1c):
\begin{enumerate}[ ]
\item  2a -- Check of $b,b,\dots,b$ inputs fails for $b\in \{-3\omega - 4, 2\omega + 2, 2\omega + 3, \omega + 3, -2\omega - 4, -2\omega - 3, -2\omega - 2, -\omega - 3\}$.
\item  2b -- Phase 2   fails because  the sequence $(0, 1, 2\omega + 1, -4\omega - 4, 4\omega + 4, 0, 1, 4\omega + 4, 0, 1, 4\omega + 4, \dots ,4\omega + 4, 0, 1, 4\omega + 4, \dots)$ leads to an infinite loop.
\item  2c -- Check of $b,b,\dots,b$ inputs fails for $b\in \{-3\omega - 3, 3\omega + 3\}$.
\item  2d -- Check of $b,b,\dots,b$ inputs fails for $b\in \{-3\omega - 4, 2\omega + 3, 2\omega + 4, \omega + 3, -2\omega - 4, -2\omega - 3, -\omega - 3, 3\omega + 4\}$.
\item  2e -- Phase 2   fails because  the sequence $(0, 1, 2\omega + 1, -4\omega - 4, 4\omega + 4, 0, 1, 4\omega + 4, 0, 1, 4\omega + 4, \dots ,4\omega + 4, 0, 1, 4\omega + 4, \dots)$ leads to an infinite loop.
\end{enumerate}


Phase 1 (methods 1d, 1e):
\begin{enumerate}[ ]
\item  2a -- Check of $b,b,\dots,b$ inputs fails for $b\in \{2\omega + 2, -2\omega - 4, -2\omega - 3, -2\omega - 2, -\omega - 3\}$.
\item  2b -- Phase 2   fails because  the sequence $(0, 0, -4\omega - 4, 3\omega + 4, 0, 4\omega + 4, 0, 4\omega + 4, 0, 4\omega + 4, \dots ,0, 4\omega + 4, 0, 4\omega + 4, \dots)$ leads to an infinite loop.
\item  2c -- Phase 2   fails because  the sequence $(0, 0, -2\omega - 1, 3\omega + 3, -4\omega - 4, \omega, -2\omega - 3, 4\omega + 4, -\omega, 2\omega + 3, -3\omega - 3, 4\omega + 4, 2\omega + 2, 4\omega + 4, 2\omega + 2, 4\omega + 4, 2\omega + 2, \dots ,4\omega + 4, 2\omega + 2, 4\omega + 4, 2\omega + 2, \dots)$ leads to an infinite loop.
\item  2d -- Check of $b,b,\dots,b$ inputs fails for $b\in \{-3\omega - 4, 2\omega + 3, 2\omega + 4, \omega + 3, -2\omega - 4, -2\omega - 3, -\omega - 3, 3\omega + 4\}$.
\item  2e -- Phase 2   fails because  the sequence $(0, -3\omega - 4, 3\omega + 3, -3\omega - 4, -4\omega - 4, 4\omega + 4, 0, 0, 4\omega + 4, -3\omega - 4, -1, -3\omega - 4, -1, -3\omega - 4, -1, \dots ,-3\omega - 4, -1, -3\omega - 4, -1, \dots)$ leads to an infinite loop.
\end{enumerate}


Phase 1 (methods 1a):
\begin{enumerate}[ ]
\item  2a -- Check of $b,b,\dots,b$ inputs fails for $b\in \{-3\omega - 3, 2\omega + 2, -2\omega - 3, -\omega - 3, 3\omega + 3\}$.
\item  2b -- Phase 2   fails because  the sequence $(0, -1, -3\omega - 3, \omega, -3\omega - 4, -3\omega - 4, -3\omega - 3, 4\omega + 4, -3\omega - 4, 2\omega + 1, 0, -1, -3\omega - 3, \omega, \dots ,0, -1, -3\omega - 3, \omega, \dots)$ leads to an infinite loop.
\item  2c -- Phase 2   fails because  the sequence $(0, 0, -4\omega - 4, 2\omega + 1, -\omega - 2, -3\omega - 4, 2\omega + 2, -\omega - 1, -\omega, 3\omega + 3, -2\omega - 4, 2\omega + 1, -\omega - 2, -3\omega - 4, 2\omega + 2, \dots ,2\omega + 1, -\omega - 2, -3\omega - 4, 2\omega + 2, \dots)$ leads to an infinite loop.
\item  2d -- Check of $b,b,\dots,b$ inputs fails for $b\in \{2\omega + 3, \omega + 3, -2\omega - 3, -\omega - 3\}$.
\item  2e -- Phase 2   fails because  the sequence $(0, -3\omega - 4, 3\omega + 3, -\omega, 2\omega + 1, -2\omega - 1, 1, 2\omega + 1, 0, -1, -3\omega - 3, \omega, -3\omega - 4, -3\omega - 4, -3\omega - 3, 4\omega + 4, -3\omega - 4, 2\omega + 1, 0, -1, -3\omega - 3, \dots ,2\omega + 1, 0, -1, -3\omega - 3, \dots)$ leads to an infinite loop.
\end{enumerate}


\end{exmp}




\begin{exmp}
\label{ex:compareAP}

\textbf{Quadratic+1+4+5\_complex1}

Phase 1 (Lemma \ref{lem:suffCondPhase1}):
\begin{enumerate}[ ]
\item  2a -- Phase 2   fails because  the sequence $(2, 2\omega - 1, -2\omega + 1, 2\omega - 2, -\omega + 2, 2, 2\omega - 2, -\omega + 2, 2, 2\omega - 2, \dots ,2\omega - 2, -\omega + 2, 2, 2\omega - 2, \dots)$ leads to an infinite loop.
\item  2b -- Check of $b,b,\dots,b$ inputs fails for $b\in \{-\omega - 2, -2\omega + 1\}$.
\item  2c -- Check of $b,b,\dots,b$ inputs fails for $b\in \{-\omega - 2\}$.
\item  2d -- Check of $b,b,\dots,b$ inputs fails for $b\in \{-\omega - 2\}$.
\end{enumerate}


Phase 1 (methods 1a):
\begin{enumerate}[ ]
\item  2a -- Check of $b,b,\dots,b$ inputs fails for $b\in \{2, -\omega - 2, -2\omega\}$.
\item  2b -- Check of $b,b,\dots,b$ inputs fails for $b\in \{-\omega - 2\}$.
\item  2c -- Check of $b,b,\dots,b$ inputs fails for $b\in \{2\omega - 2, 2\omega + 2, \omega + 3, -\omega - 2, -2\omega + 2\}$.
\item  2d -- Check of $b,b,\dots,b$ inputs fails for $b\in \{2\omega - 2, 2\omega + 2, \omega + 3, -\omega - 2, -2\omega + 2\}$.
\item  2e -- Check of $b,b,\dots,b$ inputs fails for $b\in \{-\omega - 2\}$.
\end{enumerate}


Phase 1 (methods 1b, 1c, 1d, 1e):
\begin{enumerate}[ ]
\item  2a -- Check of $b,b,\dots,b$ inputs fails for $b\in \{2, -\omega - 2, -2\omega\}$.
\item  2b -- Check of $b,b,\dots,b$ inputs fails for $b\in \{-\omega - 2\}$.
\item  2c -- Check of $b,b,\dots,b$ inputs fails for $b\in \{2\omega + 2, -\omega - 2\}$.
\item  2d -- Check of $b,b,\dots,b$ inputs fails for $b\in \{2\omega + 2, -\omega - 2\}$.
\item  2e -- Check of $b,b,\dots,b$ inputs fails for $b\in \{-\omega - 2\}$.
\end{enumerate}


\end{exmp}




\begin{exmp}
\label{ex:compareAR}

\textbf{Cubic+1+0+0+2\_integer}

Phase 1 (methods 1a, 1b, 1c, 1d, 1e):
\begin{enumerate}[ ]
\item  2a -- Phase 2   fails because  the sequence $(0, 1, 0, 1, 0, 1, 0, \dots ,0, 1, 0, 1, 0, \dots)$ leads to an infinite loop.
\item  2b -- Phase 2   fails because  the sequence $(0, 1, 0, -1, 0, 1, 0, -1, 0, \dots ,0, 1, 0, -1, 0, \dots)$ leads to an infinite loop.
\item  2c -- Phase 2   fails because  the sequence $(0, 1, 0, -1, 0, 1, 0, -1, 0, \dots ,0, 1, 0, -1, 0, \dots)$ leads to an infinite loop.
\item  2e -- Phase 2   fails because  the sequence $(0, 1, 0, -1, 0, 1, 0, -1, 0, \dots ,0, 1, 0, -1, 0, \dots)$ leads to an infinite loop.
\end{enumerate}


\end{exmp}




\begin{exmp}
\label{ex:compareAS}

\textbf{Cubic+1+0+0--2\_integer}

Phase 1 (methods 1a, 1b, 1c, 1d, 1e):
\begin{enumerate}[ ]
\item  2a -- Phase 2   fails because  the sequence $(0, 1, 0, 1, 0, 1, 0, \dots ,0, 1, 0, 1, 0, \dots)$ leads to an infinite loop.
\item  2b -- Phase 2   fails because  the sequence $(0, 1, 0, -1, 0, 1, 0, -1, 0, \dots ,0, 1, 0, -1, 0, \dots)$ leads to an infinite loop.
\item  2c -- Phase 2   fails because  the sequence $(0, 1, 0, -1, 0, 1, 0, -1, 0, \dots ,0, 1, 0, -1, 0, \dots)$ leads to an infinite loop.
\item  2e -- Phase 2   fails because  the sequence $(0, 1, 0, -1, 0, 1, 0, -1, 0, \dots ,0, 1, 0, -1, 0, \dots)$ leads to an infinite loop.
\end{enumerate}


\end{exmp}






\newpage
\subsection*{Quadratic bases with integer alphabet}
The following examples show alphabets  divided into congruence classes modulo $\beta$ and $\beta-1$ for some numeration systems in Table~\ref{tab:resultsQuadrInt}.



\begin{exmp}
\label{ex:integerAB}

An alphabet $\A$ divided into congruence classes modulo $\beta$: 
$$ \left\{\left\{-5, 6\right\}, \left\{-4\right\}, \left\{-3\right\}, \left\{-2\right\}, \left\{-1\right\}, \left\{0\right\}, \left\{1\right\}, \left\{2\right\}, \left\{3\right\}, \left\{4\right\}, \left\{5\right\}\right\} \,,$$

and modulo $\beta - 1$: $$ \left\{\left\{-5, 1\right\}, \left\{-4, 2\right\}, \left\{-3, 3\right\}, \left\{-2, 4\right\}, \left\{-1, 5\right\}, \left\{0, 6\right\}\right\} \,.$$

\end{exmp}




%
%\begin{exmp}
%\label{ex:integerAD}
%
%An alphabet $\A$ divided into congruence classes modulo $\beta$: 
%$$ \left\{\left\{-3, 4\right\}, \left\{-2\right\}, \left\{-1\right\}, \left\{0\right\}, \left\{1\right\}, \left\{2\right\}, \left\{3\right\}\right\} \,,$$
%
%and modulo $\beta - 1$: $$ \left\{\left\{-3, 1\right\}, \left\{-2, 2\right\}, \left\{-1, 3\right\}, \left\{0, 4\right\}\right\} \,.$$
%
%\end{exmp}


\begin{exmp}
\label{ex:integerAE}

An alphabet $\A$ divided into congruence classes modulo $\beta$: 
$$ \left\{\left\{-5, 2\right\}, \left\{-4, 3\right\}, \left\{-3, 4\right\}, \left\{-2, 5\right\}, \left\{-1\right\}, \left\{0\right\}, \left\{1\right\}\right\} \,,$$

and modulo $\beta - 1$: $$ \left\{\left\{-5, 2\right\}, \left\{-4, 3\right\}, \left\{-3, 4\right\}, \left\{-2, 5\right\}, \left\{-1\right\}, \left\{0\right\}, \left\{1\right\}\right\} \,.$$

\end{exmp}




%\begin{exmp}
%\label{ex:integerAG}
%
%An alphabet $\A$ divided into congruence classes modulo $\beta$: 
%$$ \left\{\left\{-1, 1\right\}, \left\{0\right\}\right\} \,,$$
%
%and modulo $\beta - 1$: $$ \left\{\left\{-1\right\}, \left\{0\right\}, \left\{1\right\}\right\} \,.$$
%
%\end{exmp}
%
%
%\begin{exmp}
%\label{ex:integerAH}
%
%An alphabet $\A$ divided into congruence classes modulo $\beta$: 
%$$ \left\{\left\{0\right\}, \left\{1, -1\right\}\right\} \,,$$
%
%and modulo $\beta - 1$: $$ \left\{\left\{0, 1, -1\right\}\right\} \,.$$
%
%\end{exmp}




\begin{exmp}
\label{ex:integerAJ}

An alphabet $\A$ divided into congruence classes modulo $\beta$: 
$$ \left\{\left\{-2, 3\right\}, \left\{-1\right\}, \left\{0\right\}, \left\{1\right\}, \left\{2\right\}\right\} \,,$$

and modulo $\beta - 1$: $$ \left\{\left\{-2, 0, 2\right\}, \left\{-1, 1, 3\right\}\right\} \,.$$

\end{exmp}


\begin{exmp}
\label{ex:integerAK}

An alphabet $\A$ divided into congruence classes modulo $\beta$: 
$$ \left\{\left\{-2, 3\right\}, \left\{-1\right\}, \left\{0\right\}, \left\{1\right\}, \left\{2\right\}\right\} \,,$$

and modulo $\beta - 1$: $$ \left\{\left\{-2, 2\right\}, \left\{-1, 3\right\}, \left\{0\right\}, \left\{1\right\}\right\} \,.$$

\end{exmp}



\begin{exmp}
\label{ex:integerAO}

An alphabet $\A$ divided into congruence classes modulo $\beta$: 
$$ \left\{\left\{-7, 6\right\}, \left\{-6, 7\right\}, \left\{-5\right\}, \left\{-4\right\}, \left\{-3\right\}, \left\{-2\right\}, \left\{-1\right\}, \left\{0\right\}, \left\{1\right\}, \left\{2\right\}, \left\{3\right\}, \left\{4\right\}, \left\{5\right\}\right\} \,,$$

and modulo $\beta - 1$: $$ \left\{\left\{-7, -1, 5\right\}, \left\{-6, 0, 6\right\}, \left\{-5, 1, 7\right\}, \left\{-4, 2\right\}, \left\{-3, 3\right\}, \left\{-2, 4\right\}\right\} \,.$$

\end{exmp}



%\begin{exmp}
%\label{ex:integerAQ}
%
%An alphabet $\A$ divided into congruence classes modulo $\beta$: 
%$$ \left\{\left\{-8, 9\right\}, \left\{-7\right\}, \left\{-6\right\}, \left\{-5\right\}, \left\{-4\right\}, \left\{-3\right\}, \left\{-2\right\}, \left\{-1\right\}, \left\{0\right\}, \left\{1\right\}, \left\{2\right\}, \left\{3\right\}, \left\{4\right\}, \left\{5\right\}, \left\{6\right\}, \left\{7\right\}, \left\{8\right\}\right\} \,,$$
%
%and modulo $\beta - 1$: $$ \left\{\left\{-8, 0, 8\right\}, \left\{-7, 1, 9\right\}, \left\{-6, 2\right\}, \left\{-5, 3\right\}, \left\{-4, 4\right\}, \left\{-3, 5\right\}, \left\{-2, 6\right\}, \left\{-1, 7\right\}\right\} \,.$$
%
%\end{exmp}


\begin{exmp}
\label{ex:integerAR}

An alphabet $\A$ divided into congruence classes modulo $\beta$: 
\begin{align*}
\{\left\{-10, 11\right\}, \left\{-9\right\}, \left\{-8\right\}, \left\{-7\right\},& \left\{-6\right\}, \left\{-5\right\}, \left\{-4\right\}, \left\{-3\right\}, \left\{-2\right\},  \\
\left\{-1\right\}, \left\{0\right\}, \left\{1\right\}, \left\{2\right\}, \left\{3\right\}, \left\{4\right\},& \left\{5\right\}, \left\{6\right\}, \left\{7\right\}, \left\{8\right\}, \left\{9\right\}, \left\{10\right\}\} \,,
\end{align*}

and modulo $\beta - 1$: \begin{align*} 
\{\left\{-10, 0, 10\right\}, \left\{-9, 1, 11\right\},& \left\{-8, 2\right\}, \left\{-7, 3\right\}, \left\{-6, 4\right\}, \\
 \left\{-5, 5\right\}, \left\{-4, 6\right\}, \{-3,\, & 7\}, \left\{-2, 8\right\}, \left\{-1, 9\right\}\} \,.
 \end{align*}

\end{exmp}




\subsection*{Killed examples}
The computation of a weight function for the following numeration systems was killed because of memory limits.
\begin{exmp}
\label{ex:killAB}

\rule{0cm}{0cm}

\begin{tabular}{ll}
$\omega=  -\frac{1}{2} \, \sqrt{37} + \frac{5}{2} $  & $\beta= -\omega - 3 = \frac{1}{2} \, \sqrt{37} - \frac{11}{2} $\\
$m_\omega(t)=  t^{2} - 5 \, t - 3 $  & $m_\beta(x)=  x^{2} + 11 \, x + 21 $\\
Real conjugate of $\beta$ greater than 1:   &  no \\
$\#\A= $ 33 $ $ & $\A$ is minimal. \\
\multicolumn{2}{l}{\begin{minipage}{\textwidth}\begin{dmath*}\A = \left\{0, 1, -1, \omega + 1, -\omega - 1, -\omega + 1, \omega - 1, \omega, -\omega, 2 \, \omega + 2, -2 \, \omega - 2, \omega + 2, -\omega - 2, -2 \, \omega + 2, 2 \, \omega - 2, 2 \, \omega + 1, -2 \, \omega - 1, -2 \, \omega + 1, 2 \, \omega - 1, 2 \, \omega, -2 \, \omega, -2 \, \omega + 3, 2 \, \omega - 3, -3 \, \omega + 3, 3 \, \omega - 3, -3 \, \omega + 2, 3 \, \omega - 2, -3 \, \omega + 1, 3 \, \omega - 1, 3 \, \omega, -3 \, \omega, -3 \, \omega + 4, 3 \, \omega - 4\right\}  \end{dmath*}\end{minipage} }\\
 & \\
Phase 1 (method  9): &
\checkmark, $\#\mathcal{Q} = $ 17 $ $ \\ 
$b,b,\dots,b$ inputs (method  2c): & \checkmark, maximal length of window: $ 3 $ \\
\multicolumn{2}{l}{\begin{minipage}{\textwidth} Computation of Phase 2 (method  2c) was killed when the length of window 5 was being proccessed. Numbers of saved combinations for each finished length are: (0, 12399, 682670, 2721482)\end{minipage} }\\
\end{tabular}

\end{exmp}



\begin{exmp}
\label{ex:killAA}

\rule{0cm}{0cm}

\begin{tabular}{ll}
$\omega=  -\frac{1}{2} \, \sqrt{29} + \frac{3}{2} $  & $\beta= 3 \, \omega + 1 = -\frac{3}{2} \, \sqrt{29} + \frac{11}{2} $\\
$m_\omega(t)=  t^{2} - 3 \, t - 5 $  & $m_\beta(x)=  x^{2} - 11 \, x - 35 $\\
Real conjugate of $\beta$ greater than 1:   &  yes \\
$\#\A= $ 49 $ $ & $\A$ is not minimal. \\
\multicolumn{2}{l}{\begin{minipage}{\textwidth}\begin{dmath*}\A = \left\{0, 1, -1, \omega + 1, -\omega - 1, -\omega + 1, \omega - 1, \omega, -\omega, 2 \, \omega + 2, -2 \, \omega - 2, \omega + 2, -\omega - 2, 2, -2, 3 \, \omega + 3, -3 \, \omega - 3, 2 \, \omega + 3, -2 \, \omega - 3, \omega + 3, -\omega - 3, 4 \, \omega + 4, -4 \, \omega - 4, 3 \, \omega + 4, -3 \, \omega - 4, 2 \, \omega + 4, -2 \, \omega - 4, 5 \, \omega + 5, -5 \, \omega - 5, 4 \, \omega + 5, -4 \, \omega - 5, 3 \, \omega + 5, -3 \, \omega - 5, 6 \, \omega + 6, -6 \, \omega - 6, 5 \, \omega + 6, -5 \, \omega - 6, 4 \, \omega + 6, -4 \, \omega - 6, 7 \, \omega + 7, -7 \, \omega - 7, 6 \, \omega + 7, -6 \, \omega - 7, 5 \, \omega + 7, -5 \, \omega - 7, -2 \, \omega + 5, 2 \, \omega - 5, 4 \, \omega + 7, -4 \, \omega - 7\right\}  \end{dmath*}\end{minipage} }\\
 & \\
Phase 1 (method  9): &
\checkmark, $\#\mathcal{Q} = $ 46 $ $ \\ 
$b,b,\dots,b$ inputs (method  21): & \checkmark, maximal length of window: $ 5 $ \\
\multicolumn{2}{l}{\begin{minipage}{\textwidth} Computation of Phase 2 (method  21) was killed. % when the length of window 1 was being proccessed. Numbers of saved combinations for each finished length are: ()
\end{minipage} }\\
\end{tabular}

\end{exmp}





\begin{exmp}
\label{ex:killAC}


\rule{0cm}{0cm}

\begin{tabular}{ll}
$\omega=  {\left(\frac{1}{9} \, \sqrt{19} \sqrt{3} + 1\right)}^{\frac{1}{3}} + \frac{2}{3 \, {\left(\frac{1}{9} \, \sqrt{19} \sqrt{3} + 1\right)}^{\frac{1}{3}}} $  & $\beta= -2 \, \omega^{2} + \omega + 2 = {\left(\sqrt{57} - \frac{197}{27}\right)}^{\frac{1}{3}} - \frac{14}{9 \, {\left(\sqrt{57} - \frac{197}{27}\right)}^{\frac{1}{3}}} - \frac{2}{3} $\\
$m_\omega(t)=  t^{3} - 2 \, t - 2 $  & $m_\beta(x)=  x^{3} + 2 \, x^{2} + 6 \, x + 18 $\\
Real conjugate of $\beta$ greater than 1:   &  no \\
$\#\A= $ 31 $ $ & $\A$ is not minimal. \\
\multicolumn{2}{l}{\begin{minipage}{\textwidth}\begin{dmath*}\A = \left\{0, 1, -1, \omega^{2} + \omega + 1, -\omega^{2} - \omega - 1, \omega + 1, -\omega - 1, -\omega^{2} + \omega + 1, \omega^{2} - \omega - 1, \omega^{2} + 1, -\omega^{2} - 1, -\omega^{2} + 1, \omega^{2} - 1, \omega^{2} - \omega + 1, -\omega^{2} + \omega - 1, -\omega + 1, \omega - 1, -\omega^{2} - \omega + 1, \omega^{2} + \omega - 1, \omega^{2} + \omega, -\omega^{2} - \omega, \omega, -\omega, \omega^{2} + 2, -\omega^{2} - 2, 2, -2, -\omega^{2} + \omega, \omega^{2} - \omega, \omega^{2}, -\omega^{2}\right\}  \end{dmath*}\end{minipage} }\\
 & \\
Phase 1 (method  9): &
\checkmark, $\#\mathcal{Q} = $ 83 $ $ \\ 
$b,b,\dots,b$ inputs (method  21): & \checkmark, maximal length of window: $ 5 $ \\
\multicolumn{2}{l}{\begin{minipage}{\textwidth} Computation of Phase 2 (method  21) was killed when the length of window 4 was being proccessed. Numbers of saved combinations for each finished length are: (0, 71, 1887261)\end{minipage} }\\
\end{tabular}

\end{exmp}


\begin{exmp}
\label{ex:killAD}


\rule{0cm}{0cm}

\begin{tabular}{ll}
$\omega=  {\left(\frac{1}{9} \, \sqrt{29} \sqrt{3} + \frac{28}{27}\right)}^{\frac{1}{3}} + \frac{1}{9 \, {\left(\frac{1}{9} \, \sqrt{29} \sqrt{3} + \frac{28}{27}\right)}^{\frac{1}{3}}} + \frac{1}{3} $ \\
\multicolumn{2}{l}{ $\beta= -\omega^{2} + \omega - 1 = {\left(\frac{2}{9} \, \sqrt{29} \sqrt{3} - 2\right)}^{\frac{1}{3}} - \frac{2}{3 \, {\left(\frac{2}{9} \, \sqrt{29} \sqrt{3} - 2\right)}^{\frac{1}{3}}} - 1 $}\\
$m_\omega(t)=  t^{3} - t^{2} - 2 $  & $m_\beta(x)=  x^{3} + 3 \, x^{2} + 5 \, x + 7 $\\
Real conjugate of $\beta$ greater than 1:   &  no \\
$\#\A= $ 16 $ $ & $\A$ is minimal. \\
\multicolumn{2}{l}{\begin{minipage}{\textwidth}\begin{dmath*}\A = \left\{0, 1, -1, \omega^{2} + \omega + 1, -\omega^{2} - \omega - 1, \omega + 1, \omega^{2} + 1, -\omega^{2} - 1, -\omega^{2} + 1, \omega^{2} + \omega, \omega, -\omega, -\omega^{2} + \omega, \omega^{2} - \omega, \omega^{2}, -\omega^{2}\right\}  \end{dmath*}\end{minipage} }\\
 & \\
Phase 1 (method  9): &
\checkmark, $\#\mathcal{Q} = $ 99 $ $ \\ 
$b,b,\dots,b$ inputs (method  21): & \checkmark, maximal length of window: $ 6 $ \\
\multicolumn{2}{l}{\begin{minipage}{\textwidth} Computation of Phase 2 (method  21) was killed when the length of window 4 was being proccessed. Numbers of saved combinations for each finished length are: (73, 5329, 315494)\end{minipage} }\\
\end{tabular}

\end{exmp}



%%%%%%%%%%%%%%%%%%%%%%%%%%%%%%%%%%%%%%%%%%%%%%%%%%%%%%%%%%%%%%%%%%%%%%%%
%     K O N E C   P R Á C E                                            %
%%%%%%%%%%%%%%%%%%%%%%%%%%%%%%%%%%%%%%%%%%%%%%%%%%%%%%%%%%%%%%%%%%%%%%%%


\end{document}